\documentclass[11pt]{article}% Shared preamble for the rigor documents (Lamport-style structured proofs).  \input this file after \documentclass.
\usepackage[T1]{fontenc}\usepackage{lmodern}\usepackage{microtype}
\usepackage[margin=1in]{geometry}
\usepackage{amsmath,amssymb,amsthm,mathtools}
\usepackage{xcolor}\usepackage{enumitem}\usepackage{booktabs}\usepackage{graphicx}\usepackage{hyperref}
\usepackage[most]{tcolorbox}
\definecolor{navy}{HTML}{17365D}\definecolor{teal}{HTML}{117A8B}\definecolor{warn}{HTML}{A33A2B}\definecolor{okgreen}{HTML}{2E7D4F}
\hypersetup{colorlinks=true,linkcolor=navy,citecolor=teal,urlcolor=teal}
\theoremstyle{plain}
\newtheorem{theorem}{Theorem}[section]\newtheorem{lemma}[theorem]{Lemma}\newtheorem{proposition}[theorem]{Proposition}\newtheorem{corollary}[theorem]{Corollary}
\theoremstyle{definition}\newtheorem{definition}[theorem]{Definition}\newtheorem{remark}[theorem]{Remark}\newtheorem{claim}[theorem]{Claim}\newtheorem{issue}[theorem]{Issue}
% ---------- Lamport structured proofs ----------
% \lstep{level}{number}{assertion}   -> indented "<level>number. assertion"
% \lproof                             -> "PROOF:" line (start of the sub-proof of the preceding step)
% \lby{justification}                 -> "BY ..." leaf justification for the preceding step
% \lqed{level}{number}                -> QED step of a sub-proof
% keywords: \lassume \lprove \lsuffices \lcase \ldefine \lpick \lnew
\newlength{\lindent}\setlength{\lindent}{1.7em}
\newcommand{\lstep}[3]{\par\smallskip\noindent\hspace*{\dimexpr\numexpr#1-1\relax\lindent\relax}{\color{navy}$\langle#1\rangle#2$.}\ #3}
\newcommand{\lproof}{\par\noindent\hspace*{\lindent}{\scshape Proof:}}
\newcommand{\lby}[1]{\par\noindent\hspace*{\lindent}{\scshape By}\ #1.}
\newcommand{\lqed}[2]{\par\smallskip\noindent\hspace*{\dimexpr\numexpr#1-1\relax\lindent\relax}{\color{navy}$\langle#1\rangle#2$.}\ {\scshape Q.E.D.}}
\newcommand{\lassume}{{\scshape Assume}\ }\newcommand{\lprove}{{\scshape Prove}\ }\newcommand{\lsuffices}{{\scshape Suffices}\ }
\newcommand{\lcase}{{\scshape Case}\ }\newcommand{\ldefine}{{\scshape Define}\ }\newcommand{\lpick}{{\scshape Pick}\ }\newcommand{\lnew}{{\scshape New}\ }
% ---------- verdict boxes ----------
\newtcolorbox{verdictok}{colback=okgreen!8,colframe=okgreen,title={\bfseries Verdict: verified},fonttitle=\small}
\newtcolorbox{verdictgap}{colback=orange!10,colframe=orange!80!black,title={\bfseries Verdict: gap / caveat},fonttitle=\small}
\newtcolorbox{verdictbreak}{colback=warn!10,colframe=warn,title={\bfseries Verdict: proof-breaking issue},fonttitle=\small}
\newtcolorbox{citedbox}[1]{colback=navy!5,colframe=navy!60,title={\bfseries Cited result: #1},fonttitle=\small,breakable}
\setlength{\parindent}{0pt}\setlength{\parskip}{4pt}\allowdisplaybreaks
\newcommand{\cA}{\mathcal A}\newcommand{\cF}{\mathcal F}\newcommand{\cM}{\mathfrak M}\newcommand{\cN}{\mathfrak N}

% extra calligraphic shortcuts (\providecommand: no clash if a document defines its own)
\providecommand{\cB}{\mathcal B}\providecommand{\cC}{\mathcal C}\providecommand{\cD}{\mathcal D}\providecommand{\cE}{\mathcal E}\providecommand{\cG}{\mathcal G}\providecommand{\cH}{\mathcal H}\providecommand{\cI}{\mathcal I}\providecommand{\cJ}{\mathcal J}\providecommand{\cK}{\mathcal K}\providecommand{\cL}{\mathcal L}\providecommand{\cO}{\mathcal O}\providecommand{\cP}{\mathcal P}\providecommand{\cQ}{\mathcal Q}\providecommand{\cR}{\mathcal R}\providecommand{\cS}{\mathcal S}\providecommand{\cT}{\mathcal T}\providecommand{\cU}{\mathcal U}\providecommand{\cV}{\mathcal V}\providecommand{\cW}{\mathcal W}\providecommand{\cX}{\mathcal X}\providecommand{\cY}{\mathcal Y}\providecommand{\cZ}{\mathcal Z}
\newcommand{\Fid}{\operatorname{F}}\newcommand{\Tr}{\operatorname{Tr}}\newcommand{\eps}{\varepsilon}\newcommand{\dd}{\,\mathrm d}

\title{Consolidated findings of the verification program\\[0.3em]\large Notes N1--N7 of the CMI / Petz-recovery series}\author{}\date{Assembled from the agent reports (rigor/findings\_entries.tex)}
\begin{document}\maketitle\tableofcontents\clearpage
\section{Scope and method}
Fourteen Opus agents re-derived every stated result of the seven notes in Lamport's structured-proof style (documents listed in \texttt{rigor/INDEX.md}), downloaded the cited references (\texttt{refs/}, 38 PDFs; a few paywalled items were replaced by open restatements, marked NOT OBTAINED), transcribed the exact cited statements with screenshots (\texttt{cited\_R1--R4.tex}), and attacked the four remaining rigor items of Note~7 (\texttt{lemma\_second\_variation}, \texttt{quadratic\_limit}, \texttt{kernel\_identity}, \texttt{collar\_and\_invariant\_formula}). Severity: BREAK = invalidates a statement as written; GAP = missing argument, likely repairable; MINOR = presentation or attribution. Each entry names the note (N1--N7), the location, the reporting agent, and the finding.

\section{Proof-breaking issues in the notes (as stated)}
\begin{enumerate}[leftmargin=1.6em,itemsep=6pt]
\item \textbf{[BREAK]} \emph{N5 \S5 (Translation paragraph, eq.\ collapse-vwz); N6 \S1, \S5; N7 Cor.~6.1} \hfill\textsc{R2}\\ For a two-dimensional CFT the Bisognano--Wichmann boost dilates the two light-ray coordinates oppositely, so the rotated Petz map with parameter \(t\) acts as \(t\) on one chirality and \(-t\) on the other: \(-\log F^{(\lambda)}=\Phi_c\bigl(z(1+e^{-\pi\lambda})\bigr)+\Phi_{\bar c}\bigl(z(1+e^{\pi\lambda})\bigr)\), even in \(\lambda\), minimal at \(\lambda=0\) (small \(\eta\): \(f_2\eta^2[2+4\cosh\pi\lambda+2\cosh2\pi\lambda]\)). Consistent with VWZ's ``largest fidelity at \(\lambda=0\)'' and with \(F^{(\lambda)}=F^{(-\lambda)}\) for real lattice states. The chiral collapse onto one function and all \(\lambda=0\) statements are unaffected; the claims ``the opposite sign of \(\lambda\) improves the recovery'' and ``the best rotated map is the geometric compression'' are chiral-only and must be withdrawn for 2D CFTs. Erratum required in N5, N6, N7.
\item \textbf{[BREAK]} \emph{N7 main box and Section~1 ``Theorem A''} \hfill\textsc{V6}\\ ``\(\lim_{\zeta\to0}\Phi/\zeta^2=r/(12\pi^2)\)'' is stronger than what Section~5 proves (Thm~5.2 gives the one-particle \(g_B\), \(g_{\mathrm{KM}}\) only); the identification with \(\lim\Phi/\zeta^2\) rests on N5 Cor.~4.2 (heuristic) and on numerics. Theorem~B is labelled conditional, Theorem~A is not. Erratum: state Theorem~A as the exact evaluation of the one-particle metrics, and the limit statement as conditional on the interchange lemma (agent P2's task).
\item \textbf{[BREAK]} \emph{N7 \S3.3, \S4.2 (application of Lemmas~3.1--3.2 to the zero-collar representative); N6 Conj.~2.1} \hfill\textsc{V6}\\ For the \(\xi\)-frame representative \(\|T(g)\Omega\|^2=\frac1{2\pi}\int\widehat W|\widehat{\widetilde g}|^2dk=\infty\) (\(\widehat W(0)=c/48\pi^2\), \(|\widehat{\widetilde g}|^2\sim4/k^2\)), so the hypothesis \(\Omega\in D(A)\) and Lemma~3.1(2)(3) fail in the intended application; only the weighted integrals converge. The tangent functional itself is unambiguous (smooth continuations of \(g\) beyond \(L\) differ by elements of \(\mathcal F(I)'\)), so Theorem~B needs a form-domain lemma equating the invariant formula for a finite-norm representative with the \(\xi\)-frame spectral integrals (this is what agent P4's \(x\)-space evaluation tests). Theorem~A (free fermion) is unaffected: the one-particle sums converge absolutely and the kernel identity is proved.
\item \textbf{[BREAK]} \emph{N5 \S7.2} \hfill\textsc{V6}\\ ``\(96f_2z\approx0.81z\)'' is a factor 2 wrong: \((-2\log F)/I=48f_2z\approx0.405z=0.20\zeta\) (consistent with Table~6.3); ``\(-\log F\simeq f_2(12/r)^2I^2\)'' has the wrong \(r\)-scaling for the per-component \(f_2\): correct is \(144f_2I^2/r=12I^2/(\pi^2r)\). N7's \(12/(\pi^2c)\) is right.
\item \textbf{[BREAK]} \emph{N5 Cor.~5.1 (last sentence)} \hfill\textsc{V6}\\ ``The modular average is worse \dots because \(\int p(t)\Phi(\zeta_t)dt\) receives contributions from arbitrarily large \(\zeta_t\)'' is invalid: by concavity that integral is an UPPER bound on the averaged channel's error (\(-\log F(\omega,\int p\,\widetilde\omega_t)\le\int p\,\Phi(\zeta_t)\)). Withdraw or reformulate (``the FHSW bound for the averaged channel is weaker than for the best fixed \(t\)'').
\item \textbf{[BREAK]} \emph{N5 Cor.~5.1 vs its VWZ paragraph} \hfill\textsc{V6}\\ The letter \(\eta\) denotes \(\eta_{\mathrm{LX}}=b(a+b+c)/((a+b)(b+c))\) (``\(z=1/\eta-1\)'') and, ten lines later, VWZ's \(\eta=x\) (``\(z=\eta/(1-\eta)\)''); \(\eta_{\mathrm{LX}}+\eta_{\mathrm{VWZ}}=1\). Rename VWZ's cross ratio (e.g.\ \(x\) or \(\eta_{\mathrm{VWZ}}\)) throughout N5--N7.
\item \textbf{[BREAK]} \emph{N6 \S2.1 Step~4} \hfill\textsc{V6}\\ ``With \(\langle TT\rangle=\frac c2(x-y)^{-4}\)'' contradicts the generator convention \(U=e^{isT(f)}\) by \((2\pi)^2\); N7 eq.~(1) (\(c/8\pi^2\)) is correct. Fix the strategy note's sketch.
\item \textbf{[BREAK]} \emph{N7 \S3.3 (proof route of Lemma~3.2) and N6 \S2 step~2} \hfill\textsc{R3}\\ ``\(\Fid(\omega,\omega_s)=\langle\Omega,\Delta_{\omega_s,\omega}^{1/2}\Omega\rangle\)'' is false: that quantity is the quasi-entropy \(Q_{1/2}(\omega_s\Vert\omega)=\Tr(\rho_{\omega_s}^{1/2}\rho_\omega^{1/2})\) (Hiai 2018 Def.~5.1 / Prop.~A.4(1)), equal to the fidelity only for commuting states (2\(\times\)2 counterexample 0.93301 vs \(F=0.93541\)); its second variation is the Wigner--Yanase metric (23.476382 vs \(g_B=21.172666\) in the test). The formula is nowhere in Uhlmann 1976. Correct standard-form statement: \(\Fid(\omega_\psi,\omega)=\sup_{u'\in U(\cM')}|\langle\Omega,u'\psi\rangle|\) (Alberti--Uhlmann 2002 Thm.~2(3)). Scope: only the announced proof route of the (still open) type-III lemma; Theorem~A and eqs.\ (11)--(12) unaffected.
\item \textbf{[BREAK]} \emph{N5 \S3.1 (sentence ``keeps all eigenvalues of \(Q_N\) and \(\widetilde Q_N\) at distance \(\gtrsim10^{-4}\) from 0 and 1 \dots double-precision evaluation stable''; compression\_box.py docstring)} \hfill\textsc{V4}\\ False: the smallest eigenvalue is \(2.8\times10^{-13}\) at \(\Lambda=60\) and \(-3.6\times10^{-16}\) at \(\Lambda=120\); double precision fails (results\_batch3.txt has \(-\log F=-8.5\)). The sentence is a leftover from before the \(t\to0\) endpoint fix; the method actually used (spectral sub-compression with 40-digit arithmetic) is the one described in \S6. Erratum required.
\item \textbf{[BREAK]} \emph{N5 \S3.1 (``monotone in \(\kappa_{\max}\)'')} \hfill\textsc{V4}\\ The sub-compressed subspaces are not nested across \(\kappa_{\max}\) (Table~1: \(3.51944\to3.51505\to3.51911\)); each raw value is a rigorous lower bound (once the basis issue below is fixed) but the sequence is not monotone. Reword.
\item \textbf{[BREAK]} \emph{N5 Cor.~5.1 eq.\ (collapse-vwz) and the following paragraph; N5 abstract, main box, \S7.2, \S8 bullet~3; \S7.1 ``collapse holds for every conformal net''} \hfill\textsc{V3}\\ Confirms [R2]: for a 2D CFT the two chiralities have \(K^{(-)}=-K^{(+)}\), hence opposite rotations \(s_t^{\pm}=\lambda(1+e^{\mp2\pi t})\) and \(-\log F^{(\lambda)}=\Phi(\zeta_-)+\Phi(\zeta_+)\), even in \(\lambda\) with minimum at \(\lambda=0\) (proved via the one-particle modular Hamiltonian and via reality of lattice states, and consistent with VWZ p.~12 item~3). ``Ordinary Petz map not optimal / best map is the geometric compression / factor 4'' are chiral-only statements and false for the Dirac fermion; ``for every conformal net'' must read ``for every chiral net''.
\item \textbf{[BREAK]} \emph{N6 \S2 (P1 step~3) and \S3 (Route A): energy bound} \hfill\textsc{R4}\\ ``\(\|T(f)\psi\|\le C\|f\|_{H^{3/2}}\|(1+L_0)\psi\|\)'' is not a theorem; the proved bound (Carpi--Weiner, = CDIT Prop.~2.2(4)) uses \(\|f\|_{3/2}=\sum_n|\hat f_n|(1+|n|^{3/2})\), and \(H^s\hookrightarrow\) that space only for \(s>2\). Rescue: CDIT Prop.~2.3 (piecewise smooth and \(C^1\) \(\Rightarrow\) finite \(\|\cdot\|_{3/2}\)), which covers the \(C^{1,1}\) field of the notes. Attribution must be Carpi--Weiner, not Fewster--Hollands (FH05 contains only the mode bound (5.2)). Also: the authors are Carpi--Del Vecchio--Iovieno--Tanimoto (not Weiner), and their Outlook gives \(D^s\), \(s>2\), for Fock-space representations with integer \(c\) [DIT19], so Route~A is already closed for the free Dirac fermion.
\item \textbf{[BREAK]} \emph{rigor/quadratic\_limit.tex Prop.~7.1, step \(\langle1\rangle5\)} \hfill\textsc{V7}\\ The claim \(2\iint_{|v|>\Lambda}|D^{(1)}|^2/N\le C\Lambda^2e^{-\Lambda/2}\) is false: the \(u\)-integration is unrestricted, so the left side is \(\Theta(\Lambda^{-2})\) (\(2.6\cdot10^{-4}\) at \(\Lambda=40\), \(1.6\cdot10^{-5}\) at \(160\)). Repair: intersect with \(\{|u|\le\Lambda\}\); the true second term is \(\Theta(\Lambda^{-4})\), constant \(\approx13.5\). Statement survives; Thms.~4.2 and 6.4 unaffected.
\item \textbf{[BREAK]} \emph{N7 \S3.1, sentence after eq.\ (13)} \hfill\textsc{V5}\\ ``both are bounded by \(1+\lambda\), which is why \(g_B\le4\|A\Omega\|^2\) \dots\ and why the second variations are finite whenever \(A\Omega\) is a vector'' is false for the Kubo--Mori weight: \((\lambda-1)\log\lambda>1+\lambda\) already at \(\lambda=e^2\) (12.778 vs 8.389) and the ratio diverges, so \(g_{\mathrm{KM}}<\infty\) needs \(A\Omega\in\mathcal Q(\Delta\log\Delta)\), exactly the extra hypothesis Lemma~3.2 imposes. Local: no computation uses it. Also ``equality for pure states'' is loose (a pure state is not faithful).
\item \textbf{[BREAK]} \emph{N6 \S2 (P1) step~3} \hfill\textsc{V5}\\ ``\(\langle\widetilde T(\xi)\widetilde T(\xi')\rangle=\frac c{32}\sinh^{-4}\frac{\xi-\xi'-i0}2\)'' is the BPZ constant, \(4\pi^2\) times the generator-normalized \(\frac c{128\pi^2}\) that the amended step~4 invokes (N7 eq.\ (1)); the two halves of one argument now use two normalizations, and a reader following step~3 into step~4 obtains \(f_2=c/3\). Fix: \(\frac c{32}\to\frac c{128\pi^2}\), and give the density with its constant, \(\widehat W=\frac c{24\pi}\frac{k(k^2+1)}{e^{2\pi k}-1}\). The step-4 sentences ``fixed by matching the free-fermion number'' / ``the first task below'' are stale.
\item \textbf{[BREAK]} \emph{phase1\_brief.md step 1: ``on \([L,\infty)\) any smooth increasing map \dots\ tending to the identity at infinity''} \hfill\textsc{PN}\\ That restriction is fatal to sharpness. Requiring \(u_s\to0\) at \(|x|=\infty\) forces the circle field to vanish to second order at the point at infinity and pins \(\inf\tfrac12E/\zeta^2\ge0.0678\), eight times \(g_{\rm B}/8\). The optimal extension is asymptotically affine and moves \(\infty\); it is a diffeomorphism of \(S^1\), not of \(\mathbb R\).
\item \textbf{[BREAK]} \emph{rigor/tail\_bound.tex Prop.~8.1 vs quadratic\_limit Problem~7.2} \hfill\textsc{PT}\\ The planned rate argument cannot work as stated: the tail bound needs \(\kappa_c>2\log(1/\zeta)\) (vacuum tail \(e^{-\kappa_c}\ll\zeta^2\)) while Problem~7.2's cubic constant \(e^{\kappa_c}\) needs \(\kappa_c<\log(1/\zeta)\). Incompatible. With \(Ce^{c\kappa_c}\zeta^3\), \(c<1/2\), one gets only \(\Phi=f_2\zeta^2[1+O(\log^{-1}(1/\zeta))]\), not \(\log^{-2}\).
\item \textbf{[BREAK]} \emph{rigor/implementation\_C11.tex Prop.~4.10 (form bound; corrects the first version)} \hfill\textsc{QB}\\ The claimed open window \(1/2\le\beta<1\) is empty. Optimising \(\psi_t=\Omega+t\lVert L_{-n}\Omega\rVert^{-1}L_{-n}\Omega\) gives \(\sup_t=\lVert L_{-n}\Omega\rVert/\sqrt{1+n}\sim(c/12)^{1/2}n\), forcing \(\beta\ge1\); \(\beta=1\) holds (CW05 Lemma~4.1). Since \(\lVert g_{\rm tot}'\rVert_1=\infty\) logarithmically, the Nelson/commutator route to \(D(L_0)\)-invariance is closed, not open.
\item \textbf{[BREAK]} \emph{open\_problems\_plan.tex \S O2 (proposed route ``Kosaki for the value'')} \hfill\textsc{QK}\\ Kosaki's variational expression (Hiai arXiv:1805.02050 Thm.~3.5) and Petz--Donald's are both \emph{suprema}: any trial function yields a lower bound on the relative entropy. Neither can supply the missing upper bound \(\limsup_s s^{-2}D\le\frac12 g_{\mathrm{KM}}\). The route as planned cannot work.
\item \textbf{[BREAK]} \emph{rigor/implementation\_C11.tex Rem.~5.2 and last summary row} \hfill\textsc{QR2}\\ The ``open window \(\tfrac12\le\beta<1\)'' for the form bound \(\lvert\langle\psi,T(f)\psi\rangle\rvert\le C\lVert f\rVert_\beta\lVert(1+L_0)^{1/2}\psi\rVert^2\) is empty. Optimising QB's own test family \(\psi_t=\Omega+tL_{-n}\Omega/\lVert L_{-n}\Omega\rVert\) over \(t\) gives quotient \(\sim(c/12)^{1/2}n\), forcing \(\beta\ge1\); \(\beta=1\) holds (three lines from Carpi--Weiner Lemma~4.1) and \(\lVert g_{\rm tot}'\rVert_1=\infty\) logarithmically. The Nelson route to \(D(L_0)\)-invariance is closed, not open.
\item \textbf{[BREAK]} \emph{Note~5 Cor.~4.2 (A2)--(A3), remainder \(\zeta^2/(120\log^2(1/\zeta))\)} \hfill\textsc{RR}\\ Refuted. With a one-parameter-group extension the Uhlmann bound gives \(\Phi\le-\tfrac12\log(1-2f_2\zeta^2)=f_2\zeta^2(1+f_2\zeta^2+O(\zeta^4))\), so the remainder's positive part is \(O(\zeta^4)\). The alleged term would add \(+4.3\%,+5.9\%,+8.5\%\) at \(\zeta=1/120,1/60,1/30\); the measured deviations are \(-0.80\%,-1.60\%,-3.18\%\).
\item \textbf{[BREAK]} \emph{rigor/quadratic\_limit.tex Problem 7.2, guess \(C(\Lambda)\sim e^{\Lambda}\)} \hfill\textsc{RR}\\ Too pessimistic. The exact cubic coefficient of the compressed curve is \(C(\kappa_c)=0.008113,\dots,0.008355\) at \(\kappa_c=9.2,\dots,23.0\): bounded, limit \(0.00838(3)=|c_3|f_2\), box-independent. The optimised lower rate is then linear, \(\Phi\ge f_2\zeta^2(1-0.99\zeta)\), not \(\log^{-2}\); bounding the two pieces separately gives \(O(\zeta^{2/3})\).
\item \textbf{[BREAK]} \emph{phase2b\_brief.md, track RK, step (c) (the Kubo--Mori gauge lemma)} \hfill\textsc{RK}\\ False. For every interval $K\supseteq I$ carrying the implementer, $g_{\mathrm{KM}}^{(K)}(T(\hat g)\Omega)=\frac{c}{12}\int(\tilde g''^2+\tilde g'^2)d\xi$ in the modular frame of $K$, and the infimum over $K$ and over extensions is $\frac{11+5\sqrt5}{12}c=1.848c$ at $R=\frac{3+\sqrt5}{2}$, i.e.\ $11.09$ times $c/6$.
\item \textbf{[BREAK]} \emph{kubo\_mori\_second\_variation.tex, hypotheses $(\mathrm U_\varepsilon)$/$(\mathrm A_\varepsilon)$, applied to $\mathcal A(K)$} \hfill\textsc{RK}\\ Both fail for every admissible gauge and every $\varepsilon>0$: locality in $K$ makes the field compactly supported in the modular frame, so its Fourier transform decays polynomially, while $\lambda^{-\varepsilon}$ weights the spectral measure by $e^{2\pi\varepsilon\lvert k\rvert}$ as $k\to-\infty$. Locality and modular analyticity are incompatible.
\item \textbf{[BREAK]} \emph{certified\_numerics.tex, Table 3 and the $\Phi_{\mathrm{off}}$ column of \S7 (script cert\_tail\_np.py)} \hfill\textsc{RCr}\\ $\Phi_{\mathrm{off}}=-\log\lvert\det(1+T)\rvert$ is evaluated in float64 on matrices of condition $10^{13}$. Under $10^{-16}$ input jitter it gives $2.60,2.94,1.49,2.38\times10^{-9}$ at $\zeta=1/120$, and $-1.16\times10^{-10}$ from the trapezoid input; the script's own identity $-\log\lvert\det M_A\rvert=\Phi_W+\Phi_{22}$ fails by $1.0\times10^{-4}$ ($180\times$). Widths $0.64$--$1.19\,\%$ unestablished.
\item \textbf{[BREAK]} \emph{rate\_of\_remainder.tex, Prop.\ 4.2 and its verdictbreak against quadratic\_limit.tex Problem 7.2} \hfill\textsc{RRr}\\ Problem 7.2 is two-sided, so its constant must dominate the quartic Taylor coefficient of $\Phi_P$, which contains $\frac14\sum_{\lvert\kappa\rvert<\kappa_c}\lvert(\mathcal D_2)_{\kappa\kappa}\rvert^2/N_{\kappa\kappa}\approx e^{\kappa_c}/(225\kappa_c^6)$: $0.023,0.087,0.42,477$ at $\kappa_c=18.4,20.7,23.0,32.2$. Exponential, as guessed. The one-sided form is what the proof uses and is unaffected.
\item \textbf{[BREAK]} \emph{cmi\_cft\_research\_program.tex ``is the Petz map optimal?''; open\_problems\_plan.tex \S O8; rigor/optimality\_second\_order.tex \S6 Q1} \hfill\textsc{S8}\\ The Petz map is NOT second-order optimal. The zero-collar geometric compression is an admissible recovery channel with \(\Phi(z)=f_2z^2\) against the Petz \(\Phi(2z)=4f_2z^2\); on the hopping chain the numerical optimum over quasi-free channels reaches \(0.270\)--\(0.280\) of the exact Petz error.
\item \textbf{[BREAK]} \emph{open\_problems\_plan.tex \S O8 ``Route''; rigor/phase3\_brief.md track S8 task (3); rigor/optimality\_second\_order.tex Prop.~5.2} \hfill\textsc{S8}\\ The proposed lower-bound route via the \(A\)--\(C\) restriction yields exactly zero. The vacuum \(A\)--\(C\) state is itself \(B\)-mediated: \(Q_{AC}=Q_{AB}T\) with \(T=Q_{AB}^{+}Q_{AC}\), residual \(10^{-17}\), \(\lVert T\rVert=0.55,0.38,0.29\) and the complete-positivity interval satisfied with margin \(0.12\)--\(0.16\). [corrected 2026-09-24: not a general zero. \(E_{AC}=0\) is proved as an implication from (Tcond), in Prop.~4.2 (\texttt{prop:AC-zero}; cited above as Prop.~5.2), and is applied only on the three hopping chains \((L_A,L_B,L_C)=(4,8,2)\), \((4,12,2)\), \((4,16,2)\), where (Tcond) holds numerically; nothing is claimed for other geometries or the continuum; card \texttt{ac-lower-bound-zero}]
\item \textbf{[BREAK]} \emph{rigor/sdp\_dual\_certificate.tex Thm.~9.1, \S12 Verdict; rigor/phase3b\_brief.md S10 target; rigor/optimality\_second\_order.tex Verdict~Q2/Q3; numerics/optimality\_all/s10\_\{circle2,theta,diag,nonpert,implement\}.py} \hfill\textsc{S10}\\ The zero-collar geometric compression is \emph{not} the second-order optimum of the quasi-free convex programme. Composing it with a Bogoliubov rotation $e^{\varepsilon\mathcal G}$ inside $B$ ($\varepsilon=O(s)$, $\lVert\varepsilon\mathcal G\rVert=O(s)$) keeps both LMI constraints active, hence stays feasible, and lowers the objective by the factor $1-\theta$ with $\theta=\sup_{\mathcal G}\cos^2\angle(u,\Pi P^{\perp}\mathcal G P)=0.300\pm0.005$: $\min_{\mathcal F}g_Q\le(0.70\pm0.01)f_2z^2$. Verified non-perturbatively (exact defect of the explicit channel, $\min_\varepsilon g_Q/g_Q(\text{compression})=0.7050,0.7046,0.7038$ at $s=10^{-3},2\times10^{-3},5\times10^{-3}$, against the predicted $0.7054$), and the same model reproduces Theorem~A ($g_Q(\delta_c)=f_2z^2$) to $1\%$. Consequence: since the Sion dual value satisfies $\mathcal D_z\le\min_{\mathcal F}g_Q$ (optimality\_all\_channels.tex Thm.~4.1), \emph{no} test operator can certify $\liminf z^{-2}E_{\mathrm{rec}}\ge f_2$; route~2 of the Phase-3b brief is closed and the Q3 coefficient obeys $\kappa_{\mathrm{opt}}\le0.71$, not $1.00\pm0.15$. The one numerical input is $E_D u E_D\neq0$ ($\lVert E_DuE_D\rVert_2^2/\lVert u\rVert_2^2=0.0143$, stable in system size, weight cap, extension and UV taper).
\item \textbf{[BREAK]} \emph{rigor/optimality\_all\_channels.tex \S5 and Verdict; rigor/sdp\_dual\_certificate.tex; numerics/optimality\_all/theta\_galerkin.py, exact\_fidelity\_rotated.py, uhlmann\_rotated.py, s11\_exactF\_L*.out} \hfill\textsc{ORCH-3b}\\ S10's second-order refutation is confirmed independently: in the block-exact modular Galerkin frame (no cap, no taper, exact \(P\) on \(I\)) the isometric-orbit gain is \(\theta=0.286,0.299,0.314,0.324,0.336\) for \(h=0.4,0.3,0.25,0.2,0.15\), increasing with resolution and independent of \(s\); with S10's spectral value \(0.300\pm0.005\) this gives \(\min_{\mathcal F}g_Q\le0.71f_2z^2\). Consequence: the zero-collar compression is NOT the minimiser of the second-order functional over quasi-free (equivalently, over all) channels; S8 Verdict Q2 superseded at that level. HOWEVER the exact-fidelity statement is NOT established: the second-order functional is a lower bound for every channel, so \(E_{\mathrm{rec}}\in[\mathcal D_z(1-o(1)),f_2z^2(1+o(1))]\) with \(\mathcal D_z\le0.71f_2z^2\), and the upper half \(-\log F\le g_Q(1+o(1))\) for the rotated channel is open. S11's windowed exact fidelities (ratio \(0.59\to0.67\)) are ratios of lower bounds; the orchestrator's full exact fidelity in the spectral model (60 digits, exact \(P\)) gives ratio \(0.94\) at \(L=96\), \(s=0.005\), with the compression's own exact error \(2.95\times\) its second-order value and the Uhlmann bound \(4.3\times\) (stable in \(L=96\)--\(384\)): the finite spectral model does not resolve the exact fidelity. The decisive open statement: does the invisible component of \(P^\perp\mathcal G_*P\) lie in the closure of the exterior variations (then the Uhlmann route gives the upper half)?
\end{enumerate}
\section{Errors found in the cited sources}
\begin{enumerate}[leftmargin=1.6em,itemsep=6pt]
\item \textbf{[BREAK-IN-SOURCE]} \emph{FHSW arXiv:2006.08002 \S4.2 eq.\ (167) and hsm definition} \hfill\textsc{V2 (and R2)}\\ With FHSW's own (25) and (166), \(\alpha^t=\operatorname{Ad}U(a(1+e^{+2\pi t}))\); as printed, (167) is \(\operatorname{Ad}U(-a(1+e^{-2\pi t}))\), which maps \(\mathcal A\) out of \(\mathcal B\), violating their (22); ``for \(t\ge0\)'' must read \(t\le0\). Documented with screenshot (rigor/check\_note4.tex Thm~3.8). Our N4 does not rely on the printed formula.
\end{enumerate}
\section{Disputed items (superseded by a resolution below)}
\begin{enumerate}[leftmargin=1.6em,itemsep=6pt]
\item \textbf{[DISPUTED]} \emph{N4 \S1.1 vs \S9 (coefficient \(r/6\) in eqs.\ (LX-MI), (CMI-LX), (norm-CMI-relation), (small-z), fidelity bounds, main box); also N1, N2} \hfill\textsc{V2 vs R1}\\ V2 claims LX's \(r/6\) is the value for \(r\) Majorana fermions (\(c=r/2\)) so that our gauge-invariant CAR over \(L^2(\mathbb R;\mathbb C^r)\) should carry \(r/3\); R1 claims LX's \(\mathcal A_r\) is \(r\) complex fermions (\(c=r\)) and \(r/6\) is correct. Resolution by the orchestrator from the LX text (Section~3.6: ``finite dimensional subspaces \(H_i\) of \(L^2(I_i,\mathbb C^r)\)'', \(\mathrm{CAR}(H_i)\); \S4.4: \(U(1)_1\), \(U(mn)_1\supset\mathrm{Spin}(2mn)_1\)): see the entry [RESOLVED] below.
\end{enumerate}
\section{Resolutions of disputed items}
\begin{enumerate}[leftmargin=1.6em,itemsep=6pt]
\item \textbf{[RESOLVED]} \emph{coefficient \(r/6\) (N1, N2, N4 \S9, N5, N7)} \hfill\textsc{orchestrator}\\ Longo--Xu Section~3.6 defines the net through \(\mathrm{CAR}(H_i)\) with \(H_i\subset L^2(I_i,\mathbb C^r)\) (finite-dimensional subspaces), i.e.\ \(r\) complex chiral fermions; Definition~3.17: \(G(I)=\frac16\bigl[\sum_{i,j}\log|b_i-a_j|-\sum_{i<j}\log|a_i-a_j|-\sum_{i<j}\log|b_i-b_j|\bigr]\); Theorem~3.18: \(S_{\mathcal A_r}(\omega,\omega_1\otimes_2\omega_2)=r\,[G(I_1)+G(I_2)-G(I_1\cup I_2)]\), ``by Lemma~3.8 it is sufficient to prove the \(r=1\) case''. For two intervals this gives \(\frac r6\log\frac{(a_2-a_1)(b_2-b_1)}{(a_2-b_1)(b_2-a_1)}\), exactly eq.\ (LX-MI) of N4. Hence V2's [BREAK] (Majorana reading, \(r/3\)) is incorrect and is withdrawn; R1's reading (complex fermions, \(c=r\)) is confirmed, consistently with the chiral half of the non-chiral \(\frac c3\log\) and with R1's lattice check. No change to any note.
\item \textbf{[RESOLVED]} \emph{graded product state (N2 eqs.\ (split-identification), (split-reference); N4 Lemma~3.1)} \hfill\textsc{V1}\\ \(\sigma_\eps\) is normal and faithful (LX Lemma~3.1(1) = Xu 2000 Prop.~2.3.1(2); \(\mathrm{dist}(\bar A_\eps,\overline{BC})=\eps>0\)) and is the unique normal state with \(\sigma(xy)=\omega(x)\omega(y)\) (\(\omega\) vanishes on odd elements), hence the relative-entropy VALUE is the same for the graded and the ordinary notation, and the Petz maps are parity preserving (\(\omega,\sigma_\eps\) are \(\operatorname{Ad}\Gamma\)-invariant). What is lost is only the factorization proof of the recovery map (N2 eqs.\ (factorization-heisenberg/schrodinger), N4 Lemma~3.1), which must be redone with the Klein transformation.
\item \textbf{[RESOLVED]} \emph{rigor/rate\_of\_remainder.tex Prop.~2.4 gap (\(f_2^\star=f_2\)), after referee RRr} \hfill\textsc{RR}\\ Closed. Use the group \(U_s=e^{s(D+ih')}\), \(h'\in\mathcal A'\): its Dyson cocycle \(\widetilde W_s^{-1}U_s\) stays localised in \(I'\) for \(s\ge0\), so it is an admissible purification family, and PA Prop.~7.6 gives the infimum \(g_B/8\) exactly. Hence \(\Phi\le-\tfrac12\log(1-2f_2\zeta^2)\) and \(c_3\le0\) are unconditional.
\item \textbf{[RESOLVED]} \emph{Withdrawal of the RR finding against quadratic\_limit.tex Problem 7.2} \hfill\textsc{RR}\\ Mis-scoped, per referee RRr. The two-sided constant must dominate the quartic Taylor coefficient, which contains \(\tfrac14\sum|(\mathcal D_2)_{\kappa\kappa'}|^2/N_{\kappa\kappa'}\approx e^{\kappa_c}/(225\kappa_c^6)\) (measured 2.8e-3, 1.1e-2, 1.1e-1, 3.7e-1, 1.31 at \(\kappa_c=9.2\dots23.0\)): P2's guess is right. Only the one-sided constant is bounded, which is all Thm.~4.2 needs.
\end{enumerate}
\section{Gaps (repairable)}
\begin{enumerate}[leftmargin=1.6em,itemsep=6pt]
\item \textbf{[GAP]} \emph{N5 Thm.~3.1 proof (citation of Ohya--Petz Cor.~5.12 / Araki77)} \hfill\textsc{R2}\\ The martingale convergence used is Araki 1977, Thm.~3.9(2), p.~181 (monotone increasing net generating \(M\), each \(N_\alpha\) approximately finite; Araki's \(S(\varphi/\psi)=D(\psi\Vert\varphi)\), arguments reversed); the AF hypothesis holds trivially for our finite-dimensional compressions but is unstated; Hiai 2018, Thm.~4.1(v), p.~12, removes AF and faithfulness. Ohya--Petz Cor.~5.12 could not be verified (book not obtainable). Replace the citation and write \((\bigcup M_n)''=M\).
\item \textbf{[GAP]} \emph{N7 \S5, last sentence of the proof of Lemma~5.1} \hfill\textsc{P3}\\ The suggested contour shift \(\xi'\to\xi'-i\pi\) is not available on the half-line contour (a vertical segment does not close); replace the remark by the series proof of rigor/kernel\_identity.tex.
\item \textbf{[GAP]} \emph{N2 eqs.\ (split-identification), (split-reference)} \hfill\textsc{R1}\\ \(\sigma_\eps=\omega_{A_\eps}\otimes\omega_{BC}\) on \(\mathcal A(A_\eps)\bar\otimes\mathcal A(BC)\) is not defined as written for a graded (fermionic) net; the product must be the twisted one of LX Lemma~3.1 (graded product \(\otimes_2\)). Fix: state the twisted product and that parity-preserving maps are unaffected.
\item \textbf{[GAP]} \emph{N1 \S3.2 (``Taking \(\eps\downarrow0\) reproduces \dots''), N2 before eq.\ (longoxu-limit)} \hfill\textsc{R1}\\ \(\lim_{\eps\downarrow0}I_\eps=F(AB,BC)\) is asserted, not argued; one line via LX Thm~4.1(5) plus continuity from inside closes it.
\item \textbf{[GAP]} \emph{N4 \S2.1 last sentence, \S1.2 item~2, eq.\ (s-t); N5 \S5 (direction of VWZ's \(\lambda\))} \hfill\textsc{R1}\\ ``Half-sided modular'' is asserted without the sign: a common LEFT endpoint gives a Wiesbrock \(+\)hsm inclusion with \(\Delta^{it}U(s)\Delta^{-it}=U(e^{+2\pi t}s)\); with \(\alpha_t=\operatorname{Ad}\Delta_{\mathcal N}^{it}\circ\alpha_0\circ\operatorname{Ad}\Delta_{\mathcal M}^{-it}\) this gives \(s_t=\lambda(1+e^{+2\pi t})\). Harmless for everything depending only on the family \(\{s_t\}\); the sign statements about VWZ's \(\lambda\) in N5 are superseded by finding [R2] (2D theory even in \(\lambda\)).
\item \textbf{[GAP]} \emph{N4 \S2.1 last sentence, \S2.2 eq.\ (s-t)} \hfill\textsc{V2}\\ The hsm sign is never identified: \(\mathcal F(B)\subset\mathcal F(D)\) is Borchers' case (B) (generator \(-K/L\le0\), \(K=xPx\)), not the ``standard'' case; \(s_t\) is asserted, not derived. V2 derived it: \(s_t=\lambda(1+e^{-2\pi t})\) is correct in the convention \(\alpha_t=\sigma^{\mathcal N}_{-t}\circ\alpha_0\circ\sigma^{\mathcal M}_t\) (two sign errors cancel); R1 obtains \(\lambda(1+e^{+2\pi t})\) in the convention \(\operatorname{Ad}\Delta_{\mathcal N}^{it}\circ\alpha_0\circ\operatorname{Ad}\Delta_{\mathcal M}^{-it}\). Since \(p(t)\) is even, nothing downstream depends on the sign; the note must state the convention and give the derivation.
\item \textbf{[GAP]} \emph{N4 \S3 Lemma~3.1, eq.\ (split-identifications)} \hfill\textsc{V2}\\ Ordinary-tensor modular data are used for a graded product; the Klein transformation is invoked but never executed; the split property is uncited. Same issue as [R1] for N2.
\item \textbf{[GAP]} \emph{N4 \S5 first display; \S10 before eq.\ (twirled-beta); \S10 Lemma~10.1} \hfill\textsc{V2}\\ (i) The \(\delta\)-term of the \(D\times D\) block and the pv-preservation under the M\"obius substitution are not discussed (both true; proved in rigor/check\_note4.tex). (ii) ``norm-dense CAR algebra'' is false (only ultraweakly dense); measurability of \(t\mapsto\overline\beta_{s_t}(x)\) is argued along nets while normality is stated for sequences (repaired via the separable predual). (iii) Alberti's variational formula is used without citation or hypotheses.
\item \textbf{[GAP]} \emph{N7 Thm~5.2 proof, last sentence} \hfill\textsc{V6}\\ ``\(\frac1{\pi^2}\int_{\mathbb R}\frac{\tanh\pi k}{k(k^2+1)}dk=\frac4{\pi^2}\)'' is off by 4: with \(u=2\pi k\), \(\int_{\mathbb R}\frac{\tanh(u/2)}{u(u^2+4\pi^2)}du=\frac1{4\pi^2}\int_{\mathbb R}\frac{\tanh\pi k}{k(k^2+1)}dk=\frac1{\pi^2}\); the boxed result \(2/(3\pi^2)\) is right.
\item \textbf{[GAP]} \emph{N5 Prop.~4.1; N5 \S6 (tail correction); C5 of the audit} \hfill\textsc{V6}\\ Prop.~4.1 is a proof sketch with the sign fixed a posteriori and numerical ratios reaching only \(1.06\); the tail correction built on it enters every reported \(\Phi\). The convention \(c(\mathcal A_r)=r\) (the ``factor 8'') is never stated. Also: \(c\) denotes both the interval length and the central charge in N3, N5 and N7 Cor.~7.1 (rename the length).
\item \textbf{[GAP]} \emph{N4 Lemma~10.1 (fidelity continuity from below), proof} \hfill\textsc{R3}\\ The variational inequality must be applied inside each \(M_n\) (Uhlmann 1976 eq.~(33)): as written only the formula in \(M\) is invoked and the convergence claimed does not bound \(\lim_nF_{M_n}\); ``it suffices to take \(m1\le x\le M_01\)'' is vacuous, Kaplansky gives a net, and the functional calculus must sit inside one \(M_n\). Repairable: for \(x\in M\) with \(\frac12(\varphi(x)+\psi(x^{-1}))\le F_M+\epsilon\), Kaplansky plus functional calculus give \(x_k\in M_{n_k}\), positive invertible, with \(\varphi(x_k)+\psi(x_k^{-1})\to\varphi(x)+\psi(x^{-1})\), whence \(\limsup_nF_{M_n}\le F_M+\epsilon\). Also: Alberti's formula needs only a W*-algebra and normal states.
\item \textbf{[GAP]} \emph{N4 Thm~7.1 (normal zero-collar map), proof} \hfill\textsc{R3}\\ Reeh--Schlieder (cyclicity of \(\Omega\) for \(\cF(I)\)) is needed to identify the restriction of the vacuum representation to \(\mathrm{CAR}(H_I)\) with the GNS representation to which Powers--St\o rmer Thm.~5.1 refers; uncited.
\item \textbf{[GAP]} \emph{N7 \S3.3 and refs/REFERENCES.md (Jen\v cov\'a 2002, Petz 1996)} \hfill\textsc{R3}\\ Jen\v cov\'a 2002 \S III opens with \(\dim H=n\) and Petz 1996 is a matrix-space theorem; neither provides monotone metrics on state spaces of general von Neumann algebras as the note suggests. N7 eqs.\ (11)--(12) do coincide with Petz's metrics for \(f=(1+t)/2\) and \((t-1)/\log t\) in Jen\v cov\'a's normalization (finite dimensions).
\item \textbf{[GAP]} \emph{numerics: defect\_matrix.py ll.~37--38 vs compression\_box.py l.~21 (N5 eq.\ (Qbox))} \hfill\textsc{V4}\\ \(\widehat D\) is built in the absolute basis \(e^{ik_j\xi}\) while \(Q_N\) is the centered formula (modes \(e^{ik_j(\xi-\xi_0)}\)); the two differ by the diagonal phases \(e^{-ik_j\xi_0}\), so \(Q_N-\widehat D\) is not exactly a compression of \(\widetilde Q_s\) and the ``rigorous lower bound'' label is void until fixed. Numerical effect on \(\Phi\): \(4.8\times10^{-5}\) relative; none on \(f_2\), \(s_2\). Fix: \(Q_N^{\mathrm{abs}}=UQ_N^{\mathrm{cent}}U^*\), \(U=\operatorname{diag}(e^{-ik_j\xi_0})\) (applied by the orchestrator, see errata).
\item \textbf{[GAP]} \emph{N5 \S6 (tail correction, ``agree to \(1.4\times10^{-4}\)''; exponents ``2.0 and 1.0'')} \hfill\textsc{V4}\\ At \(\eps=10^{-8}\) the analytic tail correction is 12\% (\(\zeta=1/15\)) to 56\% (\(\zeta=1/120\)) off because \(\kappa_c=18.4<\kappa_*(1/120)=20.7\); the agreement between the two \(\eps\) levels is a cancellation of \(+12\%/-14\%\) errors, not a validation; realistic uncertainty of the corrected \(\Phi\) is \(0.3\)--\(1\%\) at small \(\zeta\). Measured extrapolation exponents are \(1.82/1.95\) and \(0.83/1.04\), not \(2.0\) and \(1.0\); the \(s_2\) error bar \(\pm5\times10^{-4}\) is about twice too small.
\item \textbf{[GAP]} \emph{N2 \S3 eq.\ (petz-vn-equivalence)} \hfill\textsc{V1}\\ The equivalence needs the hypothesis \(D_{\cM}(\rho\Vert\sigma)<\infty\) (Jen\v cov\'a--Petz 2006 p.~14 Remark; FHSW eq.~(40)); satisfied in the application but unstated.
\item \textbf{[GAP]} \emph{N1 \S2.2 eq.\ (finite-approx); \S2.3 eq.\ (G-connected); \S3.2} \hfill\textsc{V1}\\ LX's finite-dimensional limit formula is for separated \(X,Y,Z\); the adjacent case needs LX eqs.\ (13)--(14). LX Thm~4.2 carries an unstated strong-additivity hypothesis (via LX Thm~4.4). ``Reproduces (abc-formula)'' needs LX Thm~4.1(5) plus continuity from inside.
\item \textbf{[GAP]} \emph{N5 Thm~3.1 proof (locator); Cor.~4.2} \hfill\textsc{V3}\\ Ohya--Petz Cor.~5.12 not obtainable; the correct locator is Araki 1977 Thm~3.9(2) with Thm~3.8(2)(\(\gamma\)) (transcribed with screenshot, hypotheses verified for finite-dimensional \(M_n\)). In Cor.~4.2 the finiteness of \(f_2\) needs exponential decay of the off-diagonal first-order elements, which holds because \(\widetilde F^{(1)}\) is even under \((p,p')\to(-p,-p')\) (not stated); ``fourth-order coefficient infinite'' and the collar claim are heuristic.
\item \textbf{[GAP]} \emph{N6 \S3 Observation~1 (Weil--Petersson chain)} \hfill\textsc{R4}\\ ``WP \(\Rightarrow\) implementable in the fermionic Fock space'' is not given by the cited sources: Nag--Sullivan (Thm~3.1/Cor~3.2/Prop~4.1) prove boundedness and symplecticity of the weight-0 action on \(H^{1/2}\) and their only Hilbert--Schmidt statement (Prop.~6.1) concerns \([J,M_f]\); Takhtajan--Teo Part~II (Cor.~3.9/6.4) gives HS \(\Leftrightarrow\) WP for the same weight-0 bosonic setting. Missing: the transfer to the weight-\(\frac12\) Hardy-space action and Shale--Stinespring. The first arrow (\(C^{1,1}\Rightarrow\) WP) is Shen Thm~1.1 (unnamed in the note).
\item \textbf{[GAP]} \emph{Note 5 Cor.~4.2 (upper half)} \hfill\textsc{P2}\\ \(\limsup_{\zeta\to0}\Phi/\zeta^2\le h_2=2\log2\,f_2\) is proved modulo one interchange \(\sup_n\lim_\zeta\leftrightarrow\lim_\zeta\sup_n\) (Lemma~6.7); the sharp bound \(\limsup\le g_B/8\) is open (Problem~6.8) and provably not reachable by the Hellinger route (\(1/W\ne1/N\) off-diagonal); suggested route: Bogoliubov-rotated vector representative.
\item \textbf{[GAP]} \emph{Note 5 Cor.~4.2 (rate)} \hfill\textsc{P2}\\ The stated remainder \(O(\zeta^2/\log^2(1/\zeta))\) needs a cubic estimate \(|\Phi_P-\zeta^2g_B^{(P)}/8|\le Ce^{\Lambda}\zeta^3\) for \(P\le E_\Lambda\) (Problem~7.2); together with the proved cutoff bound \(0\le g_B-g_B^{(\Lambda)}\le\frac{2r}{3\Lambda^2}(1+o(1))\) (Prop.~7.1, which confirms the fitted \(\kappa_{\max}^{-2}\) law: fit \(0.53/\Lambda^2\), bound \(0.667/\Lambda^2\)) it would give the lower half of the stated rate.
\item \textbf{[GAP]} \emph{N7 Lemma~3.1 and its proof} \hfill\textsc{P1}\\ The identity \(\eta=i(S^*-1)A\Omega\) requires the unstated hypothesis \(A\Omega\in D(S^*)=D(\Delta^{-1/2})\): the step \(\langle J\Delta^{1/2}X\Omega,A\Omega\rangle=\langle\Delta^{1/2}JA\Omega,X\Omega\rangle\) needs \(JA\Omega\in D(\Delta^{1/2})\), and conversely if \(A\Omega\notin D(S^*)\) then \emph{no} vector \(\eta\) represents \(\varphi\). Item~(2) (``if \(A\in\mathcal M\) then \(S^*A\Omega=\Delta A\Omega\)'') needs \(A\Omega\in D(\Delta)\), which is not automatic for bounded \(A\in\mathcal M\) (only \(\mathcal M\Omega\subset D(\Delta^{1/2})\) is). Item~(3), \(g_B=2\|(1-J)(1+\Delta)^{-1/2}A\Omega\|^2\), needs nothing and should be taken as the definition. No numerical conclusion of N7 is affected.
\item \textbf{[GAP]} \emph{N7 Lemma~3.2, Kubo--Mori half; N7 \S8 item 1} \hfill\textsc{P1}\\ Proved only in finite dimensions (\texttt{lemma\_second\_variation.tex} Prop.~5.5: \(D(\rho\|\rho_s)=\frac{s^2}{2}g_{KM}+O(s^3)\), by the integral representation of \(\log\) and \(\int_0^\infty(t^2-p_ip_j)(p_i+t)^{-2}(p_j+t)^{-2}dt=0\); the second-order term of \(\rho_s\) drops out). For type III no analogue of Alberti's formula is available: Kosaki/Petz's \(D=\sup_h\{\omega(h)-\log\varrho(e^h)\}\) has second variation \(\frac14 g_B\), not \(\frac12 g_{KM}\); Araki77 Thm.~3.9 gives only \(\liminf s^{-2}D\ge\frac12 g_{KM}\) (lower semicontinuity, wrong direction). Missing precisely: an upper bound \(D(\omega\|\omega_s)\le\frac{s^2}{2}g_{KM}+o(s^2)\) for \(\omega_s=\omega\circ\mathrm{Ad}\,e^{isA}\).
\item \textbf{[GAP]} \emph{N7 \S3.2--\S4 (zero-collar \(\xi\)-frame), \S8 items 2 and 4} \hfill\textsc{P1}\\ With \(\tilde g=-4\sinh^2(\xi/2)\mathbf 1_{\xi>0}\) the vector \(T(g)\Omega\) has infinite norm (\(|\hat g|^2\sim4/k^2\) at \(k\to0\), \(\hat W(0)=c/48\pi^2\ne0\); agent V6), so the lemma's hypothesis (V) fails at zero collar. The \emph{upper} bound on \(\Fid\) survives verbatim (it uses no vector, only \(g_B(\varphi)=4\sup_K[\varphi(K)-\mathrm{Var}_\omega(K)]\); Rmk.~10.1). Missing precisely: (a) finiteness of \(g_B(\varphi)\) for the zero-collar tangent functional; (b) its identification with N7's spectral integral. Prop.~10.2 reduces (a)+(b) to: cutoffs \(g_n\) with \(T(g_n)\Omega\in\mathcal H\), \(\varphi_n\to\varphi\) pointwise, \(\sup_n g_B(g_n)<\infty\), and uniformity of \(\varphi_n\to\varphi\) on Bures balls \(\{\omega(K)=0,\|K\Omega\|\le R\}\); the last is the real work.
\item \textbf{[GAP]} \emph{N7 \S4.2 (sentence after eq.\ (15)) and \S1 (``\(\xi\mapsto\xi\pm2\pi t\)'')} \hfill\textsc{V5}\\ The modular-flow sign is fixed by ``the KMS asymmetry \(d\mu(-u)=e^ud\mu(u)\) of (13)'', but that symmetry is a theorem only for \(A\in\cM\), and \(T(g)\notin\cM\) for every admissible representative (Prop.\ 3.4 of rigor/check\_notes\_6\_7.tex). The sign the note uses is nevertheless correct and is derived unconditionally there: \(\cA((-\infty,L))=\cA((L,\infty))'\) (Haag duality, BGL93 Thm~2.19(b)) plus Bisognano--Wichmann give \(\Delta^{it}\widetilde T_0(\xi)\Delta^{-it}=\widetilde T_0(\xi-2\pi t)\), hence \(u=\log\lambda=+2\pi k\). Replace the \(\pm\) by \(-\) and drop the KMS appeal.
\item \textbf{[GAP]} \emph{N7 \S4.2--4.3 (use of eq.\ (13)); Lemma~3.2 hypotheses} \hfill\textsc{V5}\\ Confirms [V6] and sharpens it: it is not only \(\|T(g)\Omega\|=\infty\) (the \(k^{-2}\) pole against \(\widehat W(0)=c/48\pi^2\)) but that NO representative with \(g=f\) near \(\overline D\) is affiliated with \(\cA(I)\) --- the \(\xi\)-chart covers only \(x<L\), so the cutoff must sit outside it --- hence eq.\ (13), the \(A\in\cM\) spectral formula that \S4.3 integrates, is unavailable. The correct general form is Lemma~3.1(3), which is representative-independent by Lemma~3.1(1); it equals eq.\ (13) only through the KMS symmetry, and \(4\int\frac{1-\lambda}{1+\lambda}d\mu\) is merely conditionally convergent at \(k=0\) while its symmetrization is absolutely convergent. So \S4 is a formal computation validated a posteriori by \S5. Lemma~3.2 separately: ``affiliated with a von Neumann algebra containing \(\cM\)'' is vacuous (true of every self-adjoint operator); ``\(o(s)\) controlled in the norm of the tangent space'' is undefined; the form-domain condition is needed only for \(g_{\mathrm{KM}}\).
\item \textbf{[GAP]} \emph{N7 \S4.4 (``the agreement of the two fixes the constant unambiguously''); Cor.\ 6.1 preamble} \hfill\textsc{V5}\\ The mathematics is right (linearity in \(c\) is proved in rigor/check\_notes\_6\_7.tex Prop.\ 4.6, with no hidden dependence on the Schwarzian constant or on the continuation of \(g\) beyond \(L\), and \(JT(g)J=T(g\circ r)\) is correct), but the narration misstates the logic: what determines \(\kappa=2/(3\pi^2)\) is linearity in \(c\) plus Theorem~5.2 at \(c=1\); \S4 is a consistency check on the continuation prescription, not an independent determination. Theorem~B is therefore conditional on Lemma~3.2 ALONE, independently of the formal status of \S4.2--4.3. Cor.\ 6.1's ``rigorous for the free fermion'' overstates: what is rigorous is the pair of metrics, not \(\lim\Phi/\zeta^2\).
\item \textbf{[GAP]} \emph{N6 \S3 Route B (compactness and identification)} \hfill\textsc{V5}\\ ``the set \(\{\psi:\langle\psi,L_0\psi\rangle\le E\}\) is norm-compact'' is false as written (it contains \(\mathbb C\Omega\)); the true statement needs \(\|\psi\|\le1\), and finite-dimensional eigenspaces alone do not give compact resolvent (one also needs eigenvalues accumulating only at \(+\infty\), which does hold for a unitary VOA of CFT type). ``the implementers of diffeomorphisms that are trivial on a region commute with its algebra'' is a theorem, uncited: \(\mathrm{supp}\,\varphi\subset K\Rightarrow U(\varphi)\in\cA(K)\), Carpi--Weiner CMP 258 (2005) Prop.~3.7 / DFK LMP 67 (2004). Strong additivity and Haag duality are used in the text but absent from Conjecture~3.2's hypotheses.
\item \textbf{[GAP]} \emph{N7 \S3.2 (``used here formally, as a computational prescription''); confirms [V5], [V6]} \hfill\textsc{P4}\\ An \(x\)-space diagnosis localises the obstruction: the Mellin--Barnes step needs \(\int_0^1 g\,u^{-\sigma}du\) (\(\mathrm{Re}\,\sigma<1\)) and \(\int_0^1 g\,u^{\sigma-4}du\) (\(\mathrm{Re}\,\sigma>3\)) simultaneously; the strips are disjoint precisely because \(g(L)\neq0\). A tangent field vanishing to second order at \(x=L\) would make the identity absolutely convergent.
\item \textbf{[GAP]} \emph{fidelity\_tables.tex Table~3 (range of the tabulated \(\Phi\))} \hfill\textsc{PL}\\ The two-branch law needs \(\Phi\) up to \(\zeta=z(1+e^{\pi\lambda})\approx112\,z\) at \(\lambda=1.5\), but \(\Phi\) is tabulated only for \(\zeta\le0.533\). The lattice supplies \(\Phi\) to \(\zeta\approx23\) and shows \(\Phi\sim\zeta^{1.1}\) there, far flatter than the \(\zeta^{1.71}\) of the last tabulated interval. A continuum computation of \(\Phi\) at \(\zeta\gtrsim1\) is needed.
\item \textbf{[GAP]} \emph{rigor/uhlmann\_upper\_bound.tex Lem.~7.3 (the tangent vector)} \hfill\textsc{PR}\\ \(d\Gamma(D)\Omega\) does not exist: the diagonal block contributes \(\operatorname{Tr}(P_+DP_+)=\infty\). The intended vector is \(:\!d\Gamma(D)\!:\Omega\leftrightarrow P_-DP_+\). In a finite-dimensional test the omission is invisible in \(\varphi\) (that trace is imaginary) but destroys \(a=Ea\) and \(\|a\|=\|P_-DP_+\|_2\), on which Prop.~7.5 rests. Two-line repair: define \(a\) directly.
\item \textbf{[GAP]} \emph{rigor/uhlmann\_upper\_bound.tex Thm.~4.1, middle term of eq.~(16)} \hfill\textsc{PR}\\ \(|\det(P_+UP_+)|\) is not a Fredholm determinant of \(1+\)trace class: \(V\in\mathfrak S_2\) makes only \(X^*X=1-V^*V\) of that form. The middle equality is neither well posed nor proved; only the outer one is used and proved. Delete it, or add the hypothesis \(P_+(U-1)P_+\in\mathfrak S_1\).
\item \textbf{[GAP]} \emph{rigor/uhlmann\_upper\_bound.tex Lem.~5.1 \(\langle1\rangle4\) (core of \(D\))} \hfill\textsc{PR}\\ Neither \(C_c^\infty\) nor \(C_c^1\) is \(\widetilde W_s\)-invariant for a \(C^{1,1}\) flow, since \(\widetilde W_sf\) involves \(\widetilde k_s''\in L^\infty\) only; Nelson's invariant-domain criterion therefore does not apply to the domain named. Use \(H^1_c(\mathbb R)\), which is dense, invariant and in \(D(D)\). Nothing downstream changes.
\item \textbf{[GAP]} \emph{rigor/tail\_bound.tex \S5 (the \(\kappa_c^{-2}\) tail of open\_problems\_plan \S2.1)} \hfill\textsc{PT}\\ The conjectured \(C\zeta^2\kappa_c^{-2}\) is NOT proved. The Bures triangle inequality contributes \(2\sqrt{\Phi\Phi_\partial}=\Theta(\zeta^2/\kappa_c)\); this is intrinsic to a metric argument. Obstruction isolated: a bound on \(\|\sqrt{\widetilde Q}-\sqrt{X_{11}\oplus X_{22}}\|_2\) quadratic in the off-diagonal block \(X_{12}\).
\item \textbf{[GAP]} \emph{rigor/tail\_bound.tex \S7 (certified error bars)} \hfill\textsc{PT}\\ Certified interval for \(\Phi(1/15)\): \([3.487,6.53]\times10^{-5}\), a factor 1.9, versus Note~5's heuristic \(0.3\)--\(1\%\); widths 82--154\% over the scan. Upper half also conditional on the box (no bound on \(\Phi_E-\Phi_{P_{\rm sub}}\)). Arithmetic margin \(\approx3\%\) at \(\eta=10^{-9}\); an interval audit of \(\widehat D,Q_N\) would give \(0.4\%\). [corrected 2026-09-24: these bars are not certified. Their upper end is \texttt{eq:final} of Cor.~5.2 at \(\kappa_{\rm eff}\), whose constants \(0.173\) and \(0.88\) rest on numerical inputs (\(A_1\le2.0\) and \(A_3\le0.0374\) in double precision, not certified; \(f=0.0085\)); rigor/tail\_bound.tex \S7 and Table~1 now say so (its Corrections~C8 and C10); card \texttt{tail-bound}]
\item \textbf{[GAP]} \emph{numerics/README.md \S"Large-zeta systematics" (window tail law)} \hfill\textsc{PZ}\\ The window tail exponent in \(\Phi^{\rm sub}=\Phi_\infty-b\kappa_c^{-p}\) is not universal: \(p=2.49,2.31,1.98,1.76,1.21\) at \(\zeta=1.07\)--\(17.07\). The small-\(\zeta\) formula \(\zeta^2/(15\kappa_c^2)\) (i.e.\ \(p=2\)) is therefore invalid for \(\zeta\gtrsim4\), and \(p=1\) extrapolations overshoot by up to 4\%.
\item \textbf{[GAP]} \emph{rigor/implementation\_C11.tex Prop.~4.9 vs Note~6 Conj.~3.2 (finite-energy vector)} \hfill\textsc{QB}\\ \(\lVert g_{\rm tot}'\rVert_{3/2}=\infty\) (partial sums \(\tfrac{2|J|}{\pi}\sqrt N\)), hence \(T(g_{\rm tot})\Omega\in D(L_0^{1/2})\setminus D(L_0)\): the second energy moment diverges logarithmically, coefficient \(cJ^2/(48\pi^2)\). ``The recovered state is the vector state of a finite-energy vector'' is established only to order \(s^2\) (coefficient \(1.6227\,c\) for \(L=1\)), not for fixed \(s\).
\item \textbf{[GAP]} \emph{rigor/implementation\_C11.tex \S6 (Route B, energy compactness)} \hfill\textsc{QB}\\ Confirmed: the compact set needs both the norm bound and finite \(L_0\)-multiplicities. The uniform energy bound is \(E(\gamma)=\frac c{24\pi}\int\{\gamma,z\}iz\,dz\), finite for \(C^{1,1}\) maps because the Schwarzian is a finite measure with singular part \(J_\gamma\delta_{\theta_0}/\gamma'(\theta_0)\); the weak-\(*\) convergence of the approximants' Schwarzians is not written out.
\item \textbf{[GAP]} \emph{universality\_normalization.tex \S3.2 and phase2\_brief.md (B-value)} \hfill\textsc{QA}\\ The brief's claim that the smooth extension makes the \(x\)-space Mellin evaluation legitimate is wrong. The disjoint-strip obstruction is caused by \(g(L)\neq0\), which the extension does not change from inside \(I\); and \(\mu\) there is the measure of \(T(g\mathbf 1_I)\Omega\), which is not a vector. Value confirmed by other routes.
\item \textbf{[GAP]} \emph{universality\_normalization.tex \S3, use of Lemma~2.1(2)} \hfill\textsc{QA}\\ \(g_B=2\int(\lambda-1)^2(\lambda+1)^{-1}d\mu\) requires \(A\in\mathcal M\); for \(A=T(g_{\rm tot})\) this fails, \(\operatorname{supp}g_{\rm tot}\not\subset I\), and the KMS symmetry \(d\mu(\lambda^{-1})=\lambda d\mu(\lambda)\) is unavailable. The unconditional form is \(\frac14 g_B=\lVert ua\rVert^2-\operatorname{Re}\langle ua,Jua\rangle\), \(u=(1+\Delta)^{-1/2}\).
\item \textbf{[GAP]} \emph{rigor/universality\_theorem\_B.tex Lem.~5.6, 5.7 (standard subspaces)} \hfill\textsc{QA}\\ The variational form \(\frac14 g_B=\inf_h\lVert T(g_{\rm tot}+h)\Omega\rVert^2\) over real \(h\) supported in \(I'\) needs one-particle Reeh--Schlieder for \(T\)-vectors and one-particle Bisognano--Wichmann for \(H_T(I')\subset\mathcal H_T\). Cited to Brunetti--Guido--Longo 1993 and Longo's notes; neither source obtainable in this session. Theorem~B does not depend on it.
\item \textbf{[GAP]} \emph{rigor/uhlmann\_upper\_bound.tex (hypothesis (H2) for the free fermion)} \hfill\textsc{QA}\\ Norm differentiability of \(s\mapsto\Gamma(\widetilde W_s)^*\Omega\) is never asserted there; the upper bound goes through the exact determinant instead. All ingredients are present (Thouless form, \(\lVert G_s-sG^{(1)}\rVert_2\le s\epsilon(s)\)), so the gap is small, but (H2) for the model case is not written down. [corrected 2026-09-24: resolved after this finding: at QA's request that document then added Section~10, whose Lemma~10.1 (\texttt{lem:H2}, ``(H2): norm differentiability on the vacuum'') gives \(\Gamma_s\Omega=\Omega+is\,T(g_{\rm tot})\Omega+o(s)\) in norm for \(|s|\) small, for \(s<0\) since that document's Correction~C6, and \(\Gamma_s^*=\Gamma_{-s}\) gives the form above; this is (H2) for one component, and only the \(r\)-fold direct sum for \(r\ge2\) stays open (universality\_theorem\_B.tex, Correction~C2)]
\item \textbf{[GAP]} \emph{rigor/kubo\_mori\_second\_variation.tex \S``What remains open'' (Gap QK1)} \hfill\textsc{QK}\\ The upper bound needs a uniform-integrability hypothesis, not just \(g_{\mathrm{KM}}<\infty\). Missing: \(\lim_N\limsup_s s^{-2}\int_{\lambda<1/N}(-\log\lambda)\,d\nu_s=0\). The analogous statement one logarithm lower is free, from the exact sum rule \(\int(\lambda-1)d\nu_s=0\). The gap is exactly one logarithm wide.
\item \textbf{[GAP]} \emph{rigor/kubo\_mori\_second\_variation.tex \S``What remains open'' (Gap QK2)} \hfill\textsc{QK}\\ Note~7's curve has \(A=T(g_{\rm tot})\) not affiliated with \(\mathcal A(I)\), so \(e^{isA}\notin\mathcal M\) and \(\Delta_{u_s^*\Omega,\Omega}\ne u_s^*\Delta u_s\): the exact formula fails and both bounds are lost. Needed: a Kubo--Mori gauge statement modulo \(\overline{\mathcal M'_{\rm sa}\Omega}\), the analogue of \(\frac14 g_B=\mathrm{dist}(a,\overline{\mathcal M'_{\rm sa}\Omega})^2\).
\item \textbf{[GAP]} \emph{rigor/implementation\_C11.tex Thm.~4.4 steps \(\langle2\rangle2\), \(\langle1\rangle5\)} \hfill\textsc{QR2}\\ \(U(\mathrm{Exp}f)=e^{iT(f)}\) ``up to a phase'' is cited from CDIT \S2.1, stated only on the irreducible \(\mathcal H(c,h)\). On a reducible positive-energy representation the discrepancy is a priori \(\bigoplus_je^{i\alpha_j}\), not a scalar, and the localisation step (hence all of (H1)) needs a scalar. Repair: global lift to the Virasoro group (FH05 Prop.~5.1), where the centre acts as the common \(c\).
\item \textbf{[GAP]} \emph{rigor/universality\_theorem\_B.tex, (H2) and Lemma 5.4 (\(a\in\mathcal H_T\))} \hfill\textsc{QR1}\\ Nothing in (H1)--(H3) says \(T(g_{\rm tot})\Omega\) is the limit of \(T(f_n)\Omega\) for smooth \(f_n\to g_{\rm tot}\) in \(H^{3/2}\); the mollification proof of Lemma 5.4 presupposes it. Since \(a\in\mathcal H_T\) carries all of Theorem 5.6, add that clause to (H2). It follows from the Carpi--Weiner energy bound.
\item \textbf{[GAP]} \emph{rigor/universality\_theorem\_B.tex, Lemma 1.3 \(\langle1\rangle2\) and the locality clause of (H1)} \hfill\textsc{QR1}\\ \(\varphi=\tilde k_s\circ(\tilde k_s^\sharp)^{-1}\) is the identity on \(J_s\) but generally moves points arbitrarily close to \(\partial J_s\), so \(\operatorname{supp}\varphi=\overline{J_s'}\) and no interval \(K\) with \(\operatorname{supp}\varphi\subset K\subseteq J_s'\) exists. Restate (H1) as \(U(\mathrm{Diff}(K))\subseteq\mathcal A(K)\), \(\mathrm{Diff}(K)=\{\varphi:\varphi|_{K'}=\mathrm{id}\}\).
\item \textbf{[GAP]} \emph{rigor/universality\_theorem\_B.tex, Theorem 6.1 \(\langle1\rangle2\) and \S7.1 (graded net)} \hfill\textsc{QR1}\\ ``\(\mathcal F\) is Haag dual (twisted duality fails for the odd part)'' is backwards: Haag duality fails, twisted duality \(\mathcal F(I)'=Z\mathcal F(I')Z^*\) holds (Foit; used in uhlmann\_upper\_bound.tex). Repair: Theorem 5.6 needs no (HD); \(\mathrm{Ad}Z\) is trivial on even operators, so \(JT(f)J=T(f\circ r)\) survives.
\item \textbf{[GAP]} \emph{rigor/rate\_of\_remainder.tex Prop.~2.4 (\(f_2^\star=f_2\) over flow extensions)} \hfill\textsc{RR}\\ Only \(f_2\le f_2^\star\le1.00036\,f_2\) is established (numerically, p1\_bogoliubov.out (4)). PA Step 5 computes the infimum over bounded commutant generators \(h'\); a field extension gives the unbounded generator \(iD_\eta\). Equality needs the minimiser in \(H_{\mathcal M'}\) to be a limit of field tangents.
\item \textbf{[GAP]} \emph{tail\_bound.tex, Thm.~5.1 and Cor.~5.2} \hfill\textsc{RC}\\ The Bures cross term $2\sqrt{\Phi\Phi_\partial}$ is removable. Choosing the optimal Uhlmann unitary of the pinched pair (block diagonal) gives $\Phi-\Phi_W\le\Phi_{22}-\log\lvert\det(1+T)\rvert$ with no term of order $\sqrt\Phi$; numerically this is $50$--$160$ times smaller than $0.173\zeta^2/\kappa_c$.
\item \textbf{[GAP]} \emph{certified\_numerics.tex, \S5 and \S7} \hfill\textsc{RC}\\ Two inputs remain unproved: the two-dimensional quadrature for $\widehat D$ (validated at $10^{-12}$ per entry, not enclosed), and the box cutoff beyond $\kappa_{\max}$, which is now the dominant uncertainty. Both routes to close them are described in \S5(a),(b).
\item \textbf{[GAP]} \emph{universality\_normalization.tex, eq.~(11) and Theorem 5.2 ($g_{\mathrm{KM}}=c/6$)} \hfill\textsc{RK}\\ The value $c/6$ is the analytic continuation of a divergent integral: in the modular frame of $I$ the field is $-4\sinh^2(\xi/2)$ for $\xi>0$, so $\int(\tilde g''^2+\tilde g'^2)=\infty$ and $\hat{\tilde g}(k)=-2i/(k(k^2+1))$ renormalises it to $2$. Any proof of $s_2=c/12$ must justify that renormalisation.
\item \textbf{[GAP]} \emph{kubo\_mori\_gauge\_lemma.tex, \S4.1 (Gap RK1)} \hfill\textsc{RK}\\ The second-order evaluation of the monotonicity bound needs uniform integrability of $\log_-$ along the orbit, $\lim_N\sup_{\lvert\sigma\rvert\le s_0}\int_{(0,1/N)}(-\log\lambda)d\rho_\sigma=0$ with $\rho_\sigma$ the spectral measure of $\hat U_\sigma T(\hat g)\Omega$. Proved at $\sigma=0$; uniformity in $\sigma$ open.
\item \textbf{[GAP]} \emph{certified\_numerics.tex, \S7 and Table 3 (lower end $\Phi_W$)} \hfill\textsc{RCr}\\ ``All numbers below use the ball $Q_N$'' is false: $\Phi_W$ is read from results\_hp2.txt, produced from the double-precision trapezoid Q\_box that \S5 declares inaccurate. Recomputation at 60 digits gives $5.744435\times10^{-7}$ (ball) against $5.744386\times10^{-7}$ (trapezoid) at $\zeta=1/120$: harmless ($8.6\times10^{-6}$ relative) but undisclosed and unbounded.
\item \textbf{[GAP]} \emph{certified\_numerics.tex, Corollary 2.6, step $\langle1\rangle3$} \hfill\textsc{RCr}\\ The sum over the eigenbasis of $X_{22}$ is split into $E^{+}$ and $E^{-}$ parts, which presupposes $E^{+}\widetilde QE^{-}=0$; $E^{\pm}$ are spectral projections of the vacuum modular Hamiltonian, not of $\widetilde Q$. Also the $E^{-}$ weight is $\le(1-\tau_-)^{-1}$, not $\le1$. Repair: split by $b_j\lessgtr1/2$, factor 2. Conclusion verified numerically (ratio $0.71$).
\item \textbf{[GAP]} \emph{rate\_of\_remainder.tex, Prop.\ 2.6 and Thm.\ 2.5 (the hypothesis $f_2^\star=f_2$)} \hfill\textsc{RRr}\\ Theorem 2.3 needs only a one-parameter unitary group with Hilbert--Schmidt $[P_-,K]$, not a flow. $U_s=e^{s(D+ih')}$, $h'$ in PA Def.\ 5.3, is such a group and is admissible: its Dyson cocycle stays localised in $I'$ because the flow shrinks $I'$. PA Prop.\ 7.6 then gives $f_2^\star=f_2$ exactly, unconditionally.
\item \textbf{[GAP]} \emph{rate\_of\_remainder.tex, Def.\ 2.1 (F3) and Lemma 2.8 $\langle1\rangle3$} \hfill\textsc{RRr}\\ Near-optimal fields are asymptotic translations ($\chi\to$ const at $+\infty$; p1\_bogoliubov.out (2)), so PA Lemma 2.4 and Prop.\ 2.5, which assume $\phi=\mathrm{id}$ off a compact set, do not apply and must be re-run. Lemma 2.8 argues ``for $y$ outside $\mathrm{supp}\,\chi$'' and so covers only compactly supported fields.
\item \textbf{[GAP]} \emph{Scope of the \(\eta^2\) recovery law (O7(a)) away from the fixed point; numerics/lattice/petz\_thermal\_rg.out \S(2b),(3)} \hfill\textsc{S9}\\ Off criticality \(-\log F^{(0)}\) is linear, not quadratic, in the CMI: the ratio \(-\log F^{(0)}/I\) rises to \(0.10\)--\(0.17\) for \(\mu L_B\ge1.6\) (from \(0.02\)--\(0.05\) at \(m=0\)) and persists to \(I\sim10^{-19}\). Any general \(\eta^2\) bound must use adjacency and conformal invariance.
\item \textbf{[GAP]} \emph{rigor/optimality\_second\_order.tex \S5.2, Lemma~5.5; open\_problems\_plan.tex \S O8 ``Definition of done''} \hfill\textsc{S8}\\ No nonzero lower bound on the optimal recovery error is established. Gauge twirling rigorously reduces to gauge-covariant channels but not to quasi-free ones, and Reeh--Schlieder makes the ``recovered \(A\)--\(C\) correlations factor through \(B\)'' argument vacuous for non-Gaussian channels. Sharpened open question: are all achievable recovered symbols quasi-free-achievable?
\item \textbf{[GAP]} \emph{rigor/optimality\_all\_channels.tex \S5.1; numerics/optimality\_all/s10\_nonpert\_cap1e8.out, s10\_nonpert\_cap1e12.out, exact\_fidelity\_rotated.py} \hfill\textsc{ORCH-3b}\\ S10's refutation of the compression's optimality is established at the level of the SLD form with weight caps \(\le10^{8}\) (gain factor \(0.705\)--\(0.74\)); with cap \(10^{12}\) the uncapped form is dominated by the modular ultraviolet (order-of-limits effect: \(g_Q(\delta_c)/\frac12s^2\lVert u\rVert^2=1.19\)--\(5.9\) for \(s=0.001\)--\(0.005\) on the \(L=512\) grid) and shows no gain at \(s=10^{-3}\). The claim therefore rests on the exact root fidelity of the rotated channel, computed by the orchestrator in high precision (stable SVD form of Note~5 Thm.~2.1, validated against brute force to 12 digits); its result decides the verdict.
\item \textbf{[GAP]} \emph{rigor/exact\_fidelity\_upper\_half.tex Thm.~7.2 (the numerical rider ``\(E_{\mathrm{rec}}\le0.71f_2z^2\)'')} \hfill\textsc{RE1}\\ The theorem proves \(E_{\mathrm{rec}}\le(1-\theta_{\mathrm{fr}})f_2z^2(1+o(1))\) with \(\theta_{\mathrm{fr}}\) the supremum of \(\cos^2\angle(u,v_{\mathcal G})\) over \emph{finite-rank} anti-self-adjoint \(\mathcal G=E_D\mathcal GE_D\), and \(\theta_{\mathrm{fr}}>0\iff E_DuE_D\ne(E_DuE_D)^{*}\). But \(\theta_{\mathrm{fr}}\le\theta\) (S10's supremum is over the larger bounded class), so quoting \(\theta=0.300\pm0.005\) gives \(1-\theta_{\mathrm{fr}}\ge1-\theta\), not \(\le\): the numerical conclusion needs \(\theta_{\mathrm{fr}}\ge0.30\), which is not stated. Repair (no new mathematics): (i) all quantities depend on \(\mathcal G\) only through \(w=P^{\perp}\mathcal GP\) and both numerator \(\Re\langle u,w\rangle\) and denominator \(\lVert\Pi w\rVert^2\) are \(\mathfrak S_2\)-continuous, so \(\theta_{\mathrm{fr}}\) equals the supremum over Hilbert--Schmidt \(\mathcal G\); (ii) \(\mathcal G=R_a=E_DuE_D\) is Hilbert--Schmidt and admissible, giving the explicit \(\theta_{\mathrm{fr}}\ge\lVert R_a\rVert_2^4/(\lVert u\rVert_2^2\langle R_a,\mathcal AR_a\rangle_r)>0\) (S10 Cor.~11.2, numerically \(0.0143\)); (iii) both discretisations that produce \(\theta\approx0.29\)--\(0.34\) (S10 \S7, optimality\_all\_channels \S5.2) are finite dimensional, so their optimisers are finite rank and \(\theta_{\mathrm{fr}}\ge0.30\) is exactly as well supported numerically as \(\theta\ge0.30\).
\item \textbf{[GAP]} \emph{rigor/exact\_fidelity\_upper\_half.tex Thm.~7.1 \(\langle1\rangle3\) and Cor.~7.3} \hfill\textsc{RE1}\\ The \(g_Q\) form of (UH), ``\(-\log F\le g_Q(\delta(s))(1+o(1))\)'', is justified by ``S10 Lemma~6.1(2) together with the continuity of the positive quadratic form \(g_Q\)''. \(g_Q(\delta)=\frac18\sup_t(\mathrm{Tr}\,t\delta)^2/\mathrm{Tr}(tQt(1-Q))\) is \emph{not} continuous in the continuum: \(\mathrm{spec}\,Q\) accumulates at \(0\) and \(1\) for an interval, which is the order-of-limits effect documented in S10 \S7(d) (uncapped \(g_Q(\delta_c)/\frac12s^2\lVert u\rVert^2=2.8,12,45\) at \(s=0.02,0.05,0.1\)) and in optimality\_all\_channels \S5.1 (uncapped \(g_Q\sim10^{283}\)). What \emph{is} proved, unconditionally, is the \(g_Q\)-free form \(\limsup_{s\downarrow0}s^{-2}[-\log F]\le\frac12\lVert\Pi P^{\perp}(D_\chi+\lambda\mathcal G)P\rVert_2^2\) (Thm.~7.1 \(\langle1\rangle1\)--\(\langle1\rangle2\), the order limsup-then-infimum being the correct one), and this is all that Thm.~7.2 consumes, so the main result is unaffected. Repair: state Thm.~7.1 and the upper half of Cor.~7.3 in the \(g_Q\)-free form and add the \(g_Q\) identification as a remark with S10's proviso. Separately, the lower half of Cor.~7.3 cites ``optimality\_all\_channels Prop.~2.4 and Thm.~4.1'' for \(E_{\mathrm{rec}}\ge\mathcal D_z(1-o(1))\); that document has no Prop.~2.4, and its Thm.~4.1(2) is conditional on test families with \(\lVert t_z\rVert_1^2/V(t_z)\) bounded and \(\Delta_z(t_z)\lVert t_z\rVert_1/V(t_z)\to0\) (unconditionally only the dual form, its eq.~(10)).
\end{enumerate}
\section{Minor issues}
\begin{enumerate}[leftmargin=1.6em,itemsep=6pt]
\item \textbf{[MINOR]} \emph{N2 \S5 eq.\ (FH-bound)} \hfill\textsc{R2}\\ The averaged form \(-2\log F(\rho,\rho|_N\circ\bar\alpha_\sigma)\) is FH22 Thm.~2 eq.~(21) (from the integrated FHSW Thm.~1 eq.~(24), \(p(t)=\pi/(\cosh2\pi t+1)\), by concavity); attribute correctly and add the concavity step. FHSW hypotheses: \(\rho,\sigma\) normal, \(\sigma\) faithful (non-faithful: Cor.~1).
\item \textbf{[MINOR]} \emph{N4 eq.\ (s-t) and the ``alternative convention'' remark} \hfill\textsc{R2}\\ FHSW \S4.2 eqs.\ (166)--(168) as printed are internally inconsistent (orientation of the hsm condition, reversed adjoints, missing sign), so the sign of \(t\) cannot be certified from FHSW; it is certified by N4's own Borchers-relation derivation (\(+\)hsm, \(\kappa=+1\)), so the hedge can be replaced by a statement of the convention actually derived.
\item \textbf{[MINOR]} \emph{N2 \S3, \S5; N5 \S6--7} \hfill\textsc{R2}\\ ``normal recovery channel'' should read ``normal unital 2-positive map'' (Petz sufficiency); SFR Thm.~2.1 gives no map structure, the universal rotated-Petz average is JRSWW Thm.~2.1; FR/SFR use \(\log_2\) (bits), our notes use nats; ``\(t\to+\infty\) optimal'' and ``factor 4'' are chiral-only statements.
\item \textbf{[MINOR]} \emph{N7 \S5, ``Weights'' paragraph} \hfill\textsc{P3}\\ \(q_\kappa-q_{\kappa'}=2\sinh(u/2)e^{v/2}/[(1+e^\kappa)(1+e^{\kappa'})]\) is missing a minus sign (it equals \(-2\sinh(u/2)\dots\) for \(u=\kappa-\kappa'\)); harmless since only the square and the correctly signed second line of eq.~(22) are used.
\item \textbf{[MINOR]} \emph{N2 main box (``\(r\) components''), N4 \S1.1; N1 \S2.1; N4 eq.\ (hardy-kernel), N5 after eq.\ (fermi)} \hfill\textsc{R1}\\ Say ``\(r\) complex (Dirac) chiral fermions, \(c=r\)''; the bound \(F_{\mathcal B}\le F_{\mathcal A_r}\) is in the proof of LX Thm~4.1(6), not its statement; name the Fourier convention and record \(\kappa=-2\pi s\) relative to Casini--Huerta's \(s\). Six misprints in the sources (LX (7), \S4.2.1 \(\eta\), Thm~4.2(1) sign, Thm~4.2(2) \(o(\eps)\), Thm~4.4 missing \(\ln\), BGL p.~9) documented; none propagates.
\item \textbf{[MINOR]} \emph{N4 \S1.1, Lemma~2.1, \S4, \S7 Thm~7.2, \S8, \S10} \hfill\textsc{V2}\\ Fourier convention unstated (\(P_+\) projects on \(k<0\), which is the positive-energy space); Lemma~2.1 states isometry, unitarity onto is what is used; a premature restriction claim in \S4; Thm~7.2 uses surjectivity of \(\beta_s\) implicitly; \(\|\Delta_t\|_1\to\infty\) as \(t\to-\infty\) should be remarked; concavity of root fidelity and Fuchs--van de Graaf uncited; strong additivity at the touching point unmentioned. Verified numerically: \(\|R_s\|_1=\tfrac12\log(1+as/(a+L))\) to 10--11 digits (three parameter sets), positivity of \(H_{\alpha,\beta}\), \(R_s\cong H_{\alpha,\beta}\) (numerics/check\_note4\_tracenorm.py).
\item \textbf{[MINOR]} \emph{N4 eq.~(45), Thm~7.1, Lemma~10.1; N5 \S1, Prop.~2.3; cross-references} \hfill\textsc{R3}\\ N4 eq.~(45) assumes \(S-T\in\mathfrak S_1\), stronger than Powers--St\o rmer Lemma~4.1 requires; automatic normality of von Neumann \(*\)-isomorphisms uncited; N5 \S1: the standard-form supremum is Alberti--Uhlmann 2002 Thm.~2(3) and the arithmetic-mean form is Uhlmann 2011 eq.~(29), not Alberti's own; N5 Prop.~2.3 uses a corollary of Daleckii--Krein, not the formula itself; N4's results are numbered Theorem~7.1 and Lemma~10.1 (N5/N7 cite ``7.2'' and ``9.2''). Verified: PS Def.~1.2 / Lemma~4.1 / Thm.~5.1 (no faithfulness needed), Araki 1971 Thm.~1 and the reduction \(S=A\oplus(1-A)\), pulled-back symbol \(=W_s^*Q_{J_s}W_s\), quasi-equivalence \(\Rightarrow\) unique normal extension, AM--GM equivalence of the Alberti forms, N5 Prop.~2.3's Bregman/DK derivation (numerics to \(1.4\times10^{-7}\)).
\item \textbf{[MINOR]} \emph{N5 \S3.1 Remark ``A trap''; Table~1 caption; \S6 text; scripts} \hfill\textsc{V4}\\ Negative eigenvalues ``\(\sim10^{-11}\) for \(\kappa_{\max}\gtrsim24\)'' not reproduced by V4 (min \(+2.2\times10^{-13}\)); caption tail \(8.78\times10^{-7}\) should be \(8.732\times10^{-7}\); the ``\(O(h)\) error of 0.2\%'' is \(+1.0\%/-7.2\%\) depending on \(n_{\mathrm{sub}}\); the M\"obius check is an identity by construction (\(\widehat D_{(2,1,0.1)}=\overline{\widehat D_{(1,2,0.2)}}\) to \(4.6\times10^{-15}\)) and carries no information; stale docstring in fidelity\_flint.py; test\_pipeline.py/test\_precision.py broken (removed \texttt{h=} argument); results\_final.txt not reproducible by the committed script (7 rows \(S=\)nan). Nine graded methodological weaknesses in \S7 of rigor/check\_note5\_numerics.tex.
\item \textbf{[MINOR]} \emph{N1 eq.\ (subnet-bound); N2 eqs.\ (rel-1), (petz-B), (norm-general), \S8} \hfill\textsc{V1}\\ (subnet-bound) is in the proof of LX Thm~4.1(6), and \(N\) is undefined; (rel-1) needs a support hypothesis; (petz-B) leaves the reference state unnamed; (norm-general) is exactly FHSW Thm~2 (uncited); \S8's ``geometrically different'' is thin. Six source typos (LX eq.~(8), Thm~4.2(1) sign, \S4.2.1 cross ratio, Thm~4.4 missing \(\ln\); FHSW eq.~(40)) documented; none affect the notes.
\item \textbf{[MINOR]} \emph{N5 \S4.2 text after Prop.~4.1; Table~1 caption} \hfill\textsc{V3}\\ Quoted tail ratios at \(\kappa=40,60\) are \(1.024,\,1.010\) (not \(1.06,1.06\)) and the \(s^2\)-scaling ratio is \(4.000\) (not \(3.97\)); caption \(8.78\to8.73\times10^{-7}\).
\item \textbf{[MINOR]} \emph{N6 \S6 items N3, N4; N3 \S6; N2 \S8, N3 \S5} \hfill\textsc{R4}\\ \(L_A,L_B\le8\) is VWZ footnote~12, p.~47 (four-interval section), not App.~A.5; N4 must use the two-chirality formula; \(c_E=3RS'\) is the modern normalization while CH04 (23) is \(C(r)=rG'(r)\); CTT (7.15) needs \(H_\gamma=-\log\rho_\gamma-S_\gamma\). In VWZ itself, observation~3 (p.~11) is consistent only on \(\lambda\ge0\), and (3.23) is Pinsker in bits.
\item \textbf{[MINOR]} \emph{Note 5 \S6 (numerics vs.\ bound)} \hfill\textsc{P2}\\ The measured \(\Phi/\zeta^2=0.008376\) at \(\zeta=1/120\) lies \(0.8\%\) below the rigorous \(\liminf\) bound \(0.0084434\); consistent, since the tabulated values are sub-compression lower bounds with the quoted \(-0.0082\zeta\) deficit.
\item \textbf{[MINOR]} \emph{N7 Lemma~3.2 hypotheses; N7 Lemma~3.1 item (1) (``\(c\) real or complex'')} \hfill\textsc{P1}\\ Three simplifications. (i) ``\(\eta\) in the form domain of \(\log\Delta/(\Delta-1)\)'' and ``\(A\) affiliated with a von Neumann algebra containing \(\mathcal M\)'' are not needed for the Bures half; only \(\Omega\in D(A)\), \(A=A^*\) on \(\mathcal H\). (ii) In \(g_B=4\inf\{\|A\Omega-B'\Omega-c\Omega\|^2\}\) the constant \(c\) is redundant (\(c\mathbf 1\in\mathcal M'_{sa}\)); allowing \(c\in\mathbb C\) lowers the infimum by exactly \((\mathrm{Im}\langle\Omega,A\Omega\rangle)^2\), hence is harmless \emph{only because} \(A=A^*\) forces \(\mathrm{Im}\langle\Omega,A\Omega\rangle=0\) (counterexample \(a=i\Omega\): \(G_{\mathbb R}=1\), \(G_{\mathbb C}=0\), \(g_B=4\)). (iii) The natural cone does \emph{not} compute the fidelity: \(\langle\xi_\rho,\xi_\sigma\rangle=\Tr(\rho^{1/2}\sigma^{1/2})\le\Fid\), so the supremum over \(U(\mathcal M')\) is genuinely needed.
\item \textbf{[MINOR]} \emph{\texttt{kernel\_identity.tex} Rem.~6.2} \hfill\textsc{V8}\\ \(Q\psi_\kappa=q_\kappa\psi_\kappa\) is labelled ``although not needed above''. It is needed: Note~7 never writes \(\psi_\kappa\), its completeness relation admits \(e^{\pm i\kappa\xi/2\pi}\) alike, and the identity is odd under \((\kappa,\kappa')\mapsto(-\kappa,-\kappa')\). The check itself is correct (reproduced). Harmless for Note~7, which uses only \(|D^{(1)}|^2\). Promote to a hypothesis in \S1.
\item \textbf{[MINOR]} \emph{\texttt{kernel\_identity.tex} Def.~1.2 / Note~7 Lemma~5.1(i)} \hfill\textsc{V8}\\ Vanishing of \(D^{(1)}\) on \(D\times D\) is postulated here and asserted in Note~7 in one clause without computation, although it changes (20) if false. Now proved (V8 Lemma~2.3): the \(\tfrac12\delta\) parts cancel and \([\widetilde Q,L]\equiv0\) for \(\xi,\xi'>0\), since \(\cosh\xi-\cosh\xi'\) and \(\sinh\xi+\sinh\xi'\) share the factor \(\sinh\frac{\xi+\xi'}2\).
\item \textbf{[MINOR]} \emph{\texttt{kernel\_identity.tex} Rem.~6.4--6.5 and closing verdict} \hfill\textsc{V8}\\ Both defects of Note~7 listed there have since been repaired (errata log, ``sign in Weights paragraph''): Note~7 now carries the minus sign in \(q_\kappa-q_{\kappa'}\) and the corrected contour sentence. The remarks and the closing \texttt{verdictgap} are stale and imply Note~7 still contains a sign error; mark them applied.
\item \textbf{[MINOR]} \emph{\texttt{kernel\_identity.tex} \S5 (cited results)} \hfill\textsc{V8}\\ All four cited boxes omit the mandatory screenshot, justified by the false claim that \texttt{rigor/snap.py} ``is not present in the repository''; it is present. The real obstacle (DLMF is a web page, not a PDF in \texttt{refs/}) is surmountable via \texttt{refs/getref.sh}. No mathematical risk: standard identities, second printed locators given, all re-derived here.
\item \textbf{[MINOR]} \emph{rigor/quadratic\_limit.tex Lem.~2.3(d) \(\langle2\rangle1\); Prop.~7.1 \(\langle1\rangle2\)} \hfill\textsc{V7}\\ Spurious factor \(2\) in \(\lVert D^{(1)}\rVert_2^2=2\iint|D^{(1)}|^2\): the modular coordinate maps all of \(I\) onto \(\mathbb R\), so \(\mathcal F\) is one unitary and there is no doubling. Harmless, only finiteness is used. Separately, Prop.~7.1 \(\langle1\rangle2\) states \(\le\) where equality holds, the projection onto \(V_\Lambda\) being multiplication by \(\mathbf 1_{B_\Lambda}\).
\item \textbf{[MINOR]} \emph{rigor/quadratic\_limit.tex title; Prop.~3.3 parenthesis} \hfill\textsc{V7}\\ The title advertises ``a rigorous upper bound \(=2\log2\cdot g_B/8\)'', but that bound is conditional on Lemma~6.7, which the document itself marks GAP; only \(\Phi\le\tau/(1-\tau)\) is rigorous. The abstract and body are scrupulous. Also, the equality half of Prop.~3.3 (Note~5 Prop.~2.3) is used in Problem~6.8 but neither proved nor flagged as imported.
\item \textbf{[MINOR]} \emph{rigor/quadratic\_limit.tex \S8 (numerics vs.\ bound), \S7.3} \hfill\textsc{V7}\\ The measured \(\Phi/\zeta^2=0.008376\) at \(\zeta=1/120\), below \(f_2=0.0084434\), is consistent: a liminf constrains no finite \(\zeta\), and the tabulated values are sub-compression lower bounds. But the offered explanation is circular (\(-0.0082\zeta\) is the fit of that same deficit), and \S7.3's substitution \(\kappa_*\approx2\log(1/\zeta)\) is off by \(4.7\) at \(\zeta=1/120\).
\item \textbf{[MINOR]} \emph{N6 eq.\ (energy); N7 after eq.\ (6), Lemma~5.1(i), Cor.\ 6.1; N6 Conjecture~2.1, Observation~1} \hfill\textsc{V5}\\ N6 eq.\ (5) is \(E(\varphi^{-1})\), not \(E(\varphi)\): the Schwarzian that appears is \(\{\varphi^{-1};\theta\}\), and the two differ (2.006827 vs 2.325495 for \(\varphi=\theta+0.21\sin2\theta\)); the relative sign IS correct, though, and the combination vanishes on \(\mathrm{PSL}(2,\mathbb R)\) to \(3\times10^{-21}\) (verified on \(\tan\frac\varphi2=\lambda\tan\frac\theta2\)), so the note's hedge can be dropped. N7: \(\langle\widetilde T\rangle=0\) in the note's own convention \(\widetilde T=\beta^2T\), so the subtraction \(\widetilde T_0\) is vacuous; Lemma~5.1(i)'s ``on \(A\times A\) and \(D\times D\) both sides vanish'' needs the M\"obius-invariance argument (\(h_s(D)\subset D\), \(P_+\) M\"obius invariant); Cor.\ 6.1's \(-\log\Fid\simeq\frac{12}{\pi^2c}I^2\) is the chiral law (2D: \(\frac{12}{\pi^2(c+\bar c)}I_{\mathrm{tot}}^2\)). N6 Conjecture~2.1 lacks the two-branch qualification; Observation~1's ``So \(U(\widetilde k)\) exists'' still reads as a deduction from what the preceding sentence declares missing.
\item \textbf{[MINOR]} \emph{N7 \S3.2, the parenthetical after eq.\ (mu) (``the KMS asymmetry \dots would fix it only for \(T(g)\in\cM\)'')} \hfill\textsc{P4}\\ The sign of the modular dilation is in fact forced without Bisognano--Wichmann: the opposite choice replaces \(\widehat W(k)\) by \(\widehat W(-k)\), multiplying both integrands by \(e^{2\pi k}\), so \(g_{\mathrm B}\) and \(g_{\mathrm{KM}}\) both diverge. Note~7's hedge understates what is available; finiteness alone selects the sign.
\item \textbf{[MINOR]} \emph{N5, paragraph after Cor.~5.1: ``the apparent exponents \(\dots\) come out as \(2.04,2.04,1.96,1.31\)''} \hfill\textsc{PL}\\ The lattice gives \(1.99,1.91,1.64,1.15\) at \(\lambda=0,0.5,1,1.5\) (fit over \(\eta_V\le0.1\), \(n\le72\)), i.e.\ VWZ's \(2.0,1.9,1.7,1.2\) almost exactly. The note's \(2.04,2.04,1.96,1.31\) come from a \(\Phi\) extrapolated above the tabulated \(\zeta\le0.53\); they overstate \(p\) at \(\lambda=0.5,1\). Recompute, or drop the numbers.
\item \textbf{[MINOR]} \emph{phase1\_brief.md step 1, kernel normalisation} \hfill\textsc{PN}\\ The correct relation is \(\|P_-W_\phi P_+\|_2^2=\tfrac12\|K\|_2^2\), not \(\|K\|_2^2\): \(K=W^*P_-W-P_-\) is a difference of projections, so \(\operatorname{Tr}K^2=\|P_-WP_+\|_2^2+\|P_+WP_-\|_2^2\). Verified against a direct Fourier-basis matrix computation of \(P_-W_\phi P_+\) and against \(\epsilon^2/(48\pi)\) for \(\phi=x+\epsilon e^{-x^2/2}\).
\item \textbf{[MINOR]} \emph{rigor/uhlmann\_upper\_bound.tex \S7.2--7.3; numerics/p1\_bogoliubov.out (PN)} \hfill\textsc{PA}\\ Confirms and locates PN's factor-8 warning. In the half-line chart the admissible circle field only needs \(\chi(x)\to\alpha\) finite at \(+\infty\) (displacement tends to a constant, not to \(0\)); \(\alpha=0\) is a proper subclass. PN's restricted class corresponds, after the M\"obius map, to a second-order vanishing at an \emph{interior} point of \(I'\). The \(\mathcal A'\)-route proved here is unaffected.
\item \textbf{[MINOR]} \emph{rigor/uhlmann\_upper\_bound.tex \S10 (imports)} \hfill\textsc{PA}\\ Three imports could not be screenshotted: Shale--Stinespring 1965 Thm~3 (and Araki RIMS 6, Lem.~6.3/Thm~6.5), Ruijsenaars 1977, and Foit RIMS 19 (1983) Thm~1 (twisted duality). The determinant formula is reproved, so only implementability and twisted duality remain as unverified citations; both are standard for the free fermion net.
\item \textbf{[MINOR]} \emph{rigor/uhlmann\_upper\_bound.tex Prop.~5.4, Rem.~9.3} \hfill\textsc{PA}\\ The error term in \(\lVert V_\zeta\rVert_2\le\zeta\lVert P_-KP_+\rVert_2+\zeta\epsilon(\zeta)+C_4\zeta^2\) contains one qualitative factor, \(\epsilon(\zeta)=\sup_{\sigma\le\zeta}\lVert \widetilde W_\sigma^*G^{(1)}\widetilde W_\sigma-G^{(1)}\rVert_2\to0\) ($L^2$-continuity of composition with the flow). So the result is \(\Phi=f_2\zeta^2+o(\zeta^2)\); a certified remainder still needs the tail analysis (O4).
\item \textbf{[MINOR]} \emph{rigor/uhlmann\_upper\_bound.tex \S6 and Thm.~7.2} \hfill\textsc{PR}\\ Thm.~8.1 uses only Def.~2.1 plus the operator class \(\mathcal A'\); the repaired circle class \(\mathcal X\) and eq.~(32) are never used, and \S6 does not say so (the reason it is safe: the distance depends on \(a\) only through \(\varphi\)). Also PA's \(\Omega(\ell)\) has the opposite linear sign from quadratic\_limit Def.~4.1, harmlessly.
\item \textbf{[MINOR]} \emph{numerics/results\_hp14.txt and numerics/README.md} \hfill\textsc{PT}\\ Confirmed and quantified: when the spectral window contains all box modes the tail must use \(\kappa_{\rm eff}=\kappa_{\max}^{\rm eff}=29.6\) and the box coefficient \(0.513/8=0.0641\) (referee\_quadratic\_limit \S7), not \(\kappa_c=32.2\) with \(1/15\). The printed ``+tail'' column is \(1.2\times10^{-3}\) relative too small; corrected \(\Phi(1/15)=3.5191\times10^{-5}\).
\item \textbf{[MINOR]} \emph{numerics/run\_largezeta\_cert.py (high-precision eigensolvers)} \hfill\textsc{PZ}\\ Arb's \texttt{acb\_mat.eig} silently returns \texttt{nan} once the Hermitian spectral range exceeds \(\sim10^{-20}\) (any precision): it fails on \(WW^{H}\) in the fidelity step and on the box symbol \(Q\) at \(\kappa_{\max}=60\). mpmath \texttt{eighe} is exact there and fast enough; boxes with \(\kappa_{\max}>45\) are currently inaccessible.
\item \textbf{[MINOR]} \emph{refs/P1\_CarpiWeiner\_math-0407190.pdf Lemma~4.6 (as used in their Prop.~5.4)} \hfill\textsc{QB}\\ ``Support contained in \(I\)'' must mean \emph{compactly} contained: the mollified approximants live in an \(\epsilon\)-enlargement. In CW05's own application (and in ours) the field vanishes on an open arc but its support is the closed complement, so one gets affiliation to \(\mathcal A(I_\epsilon)\) and must intersect using the semicontinuity axiom. Repairable in two lines.
\item \textbf{[MINOR]} \emph{strategy\_open\_problems.tex \S3 Route A; open\_problems\_plan.tex \S O5} \hfill\textsc{QB}\\ Both say Route~A is closed only for the free fermion / Fock representations with integer \(c\), the obstruction being CDIT's \(s>3\). That restriction concerns the \emph{group-level} extension to \(D^s(S^1)\); for the one-parameter flow \(\widetilde k_s=\mathrm{Exp}(s\,g_{\rm tot})\) Stone's theorem supplies the unitary, so Route~A closes for all \(c\).
\item \textbf{[MINOR]} \emph{rigor/lemma\_second\_variation.tex Thm.~``Upper bound'' vs (H3)} \hfill\textsc{QA}\\ The imported theorem assumes the first-order expansion for all \(y\in\mathcal M\); (H3) gives it only on a dense subalgebra. Closing this needs \(\mathcal M_0\) to be a \(*\)-subalgebra (so \(K^2\in\mathcal M_0\)) that is \(*\)-strongly, not merely strongly, dense, so that \(K_\alpha\Omega\to K\Omega\) for self-adjoint \(K\).
\item \textbf{[MINOR]} \emph{universality\_normalization.tex eq.~(3.27), numerical evaluation} \hfill\textsc{QK}\\ A square-box quadrature \(\lvert\kappa\rvert\le L\) of that double integral converges to \(0.147\) at \(L=30\) and \(0.161\) at \(L=200\), apparently missing \(c/6\) by \(3\)--\(12\%\). This is a box artefact (the \(\sigma\)-range is cut to \(2L-\lvert u\rvert\)); future checks should use the \((\sigma,u)\) chart.
\item \textbf{[MINOR]} \emph{phase2\_brief.md track QB, ``Trotter from \(g_A\), \(g_D\) the restrictions''} \hfill\textsc{QR2}\\ That split is unusable: \(g_{\rm tot}(L)=-L\neq0\), so the truncation to \([0,L]\) jumps at \(x=L\) and its Fourier coefficients decay like \(1/n\); the \(3/2\)-norm of both pieces is infinite. QB instead splits \(g_{\rm tot}=(g_{\rm tot}-m)+m\) with \(m=-x^2/L\) globally smooth and \(g_{\rm tot}-m\) vanishing on \(D\); both summands stay in the Carpi--Weiner class.
\item \textbf{[MINOR]} \emph{rigor/implementation\_C11.tex Thm.~4.4 \(\langle1\rangle6\); Thm.~4.5(ii); Prop.~5.1(ii),(iii)} \hfill\textsc{QR2}\\ Three repairs. (i) ``For \(m\) large the union \(\bigcup_{k\ge m}\varphi_k J\) is an interval'' is asserted; a proof needs the middle-third argument plus Hausdorff convergence. (ii) \(\varphi(X)=i\omega([T(g_{\rm tot}),X])\) is formal, since \(X\Omega\in D(T(g_{\rm tot}))\) is unknown; the sesquilinear form is the rigorous object. (iii) \(\lVert T(2\cos n\theta)\Omega\rVert^2=\frac{c}{12}n(n^2-1)\), not \(\frac{c}{3}\); and \([L_0,T(g)]=+iT(g')\).
\item \textbf{[MINOR]} \emph{numerics/README.md, spectral-window tail recipe} \hfill\textsc{QN}\\ The UV window tails of \(-\log F\) and of \(D\) obey \emph{different} laws: fits over \(\kappa_c\in[2.3,5.2]\) give \(f_2^{(\infty)}-f_2(\kappa_c)\propto\kappa_c^{-p}\) with \(p\approx2.3\), but \(p\approx1.15\) for \(s_2\). Expected: the discarded per-mode weight decays like \(\kappa^{-3}\) and the Kubo--Mori metric carries an extra modular-energy factor \(\kappa\). Using \(p=2\) for \(s_2\) undershoots by \(6\%\).
\item \textbf{[MINOR]} \emph{numerics/boson/boson\_recovery.py \S2 (finite-\(s\) scan)} \hfill\textsc{QN}\\ Double precision suffices for \(-\log F\) down to \(s\approx0.05\) (\(\Phi\sim10^{-6}\), \(2\)--\(5\%\) scatter) but not for \(D(\omega_s\Vert\omega)\), a difference of two \(O(1)\) entropies, which turns negative below \(s\approx0.1\). The second-order coefficients were therefore computed exactly at \(s=0\) from the Bures and Kubo--Mori metrics instead of extrapolated.
\item \textbf{[MINOR]} \emph{rigor/universality\_theorem\_B.tex, Prop. 3.2 \(\langle1\rangle4\)} \hfill\textsc{QR1}\\ ``\(\{K\Omega-\omega(K)\Omega\}\) is dense in \(H_{\mathcal M}\)'' is false: every such vector is real-orthogonal to \(\Omega\in H_{\mathcal M}\), so the closure has real codimension one. Conclusion unaffected, since \(\mathrm{Re}\langle-ia,\Omega\rangle=0\) makes \(\Theta(k+t\Omega)\le\Theta(k)\).
\item \textbf{[MINOR]} \emph{rigor/universality\_theorem\_B.tex, hypotheses of Thm 5.6 and (H3)} \hfill\textsc{QR1}\\ Theorem 5.6 lists (HD), (RS), but its proof uses only (BW), the (S2) two-point function and \(a\in\mathcal H_T\); trimming matters because the fermion violates (HD). Also (H2) implies (H3) with \(\mathcal M_0=\mathcal M\), so (H3) is redundant in Theorem 7.1. Haag duality is cited three times where locality suffices.
\item \textbf{[MINOR]} \emph{numerics/p1\_bogoliubov.out (3), the ``\(-0.69\,\zeta\)'' slope of the Uhlmann bound} \hfill\textsc{RR}\\ Not a property of the bound: that family rebuilds the extension at each \(s\) and is not a one-parameter group, so its bound may carry a linear term (measured slope \(-0.73,-0.71,-0.58,-0.35\)). Every group family gives an even bound. Consistent with \(c_3=-0.99\) since \(\tfrac12E/\Phi=1+0.27\zeta\).
\item \textbf{[MINOR]} \emph{rigor/rate\_of\_remainder.tex Lem.~2.7 (hypothesis (F4))} \hfill\textsc{RR}\\ The Lipschitz bound \(\epsilon(s)\le s\lVert\kappa^{(2)}\rVert_2/2\pi\) needs \(\chi\chi''\) Lipschitz. A jump of \(\chi''\) at a point where \(\chi\ne0\) makes \(\kappa^{(2)}\notin L^2\) logarithmically (hence only \(O(s\log\frac1s)\), as PA's verdict anticipated). The kink at \(x=0\) is harmless: \(\chi,\chi'\) vanish there to second order.
\item \textbf{[MINOR]} \emph{tail\_bound.tex, \S9 ``Where to attack next''} \hfill\textsc{RC}\\ The requested bound $\lVert\sqrt{\widetilde Q}-\sqrt{X_{11}\oplus X_{22}}\rVert_2^2\le C\lVert X_{12}\rVert_2^2$ is false: two-by-two positive contractions give ratio $3/\varepsilon$, $\varepsilon=\lambda_{\min}(X_{11})$. It is true with the Daleckii--Krein weight $(\sqrt{\beta}+\sqrt a)^{-2}$, where it is an identity, and positivity then bounds the weighted norm by $\tau$.
\item \textbf{[MINOR]} \emph{numerics/compression\_box.py, function \texttt{Q\_box}} \hfill\textsc{RC}\\ The $4\times10^5$-node trapezoid rule for the box symbol is accurate to only $4.7\times10^{-11}$ per entry (Euler--Maclaurin endpoint term); all off-diagonal entries lie outside their rigorous enclosures. Ball arithmetic on the $2N$ one-dimensional integrals gives $5.6\times10^{-18}$ instead.
\item \textbf{[MINOR]} \emph{certified\_numerics.tex, \S1, third and fourth bullets} \hfill\textsc{RCr}\\ The scope section announces results the body does not contain: the unconditional bound with $C_\star=6$ and $\Phi-\Phi_W\le0.85\zeta^2\kappa_c^{-2}$ (the body has $c_0\approx35$ instead); ``verified eigenvalue enclosures (Rayleigh--Ritz + residual)'' (the scripts call numpy eigh); ``the fidelity determinant in ball arithmetic'' (never done); ``certified values of $\tau$'' (not certified).
\item \textbf{[MINOR]} \emph{certified\_numerics.tex, \S5, last paragraph} \hfill\textsc{RCr}\\ ``An arithmetic margin of $8\,\%$ rather than $3\,\%$'' does not follow: with the note's own $\eta=N^{3/2}\cdot4.7\times10^{-11}=4.1\times10^{-8}$, PT's (8.1) gives $3.3\,\%$ at $\zeta=8/15$, $21.7\,\%$ at $1/15$ and $204\,\%$ at $1/120$, the relative margin growing like $1/\zeta$. Also the endpoint Taylor error improves as $d^3$, not $d^2$.
\item \textbf{[MINOR]} \emph{certified\_numerics.tex, Table 3 caption and \S8 verdict} \hfill\textsc{RCr}\\ The column ``$\Phi$ incl.\ box cutoff'' adds PT's $0.0641\zeta^2\kappa_{\max}^{-2}$, which is a one-term asymptotics of $f_2$, not a proved bound on $\Phi_E-\Phi_{P_{\mathrm{sub}}}$ (tail\_bound.tex \S7b says so explicitly); it must be labelled extrapolated, not certified. Its widths are $1.5$--$2.4\,\%$, not $1.7$--$2.4\,\%$ ($\zeta=1/30$ gives $1.54\,\%$).
\item \textbf{[MINOR]} \emph{rate\_of\_remainder.tex, all cross-references to uhlmann\_upper\_bound.tex} \hfill\textsc{RRr}\\ Citations are shifted by one section: Lemma 2.3$\to$2.4, Prop.\ 2.4$\to$2.5, Prop.\ 2.7 (Uhlmann)$\to$3.4, Thm.\ 3.1 (overlap)$\to$4.1, Lemma 4.1$\to$5.1, Lemma 4.2$\to$5.2, Def.\ 4.4$\to$5.3, Prop.\ 4.5$\to$5.4, Lemma 6.3$\to$7.5/7.6, Thm.\ 7.1 and 7.2$\to$8.1 and 8.2. The content used is always the intended result.
\item \textbf{[MINOR]} \emph{rate\_of\_remainder.tex, abstract} \hfill\textsc{RRr}\\ Two conflicts with the body: the abstract gives $c_3=-0.96(3)$ where eq.\ (16) and \S5.2 give $-0.99(1)$; and it says $C(\Lambda)$ ``grows polynomially, not exponentially, so the optimised lower rate is a power'', where Prop.\ 4.2 and the summary table say bounded and linear.
\item \textbf{[MINOR]} \emph{rate\_of\_remainder.tex, \S3.2 and the sentence after (F2)} \hfill\textsc{RRr}\\ Sign typo: $c_4^{(P)}:=(c_3^{(P)}-\mathrm{slope})/\zeta$ should be $(\mathrm{slope}-c_3^{(P)})/\zeta$; the table's $+0.874$ is right, the definition is not. And ``the flow satisfies $\widetilde k_s(x)=x/(1+sx)$ on $[0,1]$ simultaneously for all $s$'' holds only for $s\ge0$, as PA Lemma 2.2 (E2) states.
\item \textbf{[MINOR]} \emph{certified\_numerics.tex, Table~3 (first version)} \hfill\textsc{RC}\\ Self-correction after the referee report: the $\Phi_{\rm off}$ column was double-precision noise (the control identity $-\log\lvert\det M_A\rvert=\Phi_W+\Phi_{22}$ failed by $180$ times the quantity). Recomputed in mpmath at $60$ digits with the identity checked per run; Table~3 is superseded by \S9.
\item \textbf{[MINOR]} \emph{rigor/universality\_theorem\_B.tex \S5.2, step \(\langle1\rangle2\)} \hfill\textsc{PP2 (orchestrator applied)}\\ Inverted factor: with \(V[f]=(48\pi^2/c)^{1/2}T(f)\Omega\) unitary, \(V^{-1}a=(c/48\pi^2)^{1/2}[g_{\rm tot}]\), not \((48\pi^2/c)^{1/2}[g_{\rm tot}]\). Fixed in the source and in paper 2; no downstream effect (the squared norm is used).
\item \textbf{[MINOR]} \emph{fidelity\_recovered\_quasifree.tex Prop.~2.3 and Cor.~4.2; rigor/optimality\_second\_order.tex \S3.1} \hfill\textsc{S8}\\ The quantum-Fisher form is the second-order term of \(-\log F\) only for defects small compared with \(e^{-\kappa_{\max}}\). At \(n=13\) sites, \(z=0.15\), the uncapped form gives \(3.6\times10^{-2}\) against the exact \(1.3\times10^{-3}\). Minimising it on a modular window produces a minimiser \(880\) times worse than Petz; capping the weights is the fix.
\item \textbf{[MINOR]} \emph{Note~4 geometric compression, lattice transcription; numerics/optimality/compression.py} \hfill\textsc{S8}\\ Discretising the Mobius squeeze \(h(x)=x/(1+\lambda x/L)\) in the site basis does not give a lattice compression: singular values fall to \(0.47\)--\(0.62\) (modes pushed through the cutoff) and \(-\log F=0.66\)--\(2.2\), four orders worse than Petz. The lattice optimum realises the compression in the infrared only.
\item \textbf{[MINOR]} \emph{rigor/phase4\_brief.md, ``The route (Uhlmann), stated correctly''} \hfill\textsc{E1}\\ The brief says an admissible extension \(\widetilde W\) ``preserves \(L^2(I^c)\) as well'' and writes \(Z=Z_I\oplus Z_e\) block diagonally. Both are false: \(W_I\) maps \(L^2(I)\) isometrically \emph{onto the proper subspace} \(L^2(AB)\), so \(\widetilde W L^2(I^c)\supsetneq L^2(I^c)\) and no admissible generator is block diagonal. Correct bookkeeping (rigor/exact\_fidelity\_upper\_half.tex Lemma~4.1): prescribing \(\widetilde W|_{L^2(I)}\) fixes the whole column \(ZE_I\), hence also \(E_IZE_{I^c}=-(E_{I^c}ZE_I)^{*}\), and the only free block is \(E_{I^c}ZE_{I^c}\) --- which is exactly the \(\mathrm{Ext}\) parametrisation, so the conclusions stand. Reading \(Z_I\) as the truncation \(E_ID_\chi E_I\) instead is wrong: E2 measures a visible defect \(\lVert\Pi P^{\perp}(E_ID_\chi E_I)P\rVert^2/\lVert u\rVert^2=2608,6188,11052,25290\) at \(L=128,192,256,384\).
\item \textbf{[MINOR]} \emph{rigor/exact\_fidelity\_upper\_half.tex verdict box after Thm.~5.3; Lemma~4.1(c); Thm.~7.1 statement; \S8 item~3} \hfill\textsc{RE1}\\ Four local defects, none load-bearing. (1) The verdict box claims compactness is necessary and that ``\(\Gamma'\pi\Gamma'=\pi\) has nonzero bounded solutions (build \(\pi\) from an approximate eigenvector sequence)'' whenever \(\pm1\in\mathrm{spec}\,\Gamma'\). False: under the theorem's own hypothesis \(\Gamma'^n\to0\) strongly, so \(\langle u,\pi v\rangle=\langle\Gamma'^nu,\pi\Gamma'^nv\rangle\to0\) kills every bounded solution, and weak-\(*\) cluster points of \(\lvert\xi_n\rangle\langle\xi_n\rvert\) for approximate eigenvectors are \(0\). The correct reason for claiming only a closure is E2 \S2.3 (the gains \(\sqrt{d^c_{ij}/2}\) have a lower quartile falling from \(10^{-2.5}\) to \(10^{-7.5}\) between \(L=96\) and \(384\)), i.e.\ \(Z_e\mapsto P^{\perp}Z_eP\) has unbounded inverse. (2) Lemma~4.1(c) (``the only free block is \(E_{I^c}ZE_{I^c}\), an arbitrary anti-self-adjoint operator'') is false for a one-parameter group with \(W_I(\tau)\) prescribed for all \(\tau\): \(Z^2E_I=\ddot W_I(0)E_I\) contains \(\widehat Z_e(E_{I^c}ZE_I)\), so \(\widehat Z_e\) is determined on \(\mathrm{ran}(E_{I^c}ZE_I)\). The statement used downstream is the differentiated torsor statement of Lemma~1.2 (right multiplication by \(e^{sZ_e}\)), which is true. (3) Thm.~7.1 writes \(\delta(s)=E_I[P,-s(D_\chi+\lambda\mathcal G)]E_I+o(s)\), the opposite sign to S10 eq.~(24) and to the direct computation \(\delta=E_I(P-\widetilde W^{*}P\widetilde W)E_I=-sE_I[P,L_0]E_I+O(s^2)\); immaterial since \(g_Q\) is even. (4) \S8 item~3's claim that the \(o(1)\) of Thm.~6.2 is uniform over exterior generators with bounded \(\lVert Z_e\rVert\) and \(\lVert[P,Z_e]\rVert_2\) is unjustified (the \(\mathfrak S_2\)-continuity of \(\sigma\mapsto V_\sigma TV_\sigma^{*}\) is not uniform on norm-bounded sets) and is not needed, the limsup being taken before the infimum.
\item \textbf{[MINOR]} \emph{rigor/exact\_fidelity\_upper\_half.tex, citations to uhlmann\_upper\_bound.tex (PA), rate\_of\_remainder.tex (RR), sdp\_dual\_certificate.tex (S10), optimality\_all\_channels.tex} \hfill\textsc{RE1}\\ Fifteen locators point at the wrong statement (content attributed is correct in every case; only the numbers are wrong, apparently written with PA's ``Step \(k\)'' labels instead of section numbers). PA Cor.~1.9 \(\to\) 2.7; PA \S2 (Alberti--Uhlmann citedbox) \(\to\) \S3; PA Thm.~3.1 \(\to\) 4.1; PA Lemma~2.4 \(\to\) 3.2 (2.4 is the kernel-of-the-defect lemma); PA Lemma~5.5 \(\to\) 7.5; PA \S5.3 eq.~(5.5) \(\to\) \S7.3 eq.~(35); PA Prop.~5.9 \(\to\) 7.6; RR Lemma~2.2 \(\to\) Lemma~6.1 \(\langle1\rangle2\) (2.2 is the autocorrelation representation); RR Thm.~2.5 \(\langle1\rangle4\) \(\to\) Thm.~2.3 \(\langle1\rangle4\); RR Thm.~2.6 \(\langle1\rangle2\) \(\to\) Thm.~2.5 \(\langle1\rangle2\); S10 eq.~(2.1) \(\to\) eq.~(1); S10 \S9 \(\to\) \S7 (and Cor.~11.2); optimality\_all\_channels \S5.2--5.3 \(\to\) \S5.1--5.2; its Prop.~2.4 \(\to\) Prop.~2.5 with Thm.~4.1(2),(3); the sigil ``S12'' is used but never defined. Also: ``Cousin--Sylvester (CS) decomposition'' should be ``cosine--sine''; the abstract drops the finite-rank restriction on \(\mathcal N\) that Definition~3.1 imposes and advertises a Halmos normal-form proof where Thm.~5.3 gives the invariant computation; and PA Lemma~7.5 is written in PA's half-line configuration, so applying it to a bounded interval (Thm.~5.4) needs one sentence on twisted duality for a disconnected complement. Note also that the \emph{value} \(\frac12\lVert u\rVert_2^2=f_2\) used in Thm.~7.2 reaches back to PA Prop.~7.6, hence to twisted duality, contrary to the impression that only the second proof of (Q) uses it (the dependence is removable: PA Prop.~7.6 \(\langle1\rangle3\)--\(\langle1\rangle6\) uses only the \(X=I\) half of PA Lemma~7.5).
\end{enumerate}
\section{Verified items and independent confirmations}
\begin{enumerate}[leftmargin=1.6em,itemsep=6pt]
\item \textbf{[OK]} \emph{N7 Lemma~5.1(ii) (kernel identity)} \hfill\textsc{P3}\\ Fully proved analytically (rigor/kernel\_identity.tex): with \(a=-\xi\), \(b=\xi'\) the kernel is \(\sinh(a/2)\sinh(b/2)/\sinh^2((a+b)/2)\in(0,\tfrac14]\), bounded by \(e^{-(a+b)/2}\); the series \(\sinh^{-2}s=4\sum n e^{-2ns}\) and Tonelli give \(J=\sum_n n/[(n-A_+)(n-A_-)(n-B_+)(n-B_-)]\), summed by Gauss (DLMF 5.7.6) to \(J=-\frac{\sigma}{k(k^2+1)}[\psi(\tfrac12-i\alpha)-\psi(\tfrac12-i\beta)]+\frac{2}{k^2+1}\); \(\operatorname{Im}\psi(\tfrac12-ix)=-\frac\pi2\tanh\pi x\) yields exactly \(i(q_\kappa-q_{\kappa'})\mathfrak g(\kappa,\kappa')\). By-products: \(D^{(1)}(\kappa,\kappa')=(q_\kappa-q_{\kappa'})v/[u(u^2+4\pi^2)]\), \(D^{(1)}(\kappa,\kappa)=-\kappa/(8\pi^2\cosh^2(\kappa/2))\). Verified by rigor/p3\_check.py (30 digits, 14 check groups).
\item \textbf{[OK]} \emph{N1 main formula; N2, N4 eq.\ (LX-MI); N5/N7 normalizations} \hfill\textsc{R1}\\ No factor-of-two error. Longo--Xu's \(\mathcal A_r\) is \(U(r)_1\), i.e.\ \(r\) complex chiral fermions with \(c=r\) (LX \S4.4.1, Rmk~4.3); for a chiral net the two-interval mutual information is \(\frac c6\log\frac1\eta\), half the non-chiral \(\frac c3\log\frac1\eta\) (the halving is Casini--Huerta eq.~(26), \(\rho=\rho(u_+)\otimes\rho(u_-)\)). LX Thm~3.18 re-derived from Def~3.17 in Lamport form (rigor/cited\_R1.tex) and the non-chiral \(\frac13\) confirmed numerically to 6 digits (numerics/R1\_chiral\_vs\_dirac\_MI.py). Also verified in full: \(J_{\mathcal N}=U(2\lambda)J_{\mathcal M}\) (sign-robust; = Wiesbrock CMP~157 relation~2), \(\psi_\kappa\), \(q_\kappa=1/(1+e^\kappa)\) from the note's own Hardy sign, N7's \(W(u)\) and \(J:x\mapsto2L-x\), N4 eqs.\ (LX-MI), (LX-collar-MI), (delta-eps), (CMI-LX).
\item \textbf{[OK]} \emph{N7 Lemma~5.1(ii) and eq.\ (V)} \hfill\textsc{V6}\\ Independently reproduced numerically (4 points, ratio \(1.00000000\); eq.\ (V) to \(5\times10^{-16}\)).
\item \textbf{[OK]} \emph{N5 \S6 numerics, scripts} \hfill\textsc{V4}\\ Independent recomputation reproduces \(\Phi(1/15)\) exactly (raw \(3.432123265\times10^{-5}\), corrected \(3.519443495\times10^{-5}\)) by two algebraically distinct routes agreeing to 40 digits, both validated against brute-force Fock traces; defect\_matrix.py verified line by line (log-panel Jacobian, weights, phases, \(1/\Lambda\), Hermitian traceless assembly) and matches the closed form of rigor/kernel\_identity.tex to 8 digits; eq.\ (Qbox) re-derived and verified by a Plancherel route (\(2.6\times10^{-9}\)); the spectral sub-compression is a genuine compression (proved); the tail correction \(\zeta^2/(15\kappa_c^2)\) and the \(1/\kappa_c\) extrapolation formula are correctly derived; fidelity\_tables.tex provenance verified. Refits: \(f_2=0.0084472(5)\) with free exponent \(p=1.854(20)\), \(0.0084433\) with \(p=2\) fixed (analytic \(0.0084434\)); \(s_2=0.08433(32)\) with \(p=0.880(30)\), \(0.083246\) with \(p=1\) fixed (analytic \(0.083333\)).
\item \textbf{[OK]} \emph{N1 and N2, all displayed results} \hfill\textsc{V1}\\ No proof-breaking error. Every displayed result has a Lamport verification (rigor/check\_notes\_1\_2.tex, 36 pp.); all closed forms re-derived from LX Def.~3.17 and checked numerically on 2000 random configurations (max error \(2\times10^{-14}\)). The coefficient \(r/6\) is confirmed twice: from \(G((u,v))=\frac r6\ln|v-u|\) and from LX's own closed form on p.~24, \(F(A\bar\cup B_0,A\bar\cup C_0)=-\frac r6\ln|(b_2-a_2)(b_3-a_1)/((b_3-a_2)(b_2-a_1))|\), which is N1's \(-\frac r6\log\eta\).
\item \textbf{[OK]} \emph{N5 Sections 1--4 (strengthened)} \hfill\textsc{V3}\\ Verified: cross-ratio identity and M\"obius reduction; Thm~2.1, Props~2.2--2.3 (algebraic identity, Kubo--Mori/Daleckii--Krein, (qfi-weight), classical limits); Thm~3.1 (with the corrected locator); \((x-y)/\sqrt{\beta\beta'}=2\sinh\frac{\xi-\xi'}2\); the \(t\to0\) limit \(2\Lambda k\); the ``trap'' mechanism; mode normalization; the sign of \(\kappa\) by direct computation (\(q_\kappa=\frac12-\frac12\tanh\frac\kappa2\), no positivity appeal needed); Prop.~4.1 with the EXACT modular-frame kernel \(\widetilde F=\zeta S_pS_{p'}/[\sinh M(\sinh M+2\zeta S_pS_{p'})]\), \(c_2=-\zeta^2/30\) exact in \(\zeta\), the \(O(\kappa^{-5})\) remainder justified, first-order diagonal \(=C\kappa e^{-\kappa}\); the \(\kappa_*\) estimates; monotonicity of \(\Phi\) in \(\zeta\) UPGRADED TO A PROOF (Cor.~5.1 no longer needs numerics for it); the open-problem (i) identity of \S8.
\item \textbf{[OK]} \emph{VWZ conversions (N5 \S5, \S7.1; N7 Cor.~6.1); N3 thermal/RG formulas} \hfill\textsc{R4}\\ Verified with screenshots (rigor/cited\_R4.tex, 38 pp.): root fidelity; \(\lambda=2t\); \(\zeta=2\eta_{\mathrm V}/(1-\eta_{\mathrm V})\); the Dirac two-component factor; \(\beta(\lambda)d\lambda=p(t)dt\); VWZ's \(c\) is the holomorphic central charge (\(c=\bar c\)), universality stated for \(-\log F/c\) and \(D_\alpha/c\); \(d_1=2(0.44)^2/\ln2=0.5586\) (precision \(\pm0.03\) from their Fig.~19); thermal CMI and its \(e^{-2\pi b/\beta}\) asymptotics; \(c_M\ge c_E\). With the corrected two-branch law and \(f_2=c/(12\pi^2)\), VWZ's Fig.~5 exponents \(p\approx2.0,1.9,1.7,1.2\) at \(\lambda=0,0.5,1,1.5\) are reproduced as \(2.04,2.04,1.96,1.31\), and the \(\lambda\)-independent \(\eta\to1\) exponent \(c/9\) follows. Lattice lemma: any state with real reduced density matrices has \(F^{(\lambda)}=F^{(-\lambda)}\) exactly.
\item \textbf{[OK]} \emph{Note 5 Cor.~4.2 / Note 7 Thm~A (lower bound)} \hfill\textsc{P2}\\ Unconditional theorem (\texttt{quadratic\_limit.tex}, Thm~4.2): \(\liminf_{\zeta\to0}\Phi(\zeta)/\zeta^2\ge g_B/8=r/(12\pi^2)\). Proof avoids all limit interchanges: \(\Phi\ge\Phi_P\) for a single finite-rank \(P\) (data processing), real-analyticity of \(\Phi_P\), a Legendre form \(g_B=\sup_X[2\Tr(\dot\rho X)-\Tr(\rho X^2)]\) with quadratic test operators \(X=d\Gamma(\ell)\), and \(\sup_Pg_B^{(P)}=g_B\) by density of finite-rank kernels in \(L^2(N\,d\kappa\,d\kappa')\). Prop.~2.3 of Note~5 is not needed.
\item \textbf{[OK]} \emph{Note 7 Thm~5.2 (integrals)} \hfill\textsc{P2}\\ \(g_B=2r/(3\pi^2)\) now fully proved (Prop.~5.1): the integral \(\int w^2dw/(\cosh w+\cosh a)=2a(a^2+\pi^2)/(3\sinh a)\), previously only checked numerically, is proved via \(1/(\cosh w+\cosh a)=[\varphi(w+a)-\varphi(w-a)]/\sinh a\) (Lemma~5.2); both \(u\)-integrals proved (Lemma~5.3); kernel identity imported from P3.
\item \textbf{[OK]} \emph{Note 5 (new bound)} \hfill\textsc{P2}\\ Unconditional upper bound \(\Phi(\zeta)\le\tau/(1-\tau)\), \(\tau=\frac r{2\pi}\log(1+\zeta)\) (Thm~6.4), via \(F\ge\Tr\sqrt\rho\sqrt\sigma\), Ostrowski--Taussky and Powers--St\o rmer. New closed form for the Hellinger (R\'enyi-\(\frac12\)) second-order coefficient \(h_2=\frac{r\log2}{6\pi^2}=2\log2\cdot f_2\) with weight \(1/W\), \(W=\sin^2(\theta_i+\theta_k)\) (Props.~6.5/6.6).
\item \textbf{[OK]} \emph{N7 Lemma~3.2 (type-III second variation), fidelity half} \hfill\textsc{P1}\\ Proved in full generality in \texttt{rigor/lemma\_second\_variation.tex} (Thm.~7.3): for any von Neumann algebra with cyclic separating \(\Omega\) and any norm-differentiable curve of unit vectors \(\xi_s\) (in particular \(\xi_s=e^{-isA}\Omega\) with \(\Omega\in D(A)\)), \(-\log\Fid=\frac{s^2}{8}g_B+o(s^2)\) with \(g_B=2\|(1-J)(1+\Delta)^{-1/2}A\Omega\|^2\). No hyperfiniteness, no type restriction, no condition on \(\eta\). Key identity (Thm.~3.3, exact, any \(a\in\mathcal H\)): \(\frac14 g_B=\frac12\|(1-J)(1+\Delta)^{-1/2}a\|^2=\mathrm{dist}(a,\overline{\mathcal M'_{sa}\Omega})^2=\sup_{K\in\mathcal M_{sa}}[\varphi(K)-\mathrm{Var}_\omega(K)]\), with explicit real-orthogonal projection \(E'b=\Delta(1+\Delta)^{-1}b+(1+\Delta)^{-1}\Delta^{1/2}Jb\). Lower bound from \(\Fid=\sup_{v'\in U(\mathcal M')}|\langle\Omega,v'\xi_s\rangle|\), upper bound from Alberti's \(\Fid^2\le\omega(x)\omega_s(x^{-1})\) with \(x=1+sK\) (Alberti--Uhlmann arXiv:math-ph/0202038, Thm.~2 and est.~(1-4)). All six expressions verified numerically to 10 digits. The proof route printed in N7 (\(\Fid=\langle\Omega,\Delta_{\omega_s,\omega}^{1/2}\Omega\rangle\), Jen\v cov\'a) must be deleted and replaced by this one.
\item \textbf{[OK]} \emph{N7 \S8 item 4 (collar limit) and the hyperfinite route of the second-variation lemma} \hfill\textsc{P1}\\ Monotone convergence of the metric is proved: for \(\mathcal M_n\uparrow\mathcal M\) with \(\sigma\)-weakly dense union, \(g_B^{(n)}\uparrow g_B\) (Prop.~8.1) --- one line from the variational formula (a supremum over a smaller set), which is also the data-processing inequality for the SLD metric. Since \(\mathcal M_\varepsilon\) increases as \(\varepsilon\downarrow0\) and the tangent functionals are restrictions of one functional, the collar limit of \(g_B\) exists and is monotone; if the net is strongly additive the limit algebra is \(\mathcal A(I)\). The hyperfinite/finite-dimensional route suggested for the upper bound closes (Rmk.~8.3) but gives only that half; the lower bound needs the \(U(\mathcal M')\) argument in any case, and the finite-dimensional remainders are not uniform in \(n\) (\(p^{(n)}_{\min}\to0\) necessarily for type III).
\item \textbf{[OK]} \emph{N7 \S1 ``The deformation'' and Lemma~2.1; application of the second-variation lemma} \hfill\textsc{P1}\\ For positive collar \(\varepsilon>0\) the curve is exactly of the required form: with \(g=f\chi\in C^\infty_c\), the flow \(k_s\) of \(g\partial_x\) gives \(u_s=V(k_s)=e^{isT(g)}\), a strongly continuous one-parameter unitary group with \(\Omega\in D(T(g))\) (Fewster--Hollands arXiv:math-ph/0412028 \S5.1), \(\mathrm{Ad}\,u_s|_{\mathcal A(A_\varepsilon)}=\mathrm{id}\) and \(\mathrm{Ad}\,u_s|_{\mathcal A(D)}=\mathrm{Ad}\,U(s)\) (Carpi--Weiner arXiv:math/0407190 Def.~2.1(2)), so \(\widetilde\omega_\varepsilon=\langle u_s^*\Omega,\cdot\,u_s^*\Omega\rangle\) and Thm.~7.3 applies (Prop.~9.4). N7's ``\(\mathrm{id}\otimes\mathrm{Ad}\,U(s)\)'' presupposes the split property; the construction above does not. A general endomorphism version (Prop.~9.2, unitary implementation from Reeh--Schlieder) is also given but yields no group.
\item \textbf{[OK]} \emph{N4 Lemma~10.1 (continuity from below of fidelity)} \hfill\textsc{P1}\\ Verified. The variational formula used is Alberti--Uhlmann arXiv:math-ph/0202038 Cor.~2(5), p.~8 (\(\sqrt{P}=\inf_{x>0}\frac12\{\nu(x)+\varrho(x^{-1})\}\)); the reduction to \(m\mathbf 1\le x\le M_0\mathbf 1\) is automatic since a positive invertible \(x\) satisfies \(\|x^{-1}\|^{-1}\mathbf 1\le x\le\|x\|\mathbf 1\); Kaplansky density (BR Thm.~2.4.16) and normality complete it as written.
\item \textbf{[OK]} \emph{Note~7 Lemma~5.1(ii) / \texttt{kernel\_identity.tex} Thm~3.1} \hfill\textsc{V8}\\ Analytic proof verified in full: conventions match Note~7; the single sum/integral interchange is Tonelli-justified by \(\iint\rho\le4\log2\); residues, Gauss summation, reflection formula, diagonal limit and every sign re-derived independently. No principal value is needed --- the \(\sinh^{-2}\) pole lies off both blocks. Confirmed numerically to \(10^{-27}\) from the defining integral (\texttt{v8\_kernel\_check.py}).
\item \textbf{[OK]} \emph{rigor/quadratic\_limit.tex Thm.~4.2} \hfill\textsc{V7}\\ Unconditional lower bound \(\liminf\Phi/\zeta^2\ge g_B/8\) verified in all five ingredients: data processing under the subalgebra restriction; real-analyticity of \(\Phi_P\) on a two-sided neighbourhood with \(\Phi_P(0)=\Phi_P'(0)=0\); Legendre form of the SLD information in the normalization \(-\log\Fid=\eps^2g_B/8\); Wick reduction including the sign; density in \(L^2(N\mathrm d\kappa\mathrm d\kappa')\) since \(0<N<1\).
\item \textbf{[OK]} \emph{rigor/quadratic\_limit.tex Prop.~5.1, Lems.~5.2--5.3} \hfill\textsc{V7}\\ Independently re-derived and confirmed numerically: \(\int w^2dw/(\cosh w+\cosh a)=2a(a^2+\pi^2)/(3\sinh a)\) to \(25\) digits at five values of \(a\); the two \(u\)-integrals to \(25\) digits; \(g_B=2r/(3\pi^2)\), \(f_2=r/(12\pi^2)\). The three weight identities for \(N\), \(W\), \((q_\kappa-q_{\kappa'})^2\) hold to \(10^{-30}\).
\item \textbf{[OK]} \emph{rigor/quadratic\_limit.tex Thm.~6.4, Props.~6.5--6.6} \hfill\textsc{V7}\\ Unconditional \(\Phi\le\tau/(1-\tau)\) verified: \(\Fid\ge\operatorname{A}=\det(S_1S_2+C_1C_2)\), \(\operatorname{Re}A=1-\Theta/2\), Ostrowski--Taussky in the right direction with \(\operatorname{Re}A>0\) established, Powers--St\o rmer applied to both blocks, all bounds uniform in \(n\). The Hellinger second variation and \(h_2=r\log2/(6\pi^2)=2\log2\,f_2\) re-derived and confirmed on random four-mode data.
\item \textbf{[OK]} \emph{rigor/quadratic\_limit.tex Prop.~7.1 (repaired constant)} \hfill\textsc{V7}\\ Independent quadrature gives \(g_B-g_B^{(\Lambda)}=6.78\cdot10^{-3},1.36\cdot10^{-3},3.27\cdot10^{-4},8.10\cdot10^{-5},2.02\cdot10^{-5},5.01\cdot10^{-6}\) at \(\Lambda=10,20,40,80,160,320\), i.e.\ \(0.513\,r\Lambda^{-2}\) asymptotically. The rigorous \(2r/(3\Lambda^2)=0.667\,r\Lambda^{-2}\) holds for \(\Lambda\ge11\), and Note~5's fitted \(0.53\,\kappa_{\max}^{-2}\) is confirmed to \(3\) percent.
\item \textbf{[OK]} \emph{N7: eq.\ (1); eqs.\ (2),(3),(6)--(9); Lemma~3.1; Lemma~4.1; eqs.\ (16)--(19); Lemma~5.1; eqs.\ (22),(23); Thm~5.2; Cor.\ 6.1. N6: \S2 eqs.\ (2),(4); \S4.1} \hfill\textsc{V5}\\ \(f_2=c/(12\pi^2)\) and \(s_2=c/12\) are CORRECT. The load-bearing normalization \(\langle TT\rangle=\frac c{8\pi^2}(x-y-i0)^{-4}\) is proved independently for the free complex chiral fermion in the generator normalization: \(\langle\psi\psi^\dagger\rangle=\langle\psi^\dagger\psi\rangle=\frac1{2\pi i(x-y-i0)}\) forced by \([T(1),\psi]=-i\psi'\), \(P\ge0\), and Wick gives \(\langle TT\rangle=\frac12[(A')^2-AA'']=\frac1{8\pi^2}u^{-4}\) (also \(2\pi T\mapsto c/2\); flat-space \(\widehat W=\frac c{24\pi}|k|^3\theta(-k)\) confirms it a second time). Verified in full: \((x-y)/\sqrt{\beta\beta'}=2\sinh\frac{\xi-\xi'}2\); \(\widetilde g=-4\sinh^2\frac\xi2\) and \(\widehat{\widetilde g}=\frac{-2i}{k(k^2+1)}\) (the \(\eps\)-integral is RATIONAL in \(\eps\), so the continuation is unique --- worth saying); the sign of \(\mathfrak g=-dW_s/ds|_0\); eq.\ (D1-kernel) from Note~5's \(\frac i{2\pi L}\frac{uv}{(u+v)^2}\); all of Lemma~3.1 incl.\ \(\eta_{i\bar\jmath}=d\rho_{ij}/\sqrt{p_j}\), \(\eta=i(S^*-1)A\Omega\), \((1+\Delta)^{-1/2}\Delta^{1/2}J=J(1+\Delta)^{-1/2}\), \(d\mu(1/\lambda)=\lambda d\mu(\lambda)\), \(w(1/\lambda)=w(\lambda)/\lambda\); Lemma~4.1 incl.\ contour shift and its vertical edges, \(\sinh\frac{v-i\pi}2=-i\cosh\frac v2\), \(|\Gamma(2+ik)|^2=\pi k(k^2+1)/\sinh\pi k\); \(\sum4/((2n+1)(2n+3))=2\); the pointwise \((Q\mathfrak g)=\widetilde D^{(1)}\) with the transposition sign; \(\langle\psi_\kappa,\mathfrak g\psi_{\kappa'}\rangle=i\mathfrak g(\kappa,\kappa')\); both weights of eq.\ (22) with the corrected sign; \(|\mathfrak g|^2\), \(256\pi^6\), \(d\kappa d\kappa'=\frac12dudv\). The base integral \(\int w^2/(\cosh w+\cosh a)dw=\frac{2a(a^2+\pi^2)}{3\sinh a}\), only checked numerically in the note, is PROVED by a \(2\pi i\)-period rectangle with \(z^3\) and \(z\). Thm~5.2's assembly is right and I confirm the \(4/\pi^2\to1/\pi^2\) erratum; the two routes agree integrand by integrand, not just in the number. Cor.\ 6.1's VWZ conversions, the two-branch minimum at \(\lambda=0\), and N6's \(c/18\) per chirality all check. Independent numerics in numerics/v5chk.py, v5en.py (12--15 digits throughout).
\item \textbf{[OK]} \emph{N7 \S3 (eqs.\ (W),(What),(ghat),(mu),(gB-int),(gKM-int), Thm~5.2)} \hfill\textsc{P4}\\ Independent \(x\)-space check, no \(\xi\)-frame input. In \(u=L-x\) the modular dilation is \(u\mapsto e^{2\pi t}u\) and Mellin at \(s=ik\) diagonalises it. \(\mathsf M[(1-u-i0)^{-4}](2+ik)\) gives \(\widehat W\) exactly via \(z^{s-a}B(s,a-s)\), \(z=e^{i\pi}\); reassembly yields \(g_{\mathrm B}=2c/3\pi^2\), \(g_{\mathrm{KM}}=c/6\) to \(10^{-11}\). See rigor/collar\_and\_invariant\_formula.tex.
\item \textbf{[OK]} \emph{N7 \S3.1, Lemma~3.1 (thermality at \(\beta=2\pi\))} \hfill\textsc{P4}\\ The KMS asymmetry \(\widehat W(-k)=e^{2\pi k}\widehat W(k)\) is verified directly on the \(x\)-space Mellin transform of \(\frac{c}{8\pi^2}(u'-u-i0)^{-4}\), to 29 digits. Thermality at inverse temperature \(2\pi\) is thus a consequence of the vacuum two-point function plus the geometric dilation, not an independent input; likewise the two weight-2 density offsets are forced by homogeneity.
\item \textbf{[OK]} \emph{N5 Cor.~5.1 and the two-branch law \(-\log F^{(\lambda)}=\Phi(z(1+e^{-\pi\lambda}))+\Phi(z(1+e^{\pi\lambda}))\)} \hfill\textsc{PL}\\ Verified on the complex hopping chain (\(n=L_A{+}L_B{+}L_C\le72\), python-flint ball arithmetic, Gaussian Petz map validated against exact \(2^n\) many-body matrices). Within one family \(L_A=L_C=L\) the ratio lattice/two-branch is constant in \(\eta_V\) and equals \(1-0.15/L\) (\(0.962,0.976,0.983\) at \(L=4,6,8\)): exact in the continuum limit. See numerics/lattice/.
\item \textbf{[OK]} \emph{N5 \S(VWZ comparison); N6 Conjecture~2.1 (single-branch collapse)} \hfill\textsc{PL}\\ The single-branch law \(-\log F=\Phi(z(1+e^{-2\pi t}))\) fails on the 2D lattice by factors \(22\), \(4\times10^2\), \(5\times10^3\) at \(\lambda=0.5,1,1.5\), independent of \(\eta_V\). The two chiralities are both required, as N5 states; a single-chirality collapse must be tested on a chiral lattice model (O7(b)).
\item \textbf{[OK]} \emph{N5 abstract and \S(VWZ): ``\(f_2=0.008444\) \dots\ within \(4\%\) of their lattice value \(0.070\,c\,\eta^2\)''} \hfill\textsc{PL}\\ Confirmed and sharpened: \(\Phi_{\rm lat}/\Phi_{\rm cont}=1.03\)--\(1.07\) over the whole tabulated range \(0.009\le\zeta\le0.49\), the excess growing as \(L_A\) shrinks. The lattice-continuum gap is a finite-\(L\) ultraviolet effect of the touching points, not a discrepancy in \(f_2\).
\item \textbf{[OK]} \emph{phase1\_brief.md steps 3--5 (the Uhlmann/Bogoliubov route to \(\limsup\Phi/\zeta^2\le g_{\rm B}/8\))} \hfill\textsc{PN}\\ Numerically confirmed. With \(E(s)=\|P_-\widetilde W_sP_+\|_2^2\), the optimised extension gives \(\tfrac12E/\Phi=1.0022,1.0045,1.0130,1.0414,1.1206,1.2847\) at \(\zeta=1/120,\dots,4/15\): the bound holds and the ratio tends to \(1\) from above. See numerics/p1\_bogoliubov.\{py,out\}.
\item \textbf{[OK]} \emph{phase1\_brief.md step 5 (\(\inf_{h'}\tfrac12\|P_-(D_X+ih')P_+\|_2^2=g_{\rm B}/8\)) restricted to geometric \(h'\)} \hfill\textsc{PN}\\ The infimum is already attained by vector fields on \(I'\): minimising \(Q[g]=(48\pi^2)^{-1}\int_0^\infty k^3|\hat g|^2dk\) over extensions of the \(\zeta\)-normalised \(g\) gives \(9Q_{\min}=0.0168930\) versus \(g_{\rm B}/4=0.0168869\) and \(4.5Q_{\min}=0.0084465\) versus \(g_{\rm B}/8=0.0084434\) (0.04\%).
\item \textbf{[OK]} \emph{rigor/uhlmann\_upper\_bound.tex Thms.~9.1--9.2 (new)} \hfill\textsc{PA}\\ Sharp upper bound proved: \(\limsup_{\zeta\to0}\Phi/\zeta^2\le g_B/8=r/(12\pi^2)\), hence with quadratic\_limit Thm~4.2 the limit exists and equals \(r/(12\pi^2)=f_2\). No interchange of limits, no tail bound. Theorem~A is now a statement about \(\Phi\) itself, not only about the quantum Fisher information.
\item \textbf{[OK]} \emph{rigor/uhlmann\_upper\_bound.tex \S2 (Shale--Stinespring step)} \hfill\textsc{PA}\\ The extension of \(k_s\) is best built as the \emph{flow} of \(-\chi\partial_x\) with \(\chi=x^2\) on \(D\): then \(\widetilde k_s|_{[0,L]}=h_s\) exactly (the M\"obius semigroup is that flow), \(\widetilde W_s\) is a unitary group, and \(\lVert P_+-\widetilde W_s^*P_+\widetilde W_s\rVert_2^2\le\frac1{4\pi^2}[36\Lambda^2/c^2+64\delta^2+\frac{64}3\eta^2]\), explicit. The kink at \(0\) costs only \(\chi\in C^{1,1}\).
\item \textbf{[OK]} \emph{rigor/uhlmann\_upper\_bound.tex Lem.~5.2 (\(\langle1\rangle2\))} \hfill\textsc{PA}\\ The analytic core is an exact identity, not an estimate: \(\partial_s\big[\sqrt{\widetilde k_s'\widetilde k_s'}/(\widetilde k_s(x)-\widetilde k_s(y))\big]=\sqrt{\widetilde k_s'\widetilde k_s'}\,\kappa^{(1)}(\widetilde k_s(x),\widetilde k_s(y))\), i.e.\ \(G_s=\int_0^s\widetilde W_\sigma^*G^{(1)}\widetilde W_\sigma d\sigma\). No second-order Taylor expansion of the kernel (which would need \(\chi\in C^{2,1}\), false at the kink) is required.
\item \textbf{[OK]} \emph{rigor/uhlmann\_upper\_bound.tex Lem.~8.4 (Wick reduction)} \hfill\textsc{PA}\\ For every self-adjoint \(B\) in \(\cF(I)\) or in \(\cF(I)'\), the one-particle--one-hole part of \(B\Omega\) is \(P_-hP_+\) with \(h=h^*\) localized in the respective region (Wick uniqueness plus twisted duality; the Klein twist drops out on even vectors). Hence \(\operatorname{dist}(a,\overline{\cM'_{sa}\Omega})\) is attained in the two-particle sector, and \(g_B(a)\) equals quadratic\_limit's quasi-free \(g_B\). Verified in finite dimensions to \(5\cdot10^{-16}\).
\item \textbf{[OK]} \emph{rigor/uhlmann\_upper\_bound.tex Thm.~4.1 (\(|\langle\Omega,\Gamma(U)\Omega\rangle|=\det(1-V^*V)^{1/2}\))} \hfill\textsc{PA}\\ Ruijsenaars JMP 18 (1977) 517 and Berezin's book could NOT be obtained (paywalled; no open copy). The formula is therefore reproved from scratch here for \(\lVert P_-UP_+\rVert<1\), \(\lVert P_+UP_-\rVert<1\), via the two annihilation families, the singular value decomposition and the one-mode case (overlap \(\cos\theta\)).
\item \textbf{[OK]} \emph{rigor/uhlmann\_upper\_bound.tex Thm.~4.1 (vacuum overlap)} \hfill\textsc{PR}\\ Every step verified. Independent numerics (rigor/pr\_onishi\_check.py): \(|\langle\Omega,\Gamma(U)\Omega\rangle|=\det(1-V^*V)^{1/2}\) to \(10^{-14}\) on 2400 random unitaries, one-mode \(\cos t\) reproduced. Free source obtained: Derezi\'nski math-ph/0511030 Thm.~31(3), whose exponent \(1/4\) is on the doubled space (shots/PR\_onishi\_*.png).
\item \textbf{[OK]} \emph{rigor/uhlmann\_upper\_bound.tex Prop.~7.5 (\(\inf_{h'}=g_B/4\))} \hfill\textsc{PR}\\ Confirmed in an exactly solvable finite-dimensional surrogate (rigor/pr\_step5\_chain.py): \(\inf_{h'}\|P_-(D+ih')P_+\|_2^2\), \(\operatorname{dist}(a,\mathcal M'_{\rm sa}\Omega)^2\), \(\sup_{\mathcal M_{\rm sa}}[\varphi-\operatorname{Var}]\) and \(\tfrac14\sup_\ell\Omega(\ell)\) agree to \(5\cdot10^{-15}\) on six random instances, factor \(\tfrac14\) included; twisted and untwisted commutants give the same projection of \(a\).
\item \textbf{[OK]} \emph{rigor/uhlmann\_upper\_bound.tex Lem.~5.2 and Rem.~5.6} \hfill\textsc{PR}\\ Integral representation verified numerically to \(10^{-12}\) up to \(s=1\); \(\kappa^{(1)}(-u,v)=uv/(u+v)^2\) to \(10^{-16}\). Stronger consequence, not drawn in the note: \([P_-,ih']\) vanishes on \(A\times D\), so \(\inf_{h'}\|P_-(D+ih')P_+\|_2^2\ge1/(24\pi^2)=\tfrac14(g_B/4)\) unconditionally --- a real check on the claimed infimum.
\item \textbf{[OK]} \emph{rigor/uhlmann\_upper\_bound.tex Thm.~8.2 (referee verdict)} \hfill\textsc{PR}\\ No proof-breaking issue found in any of the five ingredients. Subject to the two-line repair of Lem.~7.3, and to Note~5 Lem.~1.1 (cross-ratio reduction, imported and not re-verified here), \(\lim_{\zeta\to0}\Phi(\zeta)/\zeta^2=g_B/8=c/(12\pi^2)\) stands for the free chiral fermion; the two halves use the same \(\Phi\), \(\zeta\) and root fidelity.
\item \textbf{[OK]} \emph{rigor/tail\_bound.tex Thm.~5.1 (UV tail, O4)} \hfill\textsc{PT}\\ Unconditional: \(0\le\Phi-\Phi_W\le0.173\zeta^2/\kappa_c+0.88\zeta^2/\kappa_c^2\) for \(\kappa_c\ge\max\{10,2\log(1/\zeta)+2\log(1+\kappa_c)+3\}\). Route: exact factorisation \(\Fid(\rho,\rho'')=\Fid_W\), the purification bound \(\Fid\ge\det(1-N^*N)^{1/2}\), and \(\|N\|_2^2\le5\tau+\nu\) from \(\widetilde Q\ge0\), operator concavity of \(\sqrt{t(1-t)}\) and a Schur complement. [corrected 2026-09-24: not unconditional. The shape of the bound is proved; the explicit constants \(0.173\) and \(0.88\) are \texttt{eq:final} of Cor.~5.2 (\texttt{cor:rate}), not of Thm.~5.1, and rest on numerical inputs, \(A_1\le2.0\) and \(A_3\le0.0374\) in double precision, not certified, and \(f=0.0085\); Thm.~5.1 holds with \(Y=e^{\Phi_W}y\) in place of the printed \(y\) (rigor/tail\_bound.tex, Corrections~C1, C3, C4, C10); card \texttt{tail-bound}]
\item \textbf{[OK]} \emph{rigor/tail\_bound.tex Thm.~4.2 (Note~5 Prop.~4.1 upgraded)} \hfill\textsc{PT}\\ Proved: \(\tau_+=\tau_-\) exactly (parity of \(\operatorname{Im}J\)), and \(|\int_{\kappa_c}^\infty(\widetilde Q(\kappa,\kappa)-q_\kappa)d\kappa-\zeta^2/(15\kappa_c^2)|\le0.46\zeta^2\kappa_c^{-3}+0.051\zeta(1+\kappa_c)e^{-\kappa_c}\). Re-derives \(c_2=-\zeta^2/30\) from the closed form. Numerically the remainder is 6--30 times inside the bound. [corrected 2026-09-24: proved in form, with finite \(C_2=4\pi A_3\), \(C_1=A_1/(4\pi^2)\); the values rest on \(A_3\le0.0374\) (so \(C_2\le0.47\); 0.036 and 0.46 before 2026-09-24) and \(A_1\le2.0\) (\(C_1\le0.051\)) in double precision, not certified (rigor/tail\_bound.tex, Corrections~C4, C10); with \(C_2\le0.47\) the remainder is 6--31 times inside the bound (6--30 before Correction~C10); card \texttt{tail-bound}]
\item \textbf{[OK]} \emph{rigor/tail\_bound.tex Lem.~3.1 (purification bound), numerical} \hfill\textsc{PT}\\ \(\Fid(\omega_X,\omega_Y)\ge|\det(X^{1/2}Y^{1/2}+(1-X)^{1/2}(1-Y)^{1/2})|\) verified in 40-digit arithmetic on 40 random non-commuting \(3\times3\) pairs (max ratio 1.000) and 30 perturbative \(4\times4\) pairs (bound tight to 4\%). Double precision reports spurious violations once \(\min\operatorname{spec}\lesssim10^{-9}\); use extended precision.
\item \textbf{[OK]} \emph{rigor/tail\_bound.tex Thm.~4.1 (block estimate), numerical} \hfill\textsc{PT}\\ \(\|N(\widetilde Q,Q'')\|_2^2\le5\tau+\nu\) holds on the free-fermion data with ratio 0.38--0.61 at \((\Lambda,\kappa_{\max})=(60,30),(120,45)\), \(\kappa_c=18.4,25.3,32.2\). Also confirmed: \(\tau_+=\tau_-\) to all digits at \(\Lambda=60\), and the box tail equals \(\Sigma(\kappa_c)-\Sigma(\kappa_{\max})\) to 4\%.
\item \textbf{[OK]} \emph{numerics/results\_largezeta\_certified.tex (continuum \(\Phi\) at large \(\zeta\))} \hfill\textsc{PZ}\\ Certified \(\Phi(\zeta)=-\log\Fid\): \(4.410(20)\times10^{-3}\), \(1.0797(50)\times10^{-2}\), \(2.2414(14)\times10^{-2}\), \(3.974(7)\times10^{-2}\), \(6.41(43)\times10^{-2}\) at \(\zeta=1.07,2.13,4.27,8.53,17.07\) (\(\Lambda=120\), \(\kappa_{\max}=45\), windows \(\kappa_c=18.4\)--\(43.7\)). The 1\% target is met for \(\zeta\lesssim8\).
\item \textbf{[OK]} \emph{numerics/README.md (lattice/continuum comparison)} \hfill\textsc{PZ}\\ Continuum and lattice \(\Phi\) agree to 1.1\%, 1.7\%, 3.0\% at \(\zeta=1.07,2.13,4.27\) (lattice finite-\(L\) UV effect 3--7\%), and the local log-slopes \(1.29,1.17,0.94\) vs \(1.41,1.18,0.98\) track each other; \(\Phi\) flattens monotonically and is not a power law.
\item \textbf{[OK]} \emph{rigor/implementation\_C11.tex Thm.~5.1 (H1),(H2),(H2$'$),(H3), problem O5} \hfill\textsc{QB}\\ Proved for \emph{every} diffeomorphism covariant net, any \(c\): \(T(g_{\rm tot})\) is essentially self-adjoint on the finite-energy domain (Carpi--Weiner Thm.~4.4, needing only \(\lVert g_{\rm tot}\rVert_{3/2}<\infty\)), \(U_s=e^{isT(g_{\rm tot})}\); covariance by Trotter plus localisation and by approximation with additivity and semicontinuity. Scalar phase on reducible representations: Lemma~4.2 (Fewster--Hollands Prop.~5.1).
\item \textbf{[OK]} \emph{rigor/implementation\_C11.tex Thm.~2.3 (regularity of the first-order field)} \hfill\textsc{QB}\\ \(G\) is \(C^{1,1}\), piecewise smooth, one jump \(J=-1/L\) of \(G''\) at \(0\); \((in)^3\hat G_n\to J/2\pi\); \(\lVert G\rVert_{3/2}<\infty\); \(G\in H^{5/2-\epsilon}\setminus H^{5/2}\), so \(\widetilde k_s\in D^\sigma\) iff \(\sigma<5/2\) and CDIT's \(\sigma>3\) theorem does not apply. Numerics (rigor/qb\_regularity.py) confirm all six claims to 5--6 digits.
\item \textbf{[OK]} \emph{rigor/implementation\_C11.tex Rem.~2.4 (why \(C^{1,1}\), and Gap~2 of the second-variation note)} \hfill\textsc{QB}\\ \(\lVert T(g)\Omega\rVert^2=\tfrac c{12}\sum n(n^2-1)|\hat g_n|^2\) is the Weil--Petersson \(\dot H^{3/2}\) norm; it is finite exactly on that class. A field with a corner has \(|\hat g_n|\asymp n^{-2}\) and no tangent vector, whereas \(g_{\rm tot}\) has \(n^{-3}\). One full derivative of margin, consumed again at the next order.
\item \textbf{[OK]} \emph{rigor/universality\_theorem\_B.tex Thm.~7.1 (Theorem~B)} \hfill\textsc{QA}\\ Proved: for any diffeomorphism covariant net with (HD), (BW), (RS) and (H1)--(H3), \(-\log\Fid(\omega,\omega_s)=\frac{s^2}{8}g_B+o(s^2)\), \(g_B=4\operatorname{dist}(T(g_{\rm tot})\Omega,\overline{\mathcal A(I)'_{\rm sa}\Omega})^2=2c/(3\pi^2)\), \(f_2=c/(12\pi^2)\) at \(\zeta=s\). Only (H1)--(H3) remain conditional; no collar limit and no domain condition \(T(g_{\rm tot})\Omega\in D(\Delta^{-1/2})\) is used.
\item \textbf{[OK]} \emph{rigor/universality\_theorem\_B.tex Thm.~5.5 (\(c\)-linearity)} \hfill\textsc{QA}\\ New unconditional result: the stress-tensor subspace \(\mathcal H_T\) reduces \(\Delta\) and \(J\), and the minimiser \(E'a\) of the distance lies in it, so \(g_B\) is computed inside \((\mathcal H_T,\Delta,J,a)\), whose inner product is \(c\) times a net-independent form. Hence \(g_B=c\gamma_0\); \(\gamma_0=2/(3\pi^2)\) from the rigorous free-fermion value.
\item \textbf{[OK]} \emph{rigor/universality\_theorem\_B.tex \S6.2, rigor/vB2.py (numerical check)} \hfill\textsc{QA}\\ Rayleigh--Ritz on the net-free problem \(\inf\lVert G\rVert_{3/2}^2\) over real \(G\) equal to \(-x^2\) on \((0,1)\) and \(0\) on \((-\infty,0)\), free on \((1,\infty)\): \(2.5678,2.5630,2.5579\) for boxes \(512,1024,2048\), decreasing to \(8/\pi=2.5465\) (0.45\%). The minimum is independent of the continuation to six digits.
\item \textbf{[OK]} \emph{rigor/kubo\_mori\_second\_variation.tex Thm.~``Exact formula''} \hfill\textsc{QK}\\ For \(A=A^*\) affiliated with \(\mathcal M\), \(\omega_s=\omega\circ\mathrm{Ad}\,e^{isA}\), the Araki relative entropy is exactly \(D(\omega\Vert\omega_s)=\lVert h(\Delta)^{1/2}(e^{isA}-1)\Omega\rVert^2\) with \(h(\lambda)=\lambda-1-\log\lambda\ge0\), via \(\Delta_{u\Omega,\Omega}=u\Delta u^*\). Verified to \(10^{-12}\) in finite dimensions. [corrected 2026-09-24: \texttt{thm:exact} is for bounded \(A=A^*\in\mathcal M\); for \(A\) affiliated with \(\mathcal M\) the identity is \texttt{prop:unbounded}, which is stated with \(\Omega\in D(A)\). In its proof the identity uses only \(e^{isA}\in\mathcal M\) (steps \(\langle1\rangle1\)--\(\langle1\rangle2\)), while \(\Omega\in D(A)\) enters through \texttt{lem:kms} (step \(\langle1\rangle2\)) and the lower bound (step \(\langle1\rangle3\), Stone's theorem); the second-order expansion then holds under \((\mathrm U_\varepsilon)\), since modular analyticity \(\Rightarrow(\mathrm U_\varepsilon)\) is proved only for bounded \(A\); a generator affiliated only with a larger algebra is open (Gap QK2); card \texttt{km-exact-identity}]
\item \textbf{[OK]} \emph{rigor/kubo\_mori\_second\_variation.tex Thm.~``Lower bound''} \hfill\textsc{QK}\\ \(\liminf_{s\to0}s^{-2}D(\omega\Vert\omega_s)\ge\frac12 g_{\mathrm{KM}}\) holds in \([0,\infty]\) for every von Neumann algebra with no hypothesis, because \(v\mapsto\lVert h(\Delta)^{1/2}v\rVert^2\) is a closed, hence lower semicontinuous, form and \(s^{-1}(e^{isA}-1)\Omega\to iA\Omega\). No martingale argument needed.
\item \textbf{[OK]} \emph{rigor/kubo\_mori\_second\_variation.tex Thm.~``Kubo--Mori second variation''} \hfill\textsc{QK}\\ Gap~1 of rigor/lemma\_second\_variation.tex is closed for unitary orbits: \(D=\frac{s^2}2 g_{\mathrm{KM}}+o(s^2)\) whenever \(\lVert\Delta^{-\epsilon/2}(e^{isA}-1)\Omega\rVert=O(s)\) for some \(\epsilon>0\), in particular whenever \(A\) is modularly analytic in a strip of any positive width.
\item \textbf{[OK]} \emph{universality\_normalization.tex eq.~(3.27) (free-fermion \(g_{\mathrm{KM}}\))} \hfill\textsc{QK}\\ Because \(\log q_\kappa-\log(1-q_\kappa)=-\kappa\), the two logarithmic-mean weights collapse and the double integral reduces exactly: \(\int\sigma^2[q(\frac{\sigma-u}2)-q(\frac{\sigma+u}2)]d\sigma=\frac23 u(u^2+4\pi^2)\), giving \(g_{\mathrm{KM}}=\frac13\int du/(u^2+4\pi^2)=1/6\) in closed form. Checked to \(10^{-13}\).
\item \textbf{[OK]} \emph{refs/P1\_CarpiWeiner\_math-0407190.pdf Thm.~4.4 (pp.~11--14), as used for (H1)} \hfill\textsc{QR2}\\ Proof re-read: it uses the \(\epsilon\)-uniform bound \(\lVert[L_n,e^{-\epsilon L_0}]\rVert\le q^{1/2}\lvert n\rvert^{3/2}\) and \(A^+e^{-\epsilon L_0}x\to A^*x\), not Nelson's commutator theorem; the only hypothesis is \(\sum\lvert a_n\rvert\lvert n\rvert^{3/2}<\infty\). Their Lemma~4.1 constant is explicitly independent of the lowest weight and their \S4 is written for reducible representations. QB's use is correct.
\item \textbf{[OK]} \emph{rigor/implementation\_C11.tex Thm.~5.1} \hfill\textsc{QR2}\\ Confirmed: (H1)--(H3) hold for every diffeomorphism covariant conformal net (CDIT \S4 axioms, separable Hilbert space), any \(c\), no strong additivity, no finite multiplicities. Every step checked: Trotter hypotheses, the induction \(h_aD\subseteq D\) (\(s\ge0\)), the straddling-interval limit, Stone differentiability and signs. Only the phase gap above needs a sentence.
\item \textbf{[OK]} \emph{rigor/qb\_regularity.py (all numerical claims of implementation\_C11.tex)} \hfill\textsc{QR2}\\ Re-run: \(n^3\lvert\hat G_n\rvert=0.159155\) versus \(\lvert J\rvert/2\pi=0.1591549\); \(\lVert G\rVert_{3/2}=18.70397\) with tail \(2\lvert J\rvert/(\pi\sqrt N)\); \(\lVert G'\rVert_{3/2}\) increment \(162.971\) versus \(162.975\); \(\lVert T(g)\Omega\rVert^2=0.169865c\); \(\lVert L_0^{1/2}T(g)\Omega\rVert^2=1.622737c\); \(L_0^2\) moment \(+0.1405\) per factor \(16\). All reproduced.
\item \textbf{[OK]} \emph{phase2\_brief.md track QN (O6); numerics/boson/} \hfill\textsc{QN}\\ Universality confirmed for the U(1) current net (\(c=1\)). Gaussian mode compression in the modular frame gives \(f_2=0.00844(2)\) versus \(1/(12\pi^2)=0.0084434\) and \(s_2=0.083(2)\) versus \(1/12\), matching the Dirac-fermion values \(0.008444(1)\), \(0.0837(5)\) to \(0.2\%\) and \(1\%\).
\item \textbf{[OK]} \emph{numerics/boson/README.md \S2 (bosonic modular symbol)} \hfill\textsc{QN}\\ In the modular frame \(t=\log\frac{x+a}{L-x}\) the vacuum of \(J\) restricted to \(I\) is translation invariant with covariance symbol \(M(\kappa)=\pi\kappa\coth\pi\kappa\), i.e. \(\nu(\kappa)=\coth\pi\kappa\), the Bose factor. The interior kernel \(1/(4\sinh^2\frac{t-t'}2)\) alone gives \(M-1\); the missing \(+1\) is the collar term \(2\int f g\,dt\) from \(y\notin I\).
\item \textbf{[OK]} \emph{phase2\_brief.md track QN, ``bosonic overlap = fermionic \(\det(1-V^*V)^{1/2}\)''} \hfill\textsc{QN}\\ Proved analytically, not just numerically. For \(\varphi=x+sg\) the pair blocks are \(W(p,q)=\frac{is}{4\pi}(p+q)\hat g(p-q)\) (weight \(1/2\)) and \(M(p,q)=\frac{is}{2\pi}\sqrt{|p/q|}\,q\,\hat g(p-q)\) (weight \(0\)), so \(\mathrm{Tr}(B^*B)=2\lVert P_-WP_+\rVert_2^2\) and both \(-\log|\langle\Omega,\Gamma\Omega\rangle|\) equal \(\frac{s^2}{96\pi^2}\int_0^\infty k^3|\hat g|^2dk\). Numerics: ratio \(0.99923\).
\item \textbf{[OK]} \emph{rigor/universality\_theorem\_B.tex, Props. 4.2--4.3 and Lemma 5.3} \hfill\textsc{QR1}\\ Uhlmann is used in the correct (trivial) direction; \(U(\phi')\in\mathcal A(I)'\) holds with \(\operatorname{supp}\phi'\Subset I'\). No cross term arises because the trial unitary is moved onto \(\Omega\): only two separate first-order expansions are subtracted. The Schwarzian genuinely vanishes: only the Moebius maps \(\delta_t,r\) act on \(\mathcal H_T\).
\item \textbf{[OK]} \emph{rigor/universality\_theorem\_B.tex \S6 normalisation chain; new check} \hfill\textsc{QR1}\\ Verified end to end, including the previously unwritten link: for the free fermion \(\lVert\Pi_-(f\partial+\tfrac12f')\Pi_+\rVert_2^2=\frac1{48\pi^2}\int_0^\infty p^3|\hat f|^2dp\), i.e. exactly (S2) at \(c=1\) (numerically \(1.5\times10^{-11}\)). Also \(\zeta=as/(a+L)\) re-derived: \(a=\infty\) is the Moebius normal form, not a special case.
\item \textbf{[OK]} \emph{rigor/vB2.py and Lemma 5.8(4) (admitted GAP)} \hfill\textsc{QR1}\\ vB2 reproduces 2.5678/2.5630/2.5579. A box-free recomputation with exact Fourier transforms of \(u^ke^{-\alpha u}\) gives \(\inf\lVert G\rVert_{3/2}^2=2.547971\) versus \(8/\pi=2.546479\) (\(5.9\times10^{-4}\), monotone), identical to seven digits for three different continuations. So the geometric trial family is optimal to within 0.06\%.
\item \textbf{[OK]} \emph{rigor/implementation\_C11.tex Lem.~4.4 (locality of implementers; (H2') for QR1)} \hfill\textsc{QB}\\ \(U(\mathrm{Diff}(K))\subseteq\mathcal A(K)\) for \(\mathrm{Diff}(K)=\{\gamma:\gamma|_{K'}=\mathrm{id}\}\) (well-posed form of the locality clause). And \(T(g_{\rm tot})\Omega\) lies in the \(T\)-subspace: \(\lVert T(f_n)\Omega-T(g_{\rm tot})\Omega\rVert=\lVert T(f_n-g_{\rm tot})\Omega\rVert\le r(c)\lVert f_n-g_{\rm tot}\rVert_{3/2}\to0\) for every smooth approximating sequence.
\item \textbf{[OK]} \emph{rigor/uhlmann\_upper\_bound.tex Step 4, made quantitative (rigor/rate\_of\_remainder.tex Thm.~2.3)} \hfill\textsc{RR}\\ New exact statement: for any admissible flow extension, \(\lVert P_-\widetilde W_sP_+\rVert_2^2\le s^2\lVert P_-DP_+\rVert_2^2\) for all \(s\), from \(\kappa_{\tilde k_s}=\int_0^s F_\sigma\,d\sigma\), \(\langle F_\sigma,F_\tau\rangle=R(\sigma-\tau)\), \(|R|\le R(0)\). Hence a constant-free global bound; \(E(s)\) is even; no cubic term.
\item \textbf{[OK]} \emph{Note~5 \S2, third order (rigor/rate\_of\_remainder.tex Thm.~3.1)} \hfill\textsc{RR}\\ Third-order coefficient of \(-\log\mathrm{Fid}\) for quasi-free states derived in the eigenbasis of \(Q\): \(\Phi_3=\tfrac16\mathrm{Tr}[(C^{-2}D)^3]-\mathrm{Tr}M_3+\mathrm{Tr}(M_1M_2)-\tfrac13\mathrm{Tr}M_1^3\), \(M_j=(1-Q)T_j\) from the Sylvester chain. Verified against the exact determinant formula on random examples (residuals exactly linear in \(\varepsilon\)).
\item \textbf{[OK]} \emph{Note~5 Table 3 / phase2b brief, value of \(c_3\)} \hfill\textsc{RR}\\ \(\Phi=f_2\zeta^2(1+c_3\zeta+c_4\zeta^2+\dots)\) with \(c_3=-0.99(1)\) and \(c_4=+0.9(1)\), from the third-order formula applied to the exact defect matrices in the box/modular frame (two boxes, seven windows, Richardson in \(1/\kappa_c^2\)). The brief's \(-0.96\) is the same number contaminated by the positive quartic term.
\item \textbf{[OK]} \emph{tail\_bound.tex, Table~1 (certified error bars)} \hfill\textsc{RC}\\ Improved. Rigorous: the box symbol is enclosed to $5.62\times10^{-18}$ per entry, and $\Phi_W\le\Phi^{\rm box}\le1.20\,\Phi_W$ at $\kappa_c=25.3$ (widths $17.8$--$19.9\%$, was $82$--$154\%$). The sharper computable form of the bound reaches the percent level; see certified\_numerics.tex, \S9. [corrected 2026-09-24: the bars of tail\_bound.tex, Table~1 (widths \(82\)--\(154\%\)) are not certified. Their upper end is \texttt{eq:final} of Cor.~5.2, whose constants \(0.173\) and \(0.88\) rest on numerical inputs (\(A_1\le2.0\) and \(A_3\le0.0374\) in double precision, not certified; \(f=0.0085\)); tail\_bound.tex now says so (its Corrections~C8 and C10). Of the ``rigorous'' statements only the entrywise enclosure \(5.62\times10^{-18}\) has no unproved input; \(\Phi_W\le\Phi^{\rm box}\le1.20\,\Phi_W\) holds modulo the quadrature of \(\widehat D\) and the box cutoff (the finite dimensionality of the box), with the measured constants \(c_0\in[18,42]\), \(c_1\le1.011\) of hypothesis \texttt{eq:Hyp} (certified\_numerics.tex, \S9 (RCr), items (C1), (C2), and its Correction~C2); cards \texttt{tail-bound}, \texttt{certified-numerics}]
\item \textbf{[OK]} \emph{kubo\_mori\_gauge\_lemma.tex, Theorem 3.3} \hfill\textsc{RK}\\ Unconditional and exact: $D_{\mathcal A(I)}(\omega_s\Vert\omega)\le\lVert h(\Delta_K)^{1/2}(\hat U_{-s}-1)\Omega\rVert^2$ for every admissible gauge and every $s\ge0$, by Araki monotonicity plus automorphism invariance plus QK's exact identity, the generator $T(\hat g)$ being affiliated with $\mathcal A(K)$ by Stone.
\item \textbf{[OK]} \emph{kubo\_mori\_gauge\_lemma.tex, Theorem 5.1 and Corollary 5.2} \hfill\textsc{RK}\\ For $T(f)$ affiliated with $\mathcal A(K)$ the Kubo--Mori weight cancels the KMS denominator exactly: $g_{\mathrm{KM}}^{(K)}=\frac{c}{24\pi}\int k^2(k^2+1)\lvert\hat{\tilde f}\rvert^2dk$, a local Sobolev functional, and $\mathrm{Var}^{\mathrm{KM}}=\frac{c}{48\pi^2}(\lVert\tilde f'\rVert^2+\lVert\tilde f\rVert^2)$. Both $c$-linear, net-independent; verified to $5\cdot10^{-9}$ in rk\_check.py.
\item \textbf{[OK]} \emph{kubo\_mori\_gauge\_lemma.tex, Theorem 7.1} \hfill\textsc{RK}\\ Petz--Donald's supremum formula gives the lower bound in the right direction: $\liminf s^{-2}D\ge\frac12\sup_X[2\varphi(X)-\mathrm{Var}^{\mathrm{KM}}_\omega(X)]\ge\frac12 g_B=c/(3\pi^2)$, using that the Kubo--Mori variance is dominated by the ordinary variance. Combined with the upper bound this brackets $s_2$ in $[0.034c,0.924c]$.
\item \textbf{[OK]} \emph{cmi\_cft\_research\_program.tex \S``Thermal and KMS states'', eq.~(thermal-asymp); numerics/lattice/petz\_thermal\_rg.py} \hfill\textsc{S9}\\ Exponential shielding confirmed on the hopping chain. With \(v=2\) sites per unit lattice time, \(\beta_{\rm len}=2\beta\) and \(\xi=\beta/\pi\) sites: fits of \(e^{-\kappa L_B/\beta}\) give \(\kappa_I=3.129,3.159,3.150\) (\(\pi\)) and \(\kappa_F=5.863,6.237,6.293\) (\(2\pi\)) at \(\beta=4,8,16\), independent of \(L_A=L_C\).
\item \textbf{[OK]} \emph{O7(a) two-branch law at finite temperature; numerics/lattice/petz\_thermal\_rg.out \S(1)} \hfill\textsc{S9}\\ The vacuum two-branch law holds for the KMS state at the thermal cross ratio \(\eta_\beta=\sinh(\pi a/\beta_{\rm len})\sinh(\pi c/\beta_{\rm len})/[\sinh(\pi(a+b)/\beta_{\rm len})\sinh(\pi(b+c)/\beta_{\rm len})]\), the image of the four endpoints under the exponential map: ratio \(0.98\)--\(1.01\) over five decades at \(\beta=16\).
\item \textbf{[OK]} \emph{cmi\_cft\_research\_program.tex eq.~(thermal-cmi), first-order expansion} \hfill\textsc{S9}\\ Expanding gives \(I_\beta-I_{\rm vac}=-\pi^2L_AL_C/(36\beta^2)\), independent of \(L_B\). Measured/predicted \(=0.300,0.608,0.855,0.958\) at \(\beta=4,8,16,32\) (\(L_A=L_C=8\), \(L_B=1\)), i.e. \(1-O(\beta^{-2})\); at \(\beta=32\) the value varies by 3\% between \(L_B=1\) and \(4\).
\item \textbf{[OK]} \emph{cmi\_cft\_research\_program.tex \S``A candidate Markov c-function'', eqs.~(cM), \(c_{\rm E}=3RS'\)} \hfill\textsc{S9}\\ Staggered-mass flow (\(\mu=m/2\)): \(c_{\rm E}=1.00090\), \(c_{\rm M}=1.00144\) at \(m=0\); both monotone decreasing to \(0\), \(c_{\rm M}\ge c_{\rm E}\) throughout, and \(c_{\rm M}/c_{\rm E}\to 2\mu R=mR\). The 2--3\% overshoot \(c_{\rm M}>1\) at \(R=6,8\) does not collapse in \(\mu R\): lattice artefact.
\item \textbf{[OK]} \emph{Gap geometry \(A\,|\,G\,|\,B\,|\,C\); numerics/lattice/petz\_thermal\_rg.out \S(3)} \hfill\textsc{S9}\\ For the massless Dirac field \(I(A{:}C|B)=\tfrac13\log[(a+g+b)(g+b+c)/((g+b)(a+g+b+c))]\), the adjacent law with \(L_B\to L_B+L_G\) (confirmed to 0.06\% at \(L_G=96\)), so \(I\sim L_G^{-2}\). \(-\log F^{(0)}\) decays with the same power, with \(L\)-independent ratio \(\kappa_g\approx0.09\).
\item \textbf{[OK]} \emph{rigor/optimality\_second\_order.tex Thm.~3.1 and Cor.~3.3; numerics/optimality/opt\_gaussian.py} \hfill\textsc{S8}\\ Second-order optimality over gauge-invariant quasi-free channels is an exactly convex programme: in the variables \((X,Z)\), \(Z\) the output symbol, the feasible set \(XQ_BX^*\le Z\le 1-X(1-Q_B)X^*\) is convex (operator convexity of \(X\mapsto XRX^*\)) and semidefinite representable, and the defect is affine.
\item \textbf{[OK]} \emph{rigor/optimality\_second\_order.tex Prop.~2.2; numerics/optimality/cp\_check.py} \hfill\textsc{S8}\\ Complete positivity of a gauge-invariant quasi-free channel \(Q\mapsto XQX^*+Y\) is exactly \(Y\ge0\), \(XX^*+Y\le1\); the Gaussian Stinespring dilation with \(V_E=(1-XX^*)^{1/2}\), \(Q_E=V_E^{-1}YV_E^{-1}\) was built explicitly on a four-mode Fock space and reproduces the covariance rule to \(6.7\times10^{-16}\).
\item \textbf{[OK]} \emph{rigor/optimality\_second\_order.tex \S5.2 Table~1; numerics/optimality/run\_exact.py} \hfill\textsc{S8}\\ The numerical optimum sits at the two-chirality geometric-compression value: best/\((2f_2z^2)=1.006\), \(1.009\), \(1.124\) at \((L_A,L_B,L_C)=(4,8,2)\), \((8,16,4)\), \((4,12,2)\), and no channel found anywhere goes below \(2f_2z^2\). Optimal coefficient per chirality: \(f_2z^2\), i.e. one quarter of Petz.
\item \textbf{[OK]} \emph{rigor/sdp\_dual\_certificate.tex Thm.~3.2, Lemma~11.1, Cor.~11.2} \hfill\textsc{S10}\\ Two exact structural results for the continuum programme. (i) Certificate: because the zero-collar compression is a \emph{unitary} $V_c:\mathfrak h_D\to\mathfrak h_B$ (a continuum feature; impossible on a lattice), both LMI constraints are active, the KKT multipliers are unique, $\Sigma_1=t_{DA}Q_{AB}V_c-t_{DD}V_c^*(1-Q_{BB})V_c$, and optimality holds iff $\Sigma_1=\Sigma_1^*\ge0$ and $\Sigma_1+t_{DD}\ge0$ --- necessary as well as sufficient, with no constraint qualification. (ii) Mobius rigidity: the $DD$-block of the compression defect vanishes at first order, $E_D[P,D_\chi]E_D=[E_DPE_D,E_DD_\chi E_D]=0$, because the Mobius generator commutes with the Hardy projection and is a local (first-order differential) operator. Hence the Riesz vector's corner $E_DuE_D$ is \emph{purely anti-self-adjoint} and the self-adjointness half of the certificate can only hold trivially ($E_DuE_D=0$). Also: for $t$ supported in $I$, $\operatorname{Tr}(tQt(1-Q))=\lVert P^\perp tP\rVert_2^2$ and the Legendre ratio depends on $t$ only through $P^\perp tP$, so the two $P$-diagonal blocks of a test operator are free.
\item \textbf{[OK]} \emph{rigor/optimality\_all\_channels.tex Thm.~3.5 (with Lemmas 3.1--3.4), Prop.~2.4; refs/BJL\_twisted\_duality\_math-ph-0204029.pdf Prop.~5.3, Lemma~4.10; refs/AlbertiUhlmann02 Thm.~2(1)} \hfill\textsc{ORCH-3b}\\ S8's sharpened open question (Verdict Q4) is answered affirmatively: the recovered symbol of \emph{every} recovery channel (normal, unital, CP, even) lies in the quasi-free feasible set \(\mathcal F\). Proof: Kadison--Schwarz for the pair \(a(g),a^*(g)\) gives \(Z\ge\mathrm{Gram}(\xi)\), \(1-Z\ge\mathrm{Gram}(S_B\xi)\) with \(\xi_g=\beta(a(g))\Omega\); projecting to the one-particle charge \(-1\) sector preserves these (the projections commute with \(\Delta_B\), Baumg\"artel--Jurke--Lled\'o Cor.~5.4); and \(D(S_B)\cap\{\text{one-particle}\}=\{B(q)\Omega:q\in\mathfrak h_B\oplus\Gamma\mathfrak h_B\}\) (BJL Prop.~5.3 with the closed graph of Lemma~4.10), so \(\xi_g\) becomes \(a(X^*g)\Omega\) with \(X^*g\in\mathfrak h_B\). Moreover the exact single-observable bound from the Alberti--Uhlmann variational formula gives, for any channel, \(-\log F\ge g_Q(\text{defect})(1-o(1))\) in Legendre form (Prop.~2.4). Lattice consistency: symbols of the all-channel SDP optimum, of Petz, and of random and random isometric non-Gaussian channels all lie in \(\mathcal F\) (numerics/optimality\_all/symbol\_in\_F\_test.py).
\item \textbf{[OK]} \emph{rigor/optimality\_all\_channels.tex Thm.~4.1; refs/Sion1958\_minimax.pdf Cor.~3.3} \hfill\textsc{ORCH-3b}\\ The second-order recovery problem over all channels has the same value as the quasi-free convex programme of S8 Cor.~3.3: the dual value \(\sup_t\Delta_z(t)_+^2/(8V(t))\) equals \(\min_{\mathcal F}g_Q\) by Sion's minimax theorem (the feasible set is bounded and closed in the weak operator topology because its constraints are affine positivity conditions), and every channel's error is bounded below by the dual value up to the remainder control of Prop.~2.4. Numerically the exact all-channel optimum on tiny chains (one SDP over Choi matrices, numerics/optimality\_all/all\_channels\_sdp.py) is a Gaussian state to \(10^{-6}\).
\item \textbf{[OK]} \emph{rigor/exact\_fidelity\_upper\_half.tex Thms.~5.3 and 5.4 (question (Q)), Thm.~4.3 (reduction), Thms.~6.2 and 7.1 (UH), Thm.~7.2} \hfill\textsc{E1}\\ The decisive open statement of the previous entry is \emph{proved}: with \(\mathcal N=\{P^{\perp}\ell P:\ell=\ell^{*}=E_I\ell E_I\}\) and \(\mathrm{Ext}=\{P^{\perp}Z_eP:Z_e^{*}=-Z_e=E_{I^c}Z_eE_{I^c}\}\) inside the real Hilbert space \(\mathfrak S_2(PH,P^{\perp}H)\), one has \(\overline{\mathrm{Ext}}=\mathcal N^{\perp}\), finite-rank \(Z_e\) already sufficing. Proof (invariant, two projections): for \(B\) orthogonal to both spaces put \(\Xi=B+B^{*}\), \(\Gamma=2P-1\); then \(\Gamma\Xi+\Xi\Gamma=0\), \(E_I\Xi E_I=0\), \(E_{I^c}\Gamma\Xi E_{I^c}=0\), and compressing \(\Gamma\Xi\Gamma=-\Xi\) with \(E_{I^c}\) while inserting \(1=E_I+E_{I^c}\) twice gives \(\Gamma'\pi\Gamma'=\pi\) for \(\pi=E_{I^c}\Xi E_{I^c}\), \(\Gamma'=E_{I^c}(2P-1)E_{I^c}\); generic position of \((P,E_I)\) (Hardy uniqueness) makes \(\Gamma'\) a self-adjoint contraction without eigenvalues \(\pm1\), so \(\Gamma'^n\to0\) strongly and the compact \(\pi=\Gamma'^n\pi\Gamma'^n\) vanishes; then injectivity of \(E_{I^c}PE_I\) kills the off-diagonal blocks. Independently, (Q) \emph{is} Haag (twisted) duality read in the one-particle--one-hole sector (Tomita: \(K^{\perp_{\mathbb R}}=iK'\), with PA Lemma~5.5). Consequences: the Uhlmann bound optimised over exterior extensions equals the Legendre value \(g_Q\) of the first-order defect for every isometric quasi-free channel, and depends on the channel only through that defect; with a one-parameter-group remainder estimate this gives \(-\log F(\omega,\omega^{\beta_s})\le g_Q(\delta(s))(1+o(1))\) and hence \(E_{\mathrm{rec}}\le(1-\theta)f_2z^2(1+o(1))\), i.e.\ \(\kappa_{\mathrm{opt}}\le0.71\): the zero-collar compression is \emph{not} optimal in exact fidelity. The only remaining numerical input is \(\theta>0\) (equivalently \(E_DuE_D\ne(E_DuE_D)^{*}\)); E2 reconfirms \(\mathrm{UB}_{\mathrm{rot}}/\mathrm{UB}_c=0.708\)--\(0.711\) at \(L=128\)--\(384\) and \(\mathrm{UB}_c/(\tfrac12s^2\lVert u\rVert^2)=1.0000000\), so the earlier factor \(4.3\) was an artefact of a bad exterior extension.
\item \textbf{[OK]} \emph{rigor/exact\_fidelity\_upper\_half.tex Thm.~5.3 (with Lemma~5.1 and Lemma~3.2), Thm.~4.3, Cor.~4.4, Prop.~1.3, Lemma~6.1, Thm.~6.2, Thm.~7.2} \hfill\textsc{RE1}\\ Refereed adversarially (rigor/referee\_exact\_fidelity\_upper\_half.tex). The main theorem \(\overline{\mathrm{Ext}}=\mathcal N^{\perp}\) is \emph{proved}. The pairing identity was re-derived independently from \(\Gamma\Xi+\Xi\Gamma=0\), \(E_I\Xi E_I=0\), \(E_{I^c}\Gamma\Xi E_{I^c}=0\) by inserting \(1=E_I+E_{I^c}\): \(E_{I^c}\Gamma\Xi\Gamma E_{I^c}=-\pi\) and \(=-\Gamma'\pi\Gamma'\), hence \(\Gamma'\pi\Gamma'=\pi\), with all four terms and signs as in the document. Generic position of \((P,E_I)\) is proved there directly from Hardy uniqueness for any \(I\) with \(I,I^c\) of positive measure (so the bounded-interval case is covered, BJL Prop.~3.4 being only a cross-check), and it yields both consequences used: \(\Gamma'=E_{I^c}(2P-1)E_{I^c}\) has no eigenvalue \(\pm1\) (hence \(\Gamma'^n\to0\) strongly) and \(E_{I^c}PE_I\) is injective on \(L^2(I)\) (via \(\ker(A^{1/2}(1-A)^{1/2})=\ker A\oplus\ker(1-A)\), \(A=E_IPE_I\)). The two orthogonality conditions \(E_I(B+B^{*})E_I=0\) and \(E_{I^c}(B-B^{*})E_{I^c}=0\) were recomputed and are correct, with finite-rank exterior generators already separating; the class used for the density statement coincides with the class \(\{Z_e \text{ bounded}, Z_e^{*}=-Z_e=E_{I^c}Z_eE_{I^c}, P^{\perp}Z_eP\in\mathfrak S_2\}\) that the channels can realise as \(e^{sZ_e}\). Moreover the proof is stronger than claimed: the fixed-point step needs no compactness, since \(\langle u,\pi v\rangle=\langle\Gamma'^nu,\pi\Gamma'^nv\rangle\to0\) for every bounded \(\pi\). All numbers quoted from track E2 in \S8 were checked against rigor/exterior\_extension\_numerics.tex and are accurate.
\end{enumerate}
\section{Status of the main results after the check}
\begin{small}
\begin{tabular}{p{0.30\textwidth}p{0.13\textwidth}p{0.50\textwidth}}
\hline
Result & Verdict & Comment \\ \hline
N1 Thm~4.1 (LX adjacent-interval CMI, $c=r$ complex fermions) & OK & coefficient and net identification confirmed from LX Def.~3.17/Thm~3.18 and the LX text; V2's factor-2 claim withdrawn; errata: strong additivity, $\varepsilon\to0$ via Thm~4.1(5) \\
N2 Prop.~2.x (Markov defect, factorization, Petz equivalence) & OK (errata) & hypotheses $D<\infty$, graded product, Klein transformation, FHSW integrated vs.\ FH22 averaged made explicit \\
N4 Thm~7.1 (zero-collar Petz map, trace norm $\tfrac12\log(1+\zeta)$, quasi-equivalence, fidelity bound) & OK (errata) & sign/convention of $s_t=\lambda(1+e^{-2\pi t})$ derived explicitly; Fourier convention fixed; $\delta$-term, Reeh--Schlieder and ultraweak density spelled out \\
N4 Lemma~10.1 (martingale fidelity) & OK (proof fixed) & variational inequality now applied inside each $M_{n_k}$ with Alberti--Uhlmann cited \\
N5 Thm~3.1 (fidelity/entropy as limits over compressions) & OK & Araki77 Thm~3.9(2)+3.8(2)($\gamma$) located; AF hypothesis stated \\
N5 Thm~3.2 (finite-dimensional quasi-free fidelity) & OK & determinant formula re-derived and tested against brute force \\
N5 Prop.~4.1 (UV tail $\tfrac2{15}\zeta^2\kappa^{-3}$) & OK (strengthened) & exact first-order kernel and $c_2=-\zeta^2/30$ (V3) \\
N5 Cor.~4.2 (structure of the small-$\zeta$ expansion) & limit proved & $\lim\Phi/\zeta^2=g_B/8$ (\texttt{quadratic\_limit.tex} + \texttt{uhlmann\_upper\_bound.tex}); the stated $\log^{-2}$ rate remains open \\
N5 Cor.~5.1 (collapse of the rotated family) & OK / corrected & single-chirality statement correct; 2D law is the two-branch formula $\Phi(\zeta_-)+\Phi(\zeta_+)$, even in $\lambda$, optimum at $\lambda=0$ (matches VWZ) \\
N5 numerics ($f_2$, $s_2$) & OK (corrected) & basis-phase bug fixed ($5\times10^{-5}$ effect); tail correction has 12--56\% error, realistic uncertainty 0.3--1\%; $f_2=0.008445(5)$, $s_2=0.0838(10)$ vs.\ $c/(12\pi^2)=0.0084434$, $c/12=0.08333$ \\
N7 Thm~A (metrics $g_B=2r/3\pi^2$, $g_{\rm KM}=r/6$; $\lim\Phi/\zeta^2$) & OK (theorem) & one-particle metrics proved; $\liminf\Phi/\zeta^2\ge r/(12\pi^2)$ (P2) and $\limsup\le r/(12\pi^2)$ by the Uhlmann--Bogoliubov route (\texttt{uhlmann\_upper\_bound.tex}, PA; refereed by PR, three repairs applied): $\lim\Phi/\zeta^2=c/(12\pi^2)$ is a theorem; only the rate is open \\
N7 Thm~B (universality for diffeomorphism covariant nets) & Bures half OK / KM half GAP & the fidelity half of Lemma~3.2 is now proved for every von Neumann algebra with cyclic separating $\Omega$ (\texttt{lemma\_second\_variation.tex} Thm~3.3, three-way identity $\tfrac14g_B=\operatorname{dist}(a,\overline{\cM'_{\rm sa}\Omega})^2=\sup_K[\varphi(K)-\operatorname{Var}(K)]$); Kubo--Mori half open; the $\Delta^{1/2}$ formula was false (it is $Q_{1/2}$); zero-collar representative: upper bound survives, lower bound needs the cut-off tangents (Prop.~10.2) \\
N7 Thm~5.2 (integral $=2$, $=\pi$, coefficients) & OK & factor-4 slip in the proof text fixed; integrals verified analytically and numerically; V5 re-derived the whole chain in Lamport form (\texttt{check\_notes\_6\_7.tex}, 25 pp.): $\langle TT\rangle=(c/8\pi^2)(x-y-i0)^{-4}$ proved independently from $[T(1),\psi]=-i\partial\psi$ and positivity, the base integral $\int w^2/(\cosh w+\cosh a)$ proved by a $2\pi i$-rectangle, both routes agree integrand by integrand: $f_2=c/(12\pi^2)$, $s_2=c/12$ stand \\
N7 Lemma~5.1(ii) (kernel identity) & OK (proved, refereed) & \texttt{kernel\_identity.tex} (P3): closed form via digamma; referee V8: correct, no principal value or regularization is involved, independent numerics to $10^{-22}$ \\
N7 Cor.~6.1 (VWZ comparison) & corrected & two chiralities with opposite rotations; $\ell_C$ vs.\ central charge; $\eta_{\rm V}$ \\
N6 strategy (citations) & errata & $\langle TT\rangle$ normalization, CDIT authors, energy bound via Carpi--Weiner/CDIT Prop.~2.3, WP chain gap; V5: \S2 step 3 used the BPZ normalization $c/32$ inconsistent with step 4 (fixed), Route~B energy set unbounded as written, eq.~(energy) is $E(\varphi^{-1})$ (fixed) \\
\hline
\end{tabular}
\end{small}

\section{Status of the remaining items (a)--(d) of Note~7}
\begin{enumerate}[leftmargin=1.6em]
\item[(a)] \emph{Lemma~3.2 (second variation of the fidelity in type III).} Document \texttt{lemma\_second\_variation.tex} (P1, 33 pp., 157 steps). Fidelity half proved unconditionally (any von Neumann algebra, cyclic separating $\Omega$, $\Omega\in D(A)$, $A\Omega\in D(\Delta^{-1/2})$): lower bound from the Alberti--Uhlmann supremum formula with $v'=e^{isB'}$, upper bound from Alberti's inequality with $x=1+sK$; six independent expressions agree numerically to 10 digits. Corrections to Note~7: Lemma~3.1 needs $A\Omega\in D(\Delta^{-1/2})$; the $\Delta^{1/2}$ route deleted; $c$ redundant over $\mathbb R$. Open: Kubo--Mori half beyond finite dimensions; zero-collar identification (Prop.~10.2).
\item[(b)] \emph{Quadratic limit $\lim\Phi/\zeta^2=g_B/8$ for the free fermion.} Document \texttt{quadratic\_limit.tex} (P2, 25 pp.). Done in part: $\liminf\ge g_B/8$ unconditional (Legendre form of the SLD metric, single finite-rank compression, density argument); $g_B=2r/(3\pi^2)$ fully analytic; unconditional $\Phi\le\tau/(1-\tau)$; Hellinger second-order coefficient $h_2=2\log2\,f_2$ in closed form; cutoff error $0\le g_B-g_B^{(\Lambda)}\le\frac{2r}{3\Lambda^2}$. Open: the interchange in Lemma~6.7, the sharp $\limsup$ (Problem~6.8), the cubic remainder (Problem~7.2). Refereed in \texttt{referee\_quadratic\_limit.tex} (V7): Thm~4.2 verified unconditional, all closed forms re-derived to 25--31 digits; one false step in the proof of Prop.~7.1 ($\langle1\rangle5$, an unrestricted $u$-integration) found and repaired (the stated inequality survives; exact asymptotics $0.513\,r\Lambda^{-2}$); the explanation of the $\zeta=1/120$ value in \S7.3 is circular and corrected in the appended note.
\item[(c)] \emph{Analytic proof of the kernel identity.} Done: \texttt{kernel\_identity.tex} (P3), verified numerically by \texttt{p3\_check.py} to $10^{-12}$; refereed in \texttt{referee\_kernel\_identity.tex} (V8): theorem and proof correct, no gap; independent evaluation of the defining double integral agrees with the closed form to $10^{-22}$ relative at random points (\texttt{v8\_kernel\_check.py}); the vanishing of the $D\times D$ block, asserted in Note~7 without computation, is proved there. Minor: P3's Remarks~6.4--6.5 refer to the pre-errata text of Note~7 (an orchestrator note is appended).
\item[(d)] \emph{Collar limit and the invariant formula $g_B=2c/(3\pi^2)$ in $x$-space.} Document \texttt{collar\_and\_invariant\_formula.tex} (8 pp.) + \texttt{p4\_invariant.py} (P4b; the original P4 died four times without output). Done for the invariant formula: in $u=L-x$ the Mellin transform diagonalizes the modular dilation, $\widehat W(k)=\frac{c}{24\pi}\frac{k(k^2+1)}{e^{2\pi k}-1}$ follows \emph{exactly} from $\langle TT\rangle$ and the geometric dilation (closed form via a Beta function; numerics to $10^{-23}$), KMS holds on the $x$-space transform, and the reassembled $g_B=0.0675474557603\,c$, $g_{\rm KM}=0.166666666665\,c$ agree with $2c/(3\pi^2)$ and $c/6$ to $10^{-11}$. Diagnosis of the formal $\xi$-frame step: the Mellin--Barnes strips are disjoint precisely because $g(L)\neq0$. Not proved there: identification of $g_B(\varphi)$ with the spectral integral at zero collar, exchange of the collar limit with the $s\to0$ expansion, domain bookkeeping for $\Delta^{it}T(u)\Delta^{-it}$.
\end{enumerate}
All agents have reported; see also \texttt{INDEX.md} and \texttt{FINAL\_REPORT.md}.

\section{Errata applied to the notes}
All confirmed errata have been applied to the \LaTeX\ sources of Notes 1--7 and the PDFs recompiled (second pass 2026-09-07: every one of the 93 findings re-checked against the notes); the log \texttt{errata\_log.md} lists every edit with its anchor, the agents that found it, and the timestamp (120 logged edits over two passes, 2 anchors that failed on the first attempt and were re-applied with corrected anchors). The largest corrections are: the two-chirality law for 2D CFTs (Note~5 Cor.~5.1, Note~7 Cor.~6.1, Note~6); the withdrawal of the $\Delta^{1/2}$ fidelity formula (Note~7 Lemma~3.2 route, Note~6); the explicit derivation of the rotation convention $s_t=\lambda(1+e^{-2\pi t})$ (Note~4); the corrected description of the numerical method (Note~5 \S3.1, \S6) and the basis-phase fix in the code; the relabelling of Theorem~A (Note~7).

\end{document}